%%%%%%%%%%%%%%%%%%%%%%%%%%%%%%%%%%%%%%%%%%%%%%%%%%%%%%%%%%%%%%%%%%%%%%
% Lớp 11 | Đường thẳng, mặt phẳng. Quan hệ song song trong | Hai mặt phẳng song song
% Dạng 1H4?4-6: Bài toán tổng hợp
% Số câu: 69  (sinh tự động bởi tach_cau.py; từ nay thầy sửa trực tiếp file này)
%%%%%%%%%%%%%%%%%%%%%%%%%%%%%%%%%%%%%%%%%%%%%%%%%%%%%%%%%%%%%%%%%%%%%%

%%%%%%%%%%%%%%%% TRẮC NGHIỆM 4 PHƯƠNG ÁN %%%%%%%%%%%%%%%%

% nguồn: Ok_Toan11_TracNghiem_Lan2
\begin{ex}%[1H4H4-6]%[ma:60de1995]
	Cho hình hộp $ABCD.A'B'C'D'$. Gọi $I$ là trung điểm $AB$. Mặt phẳng $(IB'D')$ cắt hình hộp theo  thiết diện là hình gì?
	\choice 
		 {Tam giác}
		{\True Hình  thang}
		{Hình bình hành} 
		 {Hình chữ nhật}
	\loigiai{\begin{center}
		\begin{tikzpicture}
				 \def\a{5}
				\def\b{2} 
				\def\h{3} 
				\path (0:0) coordinate  (A)
					++(0:\a)  coordinate (D)
					++(-130:\b)  coordinate (C)
					($(A)+(C)-(D)$) coordinate  (B)
					($(A)+(90:\h)$)  coordinate (A')
					($(B)+(90:\h)$)  coordinate (B')
					($(C)+(90:\h)$)  coordinate (C')
					($(D)+(90:\h)$)  coordinate (D')
					($(A)!1/2!(B)$) coordinate (I) 
					($(A)!1/2!(D)$) coordinate (J); 
				\draw[dashed]  (B)--(A)--(D)
					 (A)--(A')
					(I)--(J) 
					 (B)--(D)
					 (I)--(B')
					 (J)--(D');
				 \draw (C)--(C')
					(D)--(D') 
					(B)--(B') 
					 (C)--(C')
					 (B)--(C)--(D)
					(A')--(B')--(C')--(D')--cycle  (B')--(D');
				\foreach \x/\g in  {A/180,B/180,C/0,D/0,A'/180,B'/180,C'/0,D'/0,I/-60,J/-60}
				 \fill (\x) circle (1pt)
					 ($(\g:4mm)+(\x)$) node {$\x$};
			\end{tikzpicture}
	\end{center}
	Gọi $J$ là trung điểm của $AD$.  Do đó $IJ \parallel BD$ nên $IJ \parallel B'D' \Rightarrow IJ$ thuộc mặt phẳng $(IB'D')$ và $IJ = (ABCD) \cap (IB'D')$.\\
	Lại có $JD' = (IB'D') \cap (ADD'A')$, $IB' = (IB'D') \cap  (ABB'A')$ và $B'D'=(IB'D')\cap (A'B'C'D')$.\\
Vậy thiết diện cần tìm là hình thang $IJD'B'$.} 
\end{ex}

% nguồn: Ok_Toan11_TracNghiem_Lan2
\begin{ex}%[1H4H4-6]%[ma:a2de941d]
	Cho hình lập phương $ABCD.A'B'C'D'$. Mặt phẳng $(P)$ chứa $BD$ và song song với  mặt phẳng $(AB'D')$ cắt hình lập phương theo thiết diện là.
	 \choice
		{\True Một tam  giác đều}
		{Một tam  giác thường}
		{Một hình  chữ nhật}
		{Một hình  bình hành}
	\loigiai{ 
		\begin{center}
			\begin{tikzpicture}
					\def\a{3.5} 
					\path (0:0) coordinate (A) 
						++(0:\a)  coordinate (D)
						++(-130:\a/2)  coordinate (C)
						($(A)+(C)-(D)$)  coordinate (B)
						($(A)+(90:\a)$)  coordinate (A')
						 ($(B)+(90:\a)$) coordinate (B')
						($(C)+(90:\a)$)  coordinate (C')
						 ($(D)+(90:\a)$) coordinate (D');
					 \draw[dashed] (B)--(A)--(D)
						(A)--(A') 
						(A)--(B') 
						 (A)--(D')
						 (B)--(D);
					\draw (C)--(C') 
						 (D)--(D')
						(B)--(B') 
						(B)--(C)--(D) 
						(A')--(B')--(C')--(D')--cycle (B')--(D') 
						 (B)--(C')--(D);
					\foreach \x/\g in  {A/180,B/180,C/0,D/0,A'/180,B'/180,C'/0,D'/0}
					\fill  (\x) circle (1pt)
						($(\g:4mm)+(\x)$)  node {$\x$};
				\end{tikzpicture}
		\end{center}
	Do $BC'$ song  song với $AD'$, $DC'$ song song với $AB'$ nên thiết diện cần tìm là tam giác đều $BDC'$.}
\end{ex}

% nguồn: Ok_Toan11_TracNghiem_Lan2
\begin{ex}%[1H4H4-6]%[ma:cb5db451]
	Cho hình lập  phương $ABCD.A'B'C'D'$ cạnh $a$. Mặt phẳng $(\alpha)$ qua $AC$ và song song với $BB'$. Tính chu vi thiết diện của hình lập phương $ABCD.A'B'C'D'$ khi cắt bởi mặt phẳng $(\alpha)$.
	 \choice
		{\True $2(1+\sqrt{2})a$} 
		{$a^3$} 
		 {$a^2\sqrt{2}$}
		 {$(1+\sqrt{2})a$}
	\loigiai{\begin{center}
		\begin{tikzpicture}
				 \def\a{3.5}
				 \path (0:0) coordinate (A)
					++(0:\a) coordinate  (D)
					++(-130:\a/2) coordinate  (C)
					($(A)+(C)-(D)$)  coordinate (B)
					($(A)+(90:\a)$)  coordinate (A')
					 ($(B)+(90:\a)$) coordinate (B')
					($(C)+(90:\a)$)  coordinate (C')
					 ($(D)+(90:\a)$) coordinate (D');
				\draw[dashed] (B)--(A)--(D) 
					(A)--(A') 
					(A)--(C); 
				 \draw (C)--(C')
					(D)--(D') 
					(B)--(B') 
					 (B)--(C)--(D)
					 (A')--(B')--(C')--(D')--cycle (A')--(C');
				\foreach \x/\g in  {A/180,B/180,C/0,D/0,A'/180,B'/180,C'/0,D'/0}
				\fill  (\x) circle (1pt)
					($(\g:4mm)+(\x)$)  node {$\x$};
			\end{tikzpicture}
	\end{center}
	Ta dễ dàng dựng được thiết diện là tứ giác $ACC'A'$. Tứ giác $ACC'A'$ là hình chữ nhật có chiều dài là $AC=a\sqrt{2}$ và  chiều rộng $AA'=a$.\\
	Khi đó  chu vi thiết diện của hình lập phương $ABCD.A'B'C'D'$ khi cắt bởi mặt phẳng $(\alpha)$ là
	$P=2(AC+AA')=2(1+\sqrt{2})a$. 
 }
\end{ex}

% nguồn: Ok_Toan11_TracNghiem_Lan2
\begin{ex}%[1H4H4-6]%[ma:3db8d352]
	Cho tứ diện $ABCD$ có tất cả các cạnh đều bằng nhau. Gọi $I$ là trung điểm đoạn $CD$, $M$ là điểm nằm trên đoạn $BC$ ($M$ khác $B$ và $C$). $(\alpha)$ là mặt phẳng qua $M$ và song song với mặt phẳng $(ABI)$,  khi đó thiết diện của tứ diện $ABCD$ khi cắt bởi $(\alpha)$ là
	\choice 
		 {Một tam giác vuông cân}
		 {Một tam giác đều}
		{Một hình bình  hành}
		{\True  Một tam giác cân}
	\loigiai{\begin{center}
		\begin{tikzpicture}
				 \def\a{5}
				\path (0:0) coordinate  (B)
					++(0:\a) coordinate (D) 
					++(-150:\a/1.5) coordinate  (C)
					($(B)+(70:\a)$) coordinate  (A)
					($(A)!2/5!(C)$)  coordinate (Q)
					($(D)!1/2!(C)$)  coordinate (I)
					($(B)!2/5!(C)$) coordinate (M) 
					($(I)!2/5!(C)$) coordinate (P); 
				\draw[dashed] (B)--(D) 
					 (M)--(P)
					 (B)--(I);
				 \draw (B)--(A)--(D)
					 (A)--(C)
					 (B)--(C)--(D)
					 (M)--(Q)--(P)
					 (A)--(I);
				\foreach \x/\g in  {A/90,B/180,C/-45,D/0,M/-90,P/-90,I/-90,Q/30}
				\fill  (\x) circle (1pt)
					 ($(\g:3mm)+(\x)$) node {$\x$};
			\end{tikzpicture}
	\end{center}
	Theo giả thiết ta có $IA = IB$ suy ra $\triangle AIB$ cân tại $I$. Do $(\alpha)$ là mặt phẳng qua $M$ và song song với mặt phẳng $(ABI)$ nên:\begin{itemize}
	\item Trong $(BCD)$, kẻ $MP \parallel BI$, $P \in  CD$ suy ra $MP = (\alpha) \cap (BCD)$.
	\item Trong $(ACD)$, kẻ $PQ \parallel AI$, $Q \in AC$ suy ra $PQ = (\alpha) \cap  (ACD)$.
	\item $MQ = (\alpha)  \cap (ABC)$.
	\end{itemize}
	Thiết diện diện của tứ diện $ABCD$ khi cắt bởi $(\alpha)$ là $\triangle  MPQ$.\\
	Theo cách dựng ta suy ra $\triangle  MPQ$ đồng dạng với $\triangle BIA$ suy ra $\triangle MPQ$ cân tại $P$.
 }
\end{ex}

% nguồn: Ok_Toan11_TracNghiem_Lan2
\begin{ex}%[1H4H4-6]%[ma:25139dab]
	Cho hình chóp $S.ABCD$, đáy $ABCD$ là tứ giác có các cạnh đối không song song. Gọi $O=AC \cap BD$, $F=BC \cap AD$. Điểm $M$ thuộc cạnh  $SA$. Tìm giao tuyến $(d)$ của cặp mặt phẳng $(MBD)$ và $(SAC)$.
	\begin{center}
		\begin{tikzpicture}
				 \def\a{4}
				\path  (0:0) coordinate (A)
					++(0:\a) coordinate (B) 
					 ++(-130:\a/1.5) coordinate (C)
					++(165:\a/2.5) coordinate  (D)
					($(A)+(70:\a)$)  coordinate (S)
					(intersection of  A--C and B--D) coordinate (O)
					(intersection of  A--D and B--C) coordinate (F)
					($(S)!1/3!(A)$)  coordinate (M);
				\draw[dashed] (A)--(B) 
					(M)--(B)--(D) 
					 (A)--(C);
				\draw (A)--(D)--(C)--(B) 
					(A)--(S) 
					(B)--(S) 
					(C)--(S) 
					 (D)--(S)
					(D)--(M) 
					(D)--(F)--(C); 
				\foreach \x/\g in {A/180,B/0,C/-45,D/-135,S/90,M/150,O/90,F/-90} 
				\fill (\x)  circle (1pt)
					($(\g:3mm)+(\x)$) node {$\x$}; 
			\end{tikzpicture}
	\end{center}
	\choice 
		 {$d=SO$}
		{$d=SF$} 
		 {\True $d=MO$}
		{$d=MF$} 
	 \loigiai{
		\begin{center}
			\begin{tikzpicture}
					 \def\a{4}
					\path (0:0) coordinate (A) 
						 ++(0:\a) coordinate (B)
						 ++(-130:\a/1.5) coordinate (C)
						++(165:\a/2.5) coordinate  (D)
						($(A)+(70:\a)$) coordinate (S) 
						(intersection of A--C  and B--D) coordinate (O)
						(intersection of A--D and B--C)  coordinate (F)
						($(S)!1/3!(A)$) coordinate  (M);
					\draw[dashed]  (A)--(B)
						 (M)--(B)--(D)
						(A)--(C) 
						 (M)--(O);
					\draw (A)--(D)--(C)--(B) 
						 (A)--(S)
						(B)--(S) 
						(C)--(S) 
						 (D)--(S)
						 (D)--(M)
						(D)--(F)--(C); 
					\foreach  \x/\g in {A/180,B/0,C/-45,D/-135,S/90,M/150,O/90,F/-90}
					\fill  (\x) circle (1pt)
						 ($(\g:3mm)+(\x)$) node {$\x$};
				\end{tikzpicture}
		\end{center}
		Ta có $\heva{&M \in (SAC) \cap  (MBD) \\
		&O \in (SAC) \cap (MBD)\, (O=AC \cap BD)} \Rightarrow MO=(SAC) \cap  (MBD).$
	} 
\end{ex}

% nguồn: Ok_Toan11_TracNghiem_Lan2
\begin{ex}%[1H4H4-6]%[ma:f3a877a3]
	Cho hình chóp $S.ABCD$ có đáy là hình thang, $(AB \parallel CD)$. Gọi $O$ là giao điểm của $AC$ và $BD$, $I$ là giao điểm của $AD$  và $BC$. Khẳng định nào sau đây sai?
	\begin{center}
		\begin{tikzpicture}
				\def\a{4} 
				\path (0:0)  coordinate (A)
					++(0:\a) coordinate  (B)
					++(-130:\a/2)  coordinate (C)
					++(180:\a/2) coordinate  (D)
					 ($(A)+(70:\a)$) coordinate (S)
					(intersection of A--C and B--D)  coordinate (O)
					(intersection of  A--D and B--C) coordinate (I);
				 \draw[dashed] (A)--(B)
					 (A)--(C)
					(B)--(D); 
				 \draw (A)--(D)--(C)--(B)
					(A)--(S) 
					 (D)--(S)
					(C)--(S) 
					 (B)--(S)
					 (D)--(I)--(C);
				\foreach \x/\g in {A/180,D/-135,C/-45,B/0,S/90,O/90,I/-90} 
				\fill (\x) circle  (1pt)
					($(\g:3mm)+(\x)$) node {$\x$}; 
			\end{tikzpicture}
	\end{center}
	\choice 
		{Giao tuyến của hai mặt phẳng $(SAC)$ và $(SBC)$ là $SC$} 
		{Giao tuyến của  hai mặt phẳng $(SAC)$ và $(SBD)$ là $SO$}
		 {Giao tuyến của hai mặt phẳng $(SAD)$ và $(SBC)$ là $SI$}
		{\True  Giao tuyến của hai mặt phẳng $(SID)$ và $(SCO)$ là $SI$}
	\loigiai{\begin{center}
		\begin{tikzpicture}
				\def\a{4} 
				\path  (0:0) coordinate (A)
					++(0:\a) coordinate  (B)
					++(-130:\a/2)  coordinate (C)
					++(180:\a/2) coordinate (D) 
					($(A)+(70:\a)$) coordinate (S) 
					(intersection of A--C and B--D)  coordinate (O)
					(intersection of  A--D and B--C) coordinate (I);
				\draw[dashed]  (A)--(B)
					 (A)--(C)
					(B)--(D) 
					 (D)--(C);
				\draw (A)--(D) 
					(C)--(B) 
					(A)--(S) 
					(D)--(S) 
					 (C)--(S)
					 (B)--(S)
					(S)--(I) 
					(D)--(I)--(C); 
				 \foreach \x/\g in {A/180,D/-135,C/-45,B/0,S/90,O/90,I/-90}
				 \fill (\x) circle (1pt)
					 ($(\g:3mm)+(\x)$) node {$\x$};
			\end{tikzpicture}
	\end{center}
	\begin{itemize}
		\item $S$, $C$ là hai điểm chung của $(SAC)$ và $(SBC)$  nên A đúng.
		\item $S$,  $O$ là hai điểm chung của $(SAC)$ và $(SBD)$ nên B đúng.
		\item $S$,  $I$ là hai điểm chung của $(SAD)$ và $(SBC)$ nên C đúng.
		\item $S$, $A$ là hai điểm chung của $(SID)$ và $(SCO)$  nên D sai.
	\end{itemize}
} 
\end{ex}

% nguồn: Ok_Toan11_TracNghiem_Lan2
\begin{ex}%[1H4H4-6]%[ma:f78c75d6]
	Cho hình chóp $S.ABCD$ có đáy $ABCD$ là hình bình hành tâm $O$. Gọi $M$, $N$, $P$, $Q$ lần lượt là trung điểm của các cạnh $AD$, $BC$, $SC$, $SD$. Gọi $(\alpha)$ là mặt phẳng đi qua $O$ và song song với mặt phẳng $(SAB)$. Giao tuyến của $(\alpha)$ với các mặt phẳng  $(SBC)$ và $(SAD)$ lần lượt là
	\begin{center}
		\begin{tikzpicture}
				 \def\a{5}
				\path (0:0) coordinate  (A)
					++(0:\a) coordinate  (D)
					++(-130:\a/2)  coordinate (C)
					($(A)+(C)-(D)$) coordinate (B) 
					 ($(A)+(80:\a)$) coordinate (S)
					(intersection of A--C  and B--D) coordinate (O)
					($(A)!0.5!(D)$)  coordinate (M)
					($(B)!0.5!(C)$) coordinate (N) 
					($(S)!0.5!(D)$)  coordinate (Q)
					($(S)!0.5!(C)$)  coordinate (P);
				 \draw[dashed] (B)--(A)--(D)
					(A)--(S) 
					(A)--(C) 
					 (B)--(D);
				 \draw (B)-- (C)--(D)
					(B)--(S) 
					(C)--(S) 
					 (D)--(S);
				\foreach \x/\g in  {A/135,B/-135,C/-45,D/45,S/90,O/-90,M/135,N/-60,P/-120,Q/60}
				\fill (\x) circle  (1pt)
					 ($(\g:3mm)+(\x)$) node {$\x$};
			\end{tikzpicture}
	\end{center}
	 \choice
		 {$MN$ và $PN$}
		{$MN$ và $PQ$} 
		{$QP$ và $QM$} 
		{\True  $NP$ và $MQ$}
	\loigiai{\begin{center}
		\begin{tikzpicture}
				 \def\a{5}
				\path (0:0)  coordinate (A)
					++(0:\a) coordinate (D) 
					++(-130:\a/2)  coordinate (C)
					($(A)+(C)-(D)$) coordinate  (B)
					 ($(A)+(80:\a)$) coordinate (S)
					(intersection  of A--C and B--D) coordinate (O)
					 ($(A)!0.5!(D)$) coordinate (M)
					($(B)!0.5!(C)$)  coordinate (N)
					($(S)!0.5!(D)$) coordinate  (Q)
					 ($(S)!0.5!(C)$) coordinate (P);
				\draw[dashed] (B)--(A)--(D) 
					 (A)--(S)
					 (A)--(C)
					 (B)--(D)
					(N)--(M)--(Q); 
				\draw (B)--  (C)--(D)
					(B)--(S) 
					(C)--(S) 
					 (D)--(S)
					 (N)--(P)--(Q);
				\foreach  \x/\g in {A/135,B/-135,C/-45,D/45,S/90,O/-90,M/135,N/-60,P/-150,Q/60}
				\fill (\x) circle (1pt) 
					($(\g:3mm)+(\x)$)  node {$\x$};
			\end{tikzpicture}
	\end{center}
	Vì $(\alpha) \parallel (SAB)  \Rightarrow (\alpha) \parallel AB$, $(\alpha) \parallel SB$, $(\alpha)\parallel SA$.\\
	Vì $M$, $N$, $P$, $Q$ lần  lượt là trung điểm của các cạnh $AD$, $BC$, $SC$, $SD$ nên
	$MN \parallel AB$, $NP  \parallel SB$, $MQ \parallel SA$.
	Ta có\begin{itemize}
	\item $\heva{&(\alpha)\parallel AB \\& AB \subset (ABCD)  \\
	& O \in (\alpha), O \in (ABCD)} \Rightarrow (\alpha) \cap  (ABCD)=MN$ qua $O$.
	\item $\heva{&(\alpha)  \parallel SB \\
	& SB \subset (SBC)  \\
	& N \in (\alpha), N  \in (SBC)}\Rightarrow (\alpha) \cap (SBC)=NP \parallel SB$.
	\item $\heva{&(\alpha)\parallel SA \\& SA  \subset (SAD) \\
	& M \in (\alpha), M \in (SAD)} \Rightarrow (\alpha) \cap (SAD)=MQ  \parallel SA$.
	\end{itemize}
	 Vậy giao tuyến của $(\alpha)$ với các mặt phẳng $(SBC)$ và $(SAD)$ lần lượt là $NP$ và $MQ$.
 }
\end{ex}

% nguồn: Ok_Toan11_TracNghiem_Lan2
\begin{ex}%[1H4V4-6]%[ma:018cfb8d]
	Cho hình chóp $ S.ABCD$ có mặt phẳng đáy là hình bình hành tâm $ O$ và $ AC=a$, $BD=b$.  Tam giác $ SBD$ đều. Gọi $(P)$ là mặt phăng di động đi qua điểm $ I$ trên đoạn $ OC$, song song với $\left(SBD\right)$. Đặt $ AI=x~\left(\dfrac{a}{2}< x < a\right)$, cắt các cạnh $ BC$, $CD$, $SC$ lần lượt tại $ E$, $F$, $G$. Diện tích tam giác $ EFG$ bằng
	 \choice
		{$\dfrac{b^2\left(a-x\right)^2\sqrt 2}{a^2}$} 
		{$\dfrac{b^2\left(a+x\right)^2\sqrt  3}{a^2}$}
		{$\dfrac{b^2\left(a+x\right)^2}{a^2\sqrt  3}$}
		{\True $\dfrac{b^2\left(a-x\right)^2\sqrt 3}{a^2}$} 
	 \loigiai{
	 \immini{
		 Ta có $\heva{&(P)\parallel\left(SBD\right)\\&(P)\cap\left(ABCD\right)=EF\\
		&\left(SBD\right)\cap\left(ABCD\right)=EF} 
		\Rightarrow EF\parallel BD$.\\ 
		Suy  ra $\dfrac{EF}{BD}=\dfrac{CE}{CB}=\dfrac{CF}{CD}$.\\
		Chứng minh tương tự  ta được
		\begin{itemize}
			\item $EG\parallel  SB\Rightarrow\dfrac{EG}{SB}=\dfrac{CE}{CB}$.
			\item $GF\parallel  SD\Rightarrow\dfrac{GF}{SD}=\dfrac{CF}{CD}$.
		\end{itemize}
		Suy  ra $\dfrac{EF}{BD}=\dfrac{EG}{SB}=\dfrac{GF}{SD}$.\\
		Lại do $ SB=SD=BD$ nên $ EG=GF=EF$  hay tam giác $ GEF$ đều.\\
		Ta có $\dfrac{EF}{BD}=\dfrac{CI}{CO}\Rightarrow EF=\dfrac{CI}{CO}\cdot BD=\dfrac{2\left(a-x\right)}{a}\cdot  b$.\\
		Vậy $S_{GEF}=\dfrac{E{F^2}\sqrt 3}{4}=\dfrac{b^2\cdot\left(a-x\right)^2\sqrt 3}{a^2}$.}{\begin{tikzpicture}[scale=1,font=\footnotesize,line join=round,line cap=round,>=stealth]
		 \path
			(0,0) coordinate  (A)
			(-1.5,-1.5)  coordinate (D)
			(4,0) coordinate (B) 
			($(B)+(D)-(A)$) coordinate  (C)
			(1,3)  coordinate (S)
			($(B)!0.5!(D)$)  coordinate (O)
			($(C)!1/3!(O)$) coordinate (I) 
			($(C)!1/3!(B)$) coordinate (E) 
			($(C)!1/3!(D)$) coordinate  (F)
			($(C)!1/3!(S)$) coordinate  (G);
		\draw 
			 (S)--(D)--(C)--(B)--(S)--(C)
			(F)--(G)--(E); 
		 \draw[dashed]
			(S)--(A)--(C) 
			(B)--(A)--(D)--cycle 
			(F)--(E); 
		\foreach \p/\g  in {A/170,D/-90,C/-90,B/0,S/90,G/10,E/-10/,F/-90,O/90,I/90}
		\draw[fill=black] (\p) circle  (1pt) node[shift=(\g:3mm)] {$\p$};
	\end{tikzpicture}
		} 
	 }
\end{ex}


%%%%%%%%%%%%%%%% ĐÚNG/SAI %%%%%%%%%%%%%%%%

% nguồn: Ok_dungsai_lop11_Lan1
\begin{ex}%[1H4H4-6]%[ma:c65a60f5]
 Cho hình chóp $S.ABCD$ có đáy $ABCD$ là hình bình hành tâm $O$. Gọi $M$ là trung điểm của $CD$, điểm $G$ là trọng tâm tam giác $BCD$. Mặt phẳng $(\alpha)$ qua $G$, song song với $(SAB)$, cắt cạnh bên $SC$ tại $I$ và cắt trung tuyến $SM$ của tam giác $SCD$ tại  $K$.
  \choiceTF
  {\True $CD\parallel (\alpha)$} 
   {\True $\dfrac{CG}{CA}=\dfrac{1}{3}$}
  {$\dfrac{CI}{CS}=\dfrac{1}{2}$} 
  {\True  $K$ là trọng tâm tam giác $SCD$}
  \loigiai{
  \begin{center}
   \begin{tikzpicture}[scale=1,font=\footnotesize,line join=round,line cap=round,>=stealth]
     \coordinate (A)  at (0,0);
     \coordinate (B) at (-1.3,-1.7); 
     \coordinate (D) at  (4,0);
     \coordinate (C) at ($(B)+(D)-(A)$); 
     \coordinate  (S) at (1.3,4);
     \coordinate (M) at  ($(C)!1/2!(D)$);
     \coordinate (O) at  ($(B)!1/2!(D)$);
     \coordinate  (K) at ($(M)!1/3!(S)$);
     \coordinate (I) at ($(C)!1/3!(S)$); 
      \coordinate (J) at ($(D)!1/3!(S)$);
     \coordinate (G)  at ($(C)!1/3!(A)$);
      \draw (S)--(B)--(C)--(D)--(S)--(C)
      (S)--(M) 
       (I)--(J);
     \draw[dashed]  (B)--(A)--(S)
      (M)--(B)--(D)--(A)--(C) 
       (G)--(I);
     \foreach \x/\g in{A/135,B/180,C/-45,D/0,M/0,G/-135,S/90,I/180,K/-25,O/90,J/25} \fill (\x) circle (1pt)+(\g:.35)  node{$\x$};
    \end{tikzpicture}
  \end{center}
  \begin{itemchoice}
    \itemch \textbf{Đúng.}\\
   Vì $\heva{&(\alpha)\parallel (SAB)\\&CD\parallel (SAB)\\&CD\not \subset  (\alpha) \text{ (vì ngược lại ta có } G\in CD \text{)}}\Rightarrow CD\parallel (\alpha)$.
   \itemch \textbf{Đúng.}\\ 
   Vì $G$ là trọng tâm của tam giác $BCD$  nên $CG=\dfrac{2}{3}CO=\dfrac{1}{3}CA\Rightarrow \dfrac{CG}{CA}=\dfrac{1}{3}$.
   \itemch \textbf{Sai.}\\ 
   Vì ta có $\heva{&(SAB)\parallel (\alpha)\\&(SAC)\cap (SAB)=SA\\&(SAC)\cap (\alpha)=GI}  \Rightarrow GI\parallel SA$ nên $\dfrac{CI}{CS}=\dfrac{GI}{SA}=\dfrac{CG}{CA}=\dfrac{1}{3}$.
   \itemch  \textbf{Đúng.}\\
   Ta có $\heva{&CD\parallel (\alpha)\\&CD\subset (SCD)\\&(\alpha) \cap (SCD)=IJ}\Rightarrow IJ\parallel  CD \Rightarrow IK\parallel CD \ (\text{vì } IK\equiv IJ) \Rightarrow \dfrac{SK}{SM}=\dfrac{SI}{SC}=\dfrac{2}{3}$.\\
   Suy ra $K$ là trọng tâm tam giác  $SCD$.
  \end{itemchoice}
  }
\end{ex}

% nguồn: Ok_dungsai_lop11_Lan1
\begin{ex}%[1H4H4-6]%[ma:5499f16e]
 Cho hình chóp  $S.ABCD$ có đáy là hình bình hành tâm $O$. Gọi $M,\,N$ lần lượt là trung điểm của $SA$ và $SD$. Khi đó:
  \choiceTF
   {$MN \parallel (SBC)$}
   {$(OMN) \parallel (SBC)$}
  {Gọi $E$ là trung điểm đoạn $AB$ và $F$  là một điểm thuộc đoạn $ON$. Khi đó $EF$ cắt với mặt phẳng $(SBC)$}
  {Gọi $G$ là một điểm nằm trên mặt phẳng $(ABCD)$ cách đều $AB$ và $CD$. Khi  đó $GN$ cắt $(SAB)$}
  \loigiai{
   Ta xét hình vẽ:
  \begin{center}
   \begin{tikzpicture}
     \coordinate (A) at (0,0); 
     \coordinate (S) at  (1,4);
     \coordinate (D) at (4,0); 
      \coordinate (B) at (-2,-2);
     \coordinate (C) at (2,-2); 
     \coordinate (M) at  ($(S)!0.5!(A)$);
     \coordinate  (N) at ($(S)!0.5!(D)$);
     \coordinate (E) at  ($(A)!0.5!(B)$);
     \coordinate  (O) at ($(A)!0.5!(C)$);
     \coordinate (F) at  ($(N)!0.25!(O)$);
      \coordinate (P1) at ($(A)!0.5!(D)$);
      \coordinate (I) at ($(B)!0.5!(C)$);
      \coordinate (G) at ($(P1)!0.5!(O)$);
      \draw [thick] (S) -- (B);
      \draw [thick] (S) -- (C);
     \draw [thick] (S)  -- (D);
     \draw [thick] (B) -- (C) --  (D);
      \draw [dashed] (A) -- (B);
      \draw [dashed] (A) -- (C);
     \draw [dashed] (A) --  (D);
     \draw [dashed] (A) --  (S);
     \draw [dashed]  (O) -- (E) -- (M) -- (N) -- cycle;
     \draw [dashed] (E)  -- (F);
     \draw [dashed] (M) -- (O); 
     \draw [dashed] (I) -- (P1); 
     \draw [dashed] (M) -- (G); 
     \filldraw (S) circle (1pt) node[above] {$S$}; 
     \filldraw (M) circle (1pt) node[above right]  {$M$};
     \filldraw (N) circle (1pt)  node[above right] {$N$};
     \filldraw  (F) circle (1pt) node[right] {$F$};
     \filldraw (A) circle (1pt) node[above right] {$A$}; 
     \filldraw (D) circle (1pt)  node[right] {$D$};
     \filldraw (E) circle  (1pt) node[below right] {$E$};
     \filldraw (O) circle (1pt) node[below] {$O$}; 
      \filldraw (B) circle (1pt) node[below left] {$B$};
      \filldraw (C) circle (1pt) node[below right] {$C$};
     \filldraw (P1) circle (1pt) node[below right]  {};
     \filldraw (I) circle (1pt) node[below] {$I$}; 
     \filldraw (G) circle (1pt) node[below]  {$G$};
    \end{tikzpicture}
  \end{center}
  \begin{itemchoice}
   \itemch Đúng. Vì $MN$ là đường trung bình của tam giác $SAD$ nên $MN \parallel AD$\\ 
   $\Rightarrow MN \parallel BC  \Rightarrow MN \parallel (SBC)$. (1)
   \itemch Đúng. Vì tương tự ta có $O,\,N$ theo thứ tự là trung điểm của $BD,\,SD$ nên $ON$  là đường trung bình của tam giác $SBD \Rightarrow ON \parallel SB \Rightarrow ON \parallel (SBC)$. (2)\\
   Từ (1) và (2) suy ra  $(OMN) \parallel (SBC)$.
   \itemch Sai. Vì ta có $OE$ là đường trung bình của tam giác $ABD$  nên $OE \parallel AD \Rightarrow OE \parallel MN$.\\
   Do đó $E \in (OMN)$. Mặt khác $F \in ON$, $ON \subset (OMN) \Rightarrow F \in  (OMN)$.\\
   Ta có: $\heva{&EF \subset (OMN)\\&(OMN) \parallel (SBC)} \Rightarrow  EF \parallel (SBC)$.
   \itemch Sai. Vì $G$ thuộc mặt phẳng $(ABCD)$ và cách đều $AB,\,CD$ nên $G$ thuộc đường trung bình của hình bình hành  $ABCD$ (ứng với hai cạnh $AB,\,CD$).\\
   Gọi $I$ là trung điểm $BC$  thì $I,O,G$ thẳng hàng.\\
   Ta có $OI$ là đường trung bình của $\Delta  ABC$ nên $OI \parallel AB \Rightarrow OI \parallel (SAB)$. (3)\\
    Tương tự, ta có $ON \parallel SB \Rightarrow ON \parallel (SAB)$. (4)\\
   Từ (3) và (4) suy ra $(OIN) \parallel (SAB)$  mà $NG \subset (OIN)$ nên $NG \parallel (SAB)$.
  \end{itemchoice}
 } 
\end{ex}

% nguồn: Ok_dungsai_lop11_Lan2
\begin{ex}%[1H4H4-6]%[ma:c78a911a]
	Cho hình hộp $ABCD\cdot A'B'C'D'$  có các cạnh $AA'$, $BB'$, $CC'$, $DD'$ song song với nhau.
	 \choiceTF
		 {\True $(BDA')\parallel(B'D'C)$}
		{\True Đường chéo $AC'$ đi qua trọng tâm $G_1$, $G_2$ của tam giác $BDA'$ và $B'D'C$} 
		{$AG_1=2G_1  G_2$}
		{\True Mặt phẳng $(A'B'G_2)$ cắt hình hộp $ABCD\cdot A'B'C'D'$ tạo thành  một tứ giác là hình bình hành}
	\loigiai{ 
		\begin{center}
			\begin{tikzpicture}[>=stealth,line join=round,line cap=round,font=\footnotesize,scale=1]
					 \def\g{95}
					\def\h{5} 
					\path 
						(0,0) coordinate (A) 
						++(\g:\h)  coordinate (A')
						(0:6) coordinate (D) 
						++(\g:\h)  coordinate (D')
						 (-130:3) coordinate (B)
						++(\g:\h) coordinate  (B')
						($(B)+(D)-(A)$) coordinate (C) 
						 ++(\g:\h) coordinate (C')
						($(A)!0.5!(C)$) coordinate  (O)
						++(\g:\h) coordinate (O') 
						($(A)!0.5!(B)$) coordinate (H) 
						 ($(C')!0.5!(C)$) coordinate (M)
						 ($(D)!0.5!(D')$) coordinate (N)
						($(M)!0.5!(N)$)  coordinate (P)
						($(A)!1/3!(C')$)  coordinate (G_1)
						($(A)!2/3!(C')$) coordinate (G_2); 
					 \draw
						 (B)--(C)--(D)
						(B)--(B') 
						(C)--(C') 
						 (D)--(D')
						(D')--(B')--(C')--(D')--(A')--(B') 
						(A')--(C') 
						(B')--(M)--(N) 
						 (B')--(C)--(D');
					 \foreach \i in{N,D,G_2,O,A,B}
					\draw[dashed] 
						(A')--(\i) 
						(B)--(A)--(D)--(B) 
						(A)--(C)--(O') 
						 (A)--(C')
						 (B')--(P);
					 \foreach \m/\n in{A'/90,B'/90,C'/90,D'/90,A/180,B/180,C/0,D/0,M/0,N/0,P/-60,O/-90,O'/90,G_1/-20,G_2/-100}
					\fill[black] (\m)  circle (1pt)+(\n:4mm) node{$\m$};
				\end{tikzpicture}
		\end{center}
		\begin{itemchoice}
			\itemch \textbf{Đúng}.\\ 
			 Vì
			\begin{itemize}
				\item $BA'D'C$ là hình bình hành nên  $BA'\parallel D'C$.
				 \item $BB'D'D$ là hình bình hành nên $B'D'\parallel BD$.
			\end{itemize}
			Ta  có
			$\heva{&BA'\parallel D'C \subset (B'D'C)\\&BD\parallel B'D' \subset (B'D'C)\\&BA'\cap BD=B\\&BA',BD\subset (BDA')}\Rightarrow (BDA')\parallel  (B'D'C)$.
			 \itemch \textbf{Đúng}.\\
			Trong $(ABCD)$, gọi $O=AC\cap BD$.\\ 
			Suy ra $\heva{&O \in A C\\& A  C \subset(ACC'A') \\& O \in B D\\&B D \subset(BDA')} \Rightarrow O \in(ACC'A') \cap (BDA')$.\\
			Mà  $A'\in(ACC'A')\cap(BDA')$ nên $A'O=(ACC'A')\cap(BDA')$.\\
			Trong $(ACC'A')$,  gọi $E=A'O\cap AC'$.\\
			Ta có $\dfrac{A'E}{EO}=\dfrac{A'C'}{AO}=2$ (Thales).\\ 
			Suy ra $E$ trùng với trọng  tâm $G_1$ của tam giác $BDA'$.\\
			Hay $AC'$ đi qua trọng tâm $G_1$ của  $\triangle BDA'$.\hfill(1)\\
			 Trong $(A'B'C'D')$, gọi $O'=A'C'\cap B'D'$.\\
			Chứng  minh tương tự ta được $AC'$ đi qua trọng tâm $G_2$ của tam giác $B'D'C$. \hfill(2)\\
			Từ (1) và (2) suy ra $AC'$ qua trọng tâm $G_1$, $G_2$ của tam  giác $BDA'$ 'và $B'D'C$.
			 \itemch \textbf{Sai}.\\
			Trong $(AA'C'C)$, gọi $I=A'C\cap  AC'$.
			\begin{center}
				\begin{tikzpicture}[line join=round,line cap=round,>=stealth,font=\footnotesize,scale=1]
						\path 
						(0,0) coordinate (A) 
						(0:5) coordinate  (C)
						(60:3)  coordinate (A')
						($(A')+(C)-(A)$) coordinate  (C')
						($(A)!0.5!(C)$) coordinate (O) 
						($(A')!0.5!(C')$) coordinate (O') 
						($(A')!0.5!(C)$)  coordinate (I)
						($(A)!1/3!(C')$) coordinate  (G_1)
						($(A)!2/3!(C')$) coordinate  (G_2);
						 \draw (A)--(C)--(C')--(A')--(A)--(C')
						(O)--(A')--(C)--(O'); 
						\foreach \i/\g in {A/-90,C/-90,C'/90,A'/90,I/-90,G_1/-30,G_2/-30,O/-90,O'/90}{\draw[fill=black](\i) circle  (1pt)
						($(\i)+(\g:4mm)$) node[scale=1]{$\i$};} 
					\end{tikzpicture}
			\end{center}
			Ta có $G_1$ là trọng tâm của  tam giác $A'AC$ nên $AG_1=\dfrac{2}{3}AI=\dfrac{2}{3}\cdot\dfrac{1}{2}AC'=\dfrac{1}{3}AC'$.\\
			Ta có $G_2$ là trọng tâm của tam  giác $A'CC'$ nên $AG_2=\dfrac{2}{3}C'I=\dfrac{2}{3}\cdot\dfrac{1}{2}AC'=\dfrac{1}{3}AC'$.\\
			 Ta có $G_1 G_2=AC'-AG_1-C'G_2=AC'-\dfrac{1}{3}AC'-\dfrac{1}{3}AC'=\dfrac{1}{3}AC'$.\\
			Từ  đó ta có $AG_1=G_1 G_2=G_2 C'\left(=\dfrac{1}{3}AC'\right)$.\\
			Suy ra $G_1$, $G_2$ chia đoạn $AC'$ ' thành ba  phần bằng nhau
			\itemch \textbf{Đúng}.\\ 
			Xét tam  giác $B'D'C$, gọi $P=B'G_2\cap CD'$.\\
			Suy ra $P$ là trung điểm của  $CD'$.\\
			Trong $(A'B'G_2)$, vẽ $Px\parallel A'B'$.\\ 
			Trong $(CC'D'D)$, $Px$ cắt $CC'$, $DD'$ lần lượt tại $M$  và $N$.\\
			Suy ra $MN\parallel  CD$ và $MN=CD$.\\
			Ta có $\heva{&(A'B'G_2) \cap(A'B'C'D')=A'B'\\& (A'B'G_2) \cap(B'C'CB)=B'M \\& (A'B'G_2)  \cap(C C'D'D)=MN \\& (A'B'G_2) \cap(A A'D'D)=N A'}$
			suy ra thiết diện của $(A'B'G_2)$  và hình hộp $ABCD.A'B'C'D'$ là tứ giác $A'B'MN$.\\
			Ta  có $\heva{&MN\parallel A'B'\\&MN=A'B'(=CD)}$ suy ra thiết diện $A'B'MN$ là hình bình hành.
		\end{itemchoice}
	} 
\end{ex}

% nguồn: Ok_dungsai_lop11_Lan2
\begin{ex}%[1H4H4-6]%[ma:e6b06bff]
	Cho hình chóp $S.ABCD$ có đáy $ABCD$ là hình bình hành tâm $O$. Gọi $M$, $N$ lần lượt là  trung điểm của $SA$ và $SD$.
	 \choiceTF
		{$ON$ chéo  nhau với $SB$}
		 {\True $(OMN)\parallel(SBC)$}
		{Gọi $P$  và $Q$ là trung điểm của $AB$ và $ON$. Khi đó $PQ$ cắt $(SBC)$}
		{\True Gọi $R$ là trung điểm  $AD$. Khi đó $(MOR)\parallel(SCD)$}
	 \loigiai{
		\begin{center}
			\begin{tikzpicture}[>=stealth,line join=round,line cap=round,font=\footnotesize,scale=1]
					\path 
						(0,0) coordinate  (A)
						(0:6) coordinate (D) 
						(-140:4) coordinate (B) 
						($(B)+(D)-(A)$) coordinate (C) 
						 (70:5) coordinate (S)
						($(A)!0.5!(C)$) coordinate (O) 
						($(A)!0.5!(S)$) coordinate  (M)
						($(S)!0.5!(D)$) coordinate (N) 
						($(A)!0.5!(B)$) coordinate (P) 
						($(O)!0.5!(N)$) coordinate  (Q)
						($(A)!0.5!(D)$) coordinate (R); 
					\draw 
						 (S)--(B)--(C)--(D)--(S)--(C);
					\draw[dashed]  (S)--(A)--(B)--(D)--(A)--(C)
						(P)--(M)--(N)--(O)--(P)--(Q) 
						(M)--(O)--(R)--(M); 
					 \foreach \m/\n in{S/90,M/30,N/30,D/0,C/-90,B/-90,A/180,O/-90,P/-90,Q/0,R/50}
					 \fill[black] (\m) circle (1pt)+(\n:3mm) node{$\m$};
				\end{tikzpicture}
		\end{center}
		\begin{itemchoice}
			 \itemch \textbf{Sai}.\\
			Do 
			\begin{itemize}
				 \item $O$ là trung điểm của $BD$.
				\item $N$ là trung điểm của $SD$. 
			\end{itemize}
			Suy ra $ON$ là đường trung bình của $\triangle BDS$.\\ 
			Suy ra  $ON \parallel SB$.
			 \itemch \textbf{Đúng}.\\
			 Ta có
			\begin{itemize}
				\item $OM\parallel SC$  (đường trung bình tam giác $SAC$).
				\item $ON\parallel  SB$ (đường trung bình tam giác $SBD$).
			\end{itemize}
			Vì $\heva{&ON\parallel SB;OM\parallel SC\\&OM, ON\subset(OMN)\\& OM\cap ON=O\\&SB, SC\subset(SBC)\\&  SB\cap SC=S}
			\Rightarrow  (OMN)\parallel (SBC)$.
			\itemch  \textbf{Sai}.\\
			Gọi $P$ và $Q$ lần lượt là  trung điểm của $AB$ và $ON$.\\
			Ta có $\heva{&OP\parallel AD\\&OP=\dfrac{1}{2}AD\\&AD\parallel MN\\&MN=\dfrac{1}{2}AD} \Rightarrow \heva{&OP\parallel MN\\&OP=MN} \Rightarrow OMPN$ là hình  bình hành $ \Rightarrow P\in(OMN)$.\\
			Ta có  $\heva{&PQ\subset(OMN)\\&(OMN)\parallel(SBC)}\Rightarrow PQ\parallel(SBC)$.
			 \itemch \textbf{Đúng}.\\
			Gọi $R$ là  trung điểm $AD$.\\
			Ta có $OR\parallel CD$ (đường  trung bình của tam giác $ACD$).\\
			Vì $\heva{&OM\parallel SC\\&OR\parallel CD\\&OM, OR\subset(MOR)\\&OM\cap OR=O\\&SC,  SD\subset(SCD)\\&SC\cap SD=S} \Rightarrow (MOR)\parallel(SCD)$.
		\end{itemchoice}
	 }
\end{ex}

% nguồn: Ok_dungsai_lop11_Lan2
\begin{ex}%[1H4H4-6]%[ma:08e2e041]
	 Cho hình chóp $S.ABCD$ có đáy là hình bình hành tâm $O$. Gọi $M$, $N$ lần lượt là trung điểm của $SA$ và $SD$.
	 \choiceTF
		 {\True $MN\parallel(SBC)$}
		{\True $(OMN)\parallel(SBC)$} 
		{Gọi $E$ là trung điểm đoạn $AB$ và $F$ là một  điểm thuộc đoạn $ON$. Khi đó $EF$ cắt với mặt phẳng $(SBC)$}
		{Gọi $G$ là một điểm trên mặt phẳng $(ABCD)$ cách đều $AB$ và $CD$. Khi đó $GN$ cắt  $(SAB)$}
	\loigiai{ 
		\begin{center}
			\begin{tikzpicture}[line join=round,line cap=round,>=stealth,font=\footnotesize,scale=1]
					\path 
						(0,0) coordinate  (A)
						(-135:3.5) coordinate  (B)
						(0:6) coordinate (D) 
						($(B)+(D)-(A)$) coordinate (C) 
						(80:4) coordinate (S) 
						($(S)!0.5!(A)$) coordinate  (M)
						($(S)!0.5!(D)$)  coordinate (N)
						($(A)!0.5!(B)$) coordinate  (E)
						($(A)!0.5!(C)$) coordinate  (O)
						($(B)!0.5!(C)$) coordinate  (I)
						($(A)!0.5!(D)$) coordinate  (K)
						($(O)!3/4!(N)$) coordinate (F) 
						($(I)!4/5!(K)$) coordinate  (G);
					\draw  (S)--(B)--(C)--(D)--(S)--(C);
					\draw[dashed,thin]  (S)--(A)--(B)
						 (C)--(A)--(D)
						 (M)--(N)--(O)--(E)--(M)--(O)
						(E)--(F) 
						(I)--(K) 
						 (N)--(G)
						 (B)--(D);
					\foreach \i/\g in {S/90,A/60,B/-90,C/-90,D/60,M/180,N/60,F/0,O/-90,E/-90,I/-90,G/-90}{\draw[fill=black](\i)  circle (1pt)
						 ($(\i)+(\g:3mm)$) node[scale=1]{$\i$};}
				\end{tikzpicture}
		\end{center}
		\begin{itemchoice}
			\itemch \textbf{Đúng}.\\ 
			Do $M$, $N$ lần lượt là trung điểm của $SA$ và $SD$ nên $MN$ là đường trung bình của tam giác $SAD$. Suy ra $MN\parallel AD \Rightarrow  MN\parallel BC\subset (SBC) \Rightarrow MN\parallel(SBC)$.
			 \itemch \textbf{Đúng}.\\
			Tương tự, ta có $O$,  $N$ theo thứ tự là trung điểm của $BD$, $SD$ nên $ON$ là đường trung bình của tam giác $SBD \Rightarrow ON\parallel SB\subset (SBC) \Rightarrow ON\parallel(SBC)$.\\
			Do $\heva{&MN\parallel(SBC)\\&ON\parallel(SBC)\\&MN\cap ON=N\\&MN,ON\subset (OMN)}\Rightarrow  (OMN)\parallel(SBC)$.
			 \itemch \textbf{Sai}.\\
			Ta có $OE$ là  đường trung bình của tam giác $ABD$ nên $OE\parallel AD \Rightarrow OE\parallel MN$. Do đó $E\in(OMN)$.\\
			Mặt khác $F\in ON$, $ON\subset(OMN) \Rightarrow F\in(OMN)$.\\ 
			Vì $\heva{&EF\subset(OMN)\\&(OMN)\parallel(SBC)} \Rightarrow  EF\parallel(SBC)$.
			\itemch \textbf{Sai}.\\ 
			Vì  $G$ thuộc mặt phẳng $(ABCD)$ và cách đều $AB$, $CD$ nên $G$ thuộc đường trung bình của hình bình hành $ABCD$ (ứng với hai cạnh $AB$, $CD$).\\
			Gọi  $I$ là trung điểm $BC$ thì $I$, $O$, $G$ thẳng hàng.\\
			Ta có $OI$  là đường trung bình của $\triangle ABC$ nên $OI\parallel AB \Rightarrow OI\parallel(SAB)$.\\
			Tương tự, ta  có $ON\parallel SB \Rightarrow ON\parallel(SAB)$.\\
			Ta có $\heva{&ON\parallel(SAB)\\&OI\parallel(SAB)\\&ON\cap OI=O\\&OI,ON \subset (OIN)}$ suy ra $(OIN)\parallel(SAB)$ mà $NG\subset(OIN)$ nên $NG\parallel(SAB)$. 
		\end{itemchoice}
	 }
\end{ex}

% nguồn: Ok_dungsai_lop11_Lan2
\begin{ex}%[1H4H4-6]%[ma:079fea50]
	Cho lăng trụ tam giác $ABC\cdot A'B'C'$ có $I$, $K$,  $G$ lần lượt là trọng tâm các tam giác $ABC$, $A'B'C'$, $ACC'$. Gọi $M$, $M'$ và $N$ lần lượt là trung điểm của $BC$, $B'C'$ và $CC'$.
	 \choiceTF
		{\True $AMM'A'$ là  hình bình hành}
		{$\dfrac{AI}{AM}=\dfrac{AG}{AN}=\dfrac{1}{3}$} 
		{$(IKG)$  cắt $(BCC'B')$}
		{\True  $(A'KG)\parallel(AIB')$}
	 \loigiai{
		\begin{center}
			\begin{tikzpicture}[line join=round,line cap=round,>=stealth]
					\def\g{80} 
					\def\h{5} 
					\path 
						(0,0)  coordinate (A)
						 (-50:3.5) coordinate (B)
						(0:7) coordinate  (C)
						(A)++(\g:\h) coordinate (A') 
						 (B)++(\g:\h) coordinate (B')
						 (C)++(\g:\h) coordinate (C')
						 ($(C)!0.5!(B)$) coordinate (M)
						($(A)!2/3!(M)$)  coordinate (I)
						(M)++(\g:\h) coordinate (M') 
						(I)++(\g:\h) coordinate (K) 
						($(C)!0.5!(C')$) coordinate  (N)
						 ($(A)!2/3!(N)$) coordinate (G);
					\draw (A')--(A)--(B)--(C)--(C')--(B')--(A')--(C') 
						 (B)--(B')
						 (M)--(N)
						(A')--(M')--(M) 
						 (M')--(C)
						 (A)--(B')--(M);
					 \draw[dashed](A)--(C)
						(A)--(M) 
						(A)--(N) 
						 (I)--(K)
						(I)--(G)--(K) 
						 (A')--(G)
						(B')--(I) 
						 (C)--(G);
					\foreach \i/\g  in {A/180,B/-90,C/0,A'/90,B'/-170,C'/90,I/-90,K/90,M'/90,N/0,G/170,M/-60}{
						\draw[fill=black] (\i)  circle (1pt)
						($(\i)+(\g:4mm)$) node[scale=1]  {$\i$};
					 }
				\end{tikzpicture}
		\end{center}
		\begin{itemchoice}
			\itemch  \textbf{Đúng}.\\
			Gọi  $M$, $M'$ lần lượt là trung điểm của $BC$, $B'C'$.\\
			Suy ra $MM'$ là đường trung bình của  hình bình hành $BCC'B'$ nên
			\[\heva{&MM' \parallel BB'  \\&MM'=BB'} \Rightarrow \heva{&MM'\parallel AA' \\&MM'=AA'.}\]
			Suy ra $ AMM'A'$  là hình bình hành.
			 \itemch \textbf{Sai}.\\
			Vì $I$ và $G$ lần lượt  là trọng tâm các tam giác $ABC$, $ACC'$ nên
			 \[\dfrac{AI}{AM}=\dfrac{AG}{AN}=\dfrac{2}{3}.\]
			 \itemch \textbf{Sai}.\\
			Vì $I$, $K$ theo thứ  tự là trọng tâm các tam giác $ABC$, $A'B'C'$ nên
			\[IM=KM'=\dfrac{1}{3}A'M'=\dfrac{1}{3}AM.\] 
			Mà  $IM\parallel KM'$ nên $IKM'M$ là hình bình hành.\\
			 Suy ra $\heva{&IK\parallel MM'\\&MM'\subset(BCC'B')} \Rightarrow IK\parallel(BCC'B')$.\\
			Lại có $\heva{&IG\parallel  MN\\&MN\subset(BCC'B')}$ nên $IG\parallel(BCC'B')$.\\
			Ta có $\heva{&IK\parallel(BCC'B')\\&IG\parallel(BCC'B')\\&IK \cap  IG=I\\&IK,IG\subset (IKG)}$ suy ra $(IKG)\parallel(BCC'B')$.
			\itemch \textbf{Đúng}.\\ 
			Vì $(A'KG)\equiv(A'M'C)$, $(AIB')\equiv(AMB')$,  ta cần chứng minh
			\[(A' M' C)\parallel  (AMB').\]
			Vì $AMM'A'$ là hình  bình hành nên $AM\parallel A'M'$ mà $A'M'\subset(A'M'C)$ nên $AM\parallel(A'M'C)$.\\
			Ta có $\heva{&CM\parallel  B'M'\\&CM=\dfrac{1}{2}BC=\dfrac{1}{2}C'B'=B'M'}$ suy ra $CMB'M'$ là hình bình hành nên $B'M\parallel CM'$.\\
			Mà $CM'\subset(A'M'C)$ nên  $ B'M\parallel(A'M'C)$.\\
			Ta có $\heva{&AM\parallel(A'M'C)\\&B'M\parallel(A'M'C)\\&AM\cap B'M=M\\&AM,  B'M\subset (AMB')}$ suy ra $(A'M'C)\parallel(AMB')$, hay $(A'KG)\parallel(AIB')$.
		\end{itemchoice}
	 }
\end{ex}

% nguồn: Ok_dungsai_lop11_Lan2
\begin{ex}%[1H4H4-6]%[ma:4436f785]
	Cho lăng trụ tam giác $ABC\cdot A'B'C'$. Gọi $I$ và $I'$ lần lượt là  trung điểm của $BC$ và $B'C'$.
	 \choiceTF
		{\True $II'\parallel BB'$} 
		{\True $AA'I'I$ là  hình bình hành}
		{$IA'$ song song  $(AB'C')$}
		{Giao tuyến của $(AB'C')$ và $(A'BC')$ là đường thẳng đi  qua giao điểm của hai đường thẳng $AI'$, $A'I$}
	 \loigiai{
		\begin{center}
			\begin{tikzpicture}[line join=round,line cap=round,>=stealth,font=\footnotesize,scale=1]
					\def\h{4} 
					 \path
						(0,0)  coordinate (A)
						(0:6) coordinate  (C)
						(-40:5) coordinate (B) 
						($(B)!0.5!(C)$) coordinate (I); 
					\foreach \m in {A,B,C,I}{\coordinate (\m') at ($(\m)+(90:\h)$);} 
					\path 
						($(A)!0.5!(B')$) coordinate  (F)
						($(A)!0.5!(I')$) coordinate (E); 
					 \draw (A')--(B')--(C')--(A')--(A)--(B)--(C)--(C')--(B)--(A')--(I')--(I)
						 (A)--(B')--(B);
					\draw[dashed] (A)--(C) 
						(A)--(I)--(A') 
						(A)--(I') 
						(A)--(C'); 
					\foreach \i/\g  in {A'/90,B'/90,C'/90,A/180,B/-90,C/0,I/-90,I'/90,F/-90,E/-90}{\draw[fill=black](\i) circle (1pt)
						($(\i)+(\g:3mm)$)  node[scale=1]{$\i$};}
				\end{tikzpicture}
		\end{center}
		\begin{itemchoice}
			\itemch  \textbf{Đúng}.\\
			Ta có $I'$, $I$ là trung điểm của  $B'C'$ và $BC$.\\
			Suy ra $II'$ là đường trung bình của  hình bình hành $BB'C'C$.\\
			Suy ra $II'=BB'$  và $II'\parallel BB'$.
			\itemch  \textbf{Đúng}.\\
			Ta có $\heva{&II'\parallel AA'(\parallel BB')\\&II'=AA'(=BB')}$  suy ra $AA'I'I$ là hình bình hành.
			 \itemch \textbf{Sai}.\\
			Trong  $(IAA'I')$, gọi $E=AI'\cap A'I$.\\
			Ta có $\heva{&E\in  AI';AI'\subset(AB'C')\\&E\in A'I}$ suy ra $E=A'I\cap(AB'C')$.
			\itemch \textbf{Sai}.\\ 
			Trong $(AA'B'B)$,  gọi $F = AB' \cap A'B$. \\
			$\Rightarrow \heva{&F \in AB'; AB' \subset (ABC') \\ &F \in A'B; A'B \subset (A'BC')} \Rightarrow  F \in (ABC') \cap (A'BC')$. \hfill(1) \\
			Ta có $E = AI' \cap  A'I$. \\
			$\Rightarrow \heva{&C' \in (ABC') \\ &C' \in (A'BC')} \Rightarrow C' \in (ABC')  \cap (A'BC')$. \hfill(2) \\
			Từ $(1)$ và $(2)$ suy ra $C'F = (ABC')  \cap (A'BC')$.\\
			Ta thấy $\heva{&EF\parallel  C'B'\\&C'F\not\parallel C'B'} \Rightarrow E$, $F$, $C'$ không thẳng hàng.\\
			Hay giao tuyển của hai mặt  phẳng không đi qua $E$.
		\end{itemchoice}
	 }
\end{ex}

% nguồn: Ok_dungsai_lop11_Lan2
\begin{ex}%[1H4H4-6]%[ma:5255b990]
	Cho hình bình  hành $ABCD$ và $ABEF$ nằm ở hai mặt phẳng khác nhau. Gọi $M$ là trọng tâm $\triangle ABE$. Gọi $(P)$ là mặt phẳng đi qua $M$ và song song với mặt $(ADF)$. Lấy $N$ là giao điểm của $(P)$ và $AC$.
	 \choiceTF
		 {\True $EFDC$ là hình thang}
		 {\True $FD\parallel EC$}
		{\True  $(ADF)\parallel(BCE)$}
		{$\dfrac{AN}{NC}=3$} 
	 \loigiai{
		\begin{center}
			\begin{tikzpicture}[>=stealth,line join=round,line cap=round,font=\footnotesize,scale=1]
					\path 
						(0,0)  coordinate (A)
						 (0:5) coordinate (D)
						(-130:3) coordinate  (B)
						 ($(B)+(D)-(A)$) coordinate (C)
						 (B)++(70:4) coordinate (E)
						($(A)+(E)-(B)$) coordinate  (F)
						($(A)!2/3!(B)$)  coordinate (I)
						 ($(F)!2/3!(E)$) coordinate (J)
						($(D)!2/3!(C)$)  coordinate (K)
						($(B)!0.5!(E)$) coordinate (Q) 
						 ($(A)!2/3!(Q)$) coordinate (M)
						(intersection of I--K and A--C)  coordinate (N);
					\draw (E)--(B)--(C)--(E)--(F)--(D)--(C) 
						 (J)--(K);
					\draw[dashed] (F)--(A)--(B) 
						 (A)--(D)
						(J)--(I)--(K) 
						 (C)--(A)--(E)
						(A)--(Q); 
					 \foreach \m/\n in{E/90,F/90,A/-90,B/180,C/0,D/0,M/-50,Q/180,J/90,K/0,N/-90,I/-90}
					\fill[black]  (\m) circle (1pt)+(\n:3mm) node{$\m$};
				\end{tikzpicture}
		\end{center}
		\begin{itemchoice}
			\itemch \textbf{Đúng}.\\ 
			Do $ABCD$ và $ABEF$ là các hình bình hành nên $\heva{&EF\parallel CD&(\parallel AB)\\&EF=CD&(=AB)}$ suy ra $EFDC$  là hình bình hành.
			\itemch  \textbf{Đúng}.\\
			Do $EFDC$ là hình bình hành  nên $FD\parallel EC$.
			 \itemch \textbf{Đúng}.\\
			Do $\heva{&AD\parallel  BC\subset (BCE)\\&AF\parallel BE\subset (BCE)\\&AD, AF\subset(ADF)\\&AD\cap AF=A} \Rightarrow (ADF)\parallel(BCE)$.
			 \itemch \textbf{Đúng}.\\
			Vẽ mp$(P)$ chứa  $M$ và $(P)\parallel(ADF)$.
			\begin{itemize}
				\item Trong  $(ABEF)$ kẻ $IJ\parallel AF$ qua $M$ với $I\in AB$, $J\in EF$.
				\item Từ  $I$ kẻ $IK\parallel AD$ với $K\in CD$.
				\item Nối $J$ với  $K$.
			\end{itemize}
			Ta được $(P)$ cắt $AB$, $AC$, $CD$, $EF$ lần  lượt tại $I$, $N$, $K$, $J$.\\
			Gọi $Q$ là trung điểm  của $BE$, mà $M$ là trọng tâm $\triangle ABE$ nên $AM=2MQ$.\\
			Mà $MI\parallel AF\parallel BE$ suy ra $\dfrac{AM}{QM}=\dfrac{AI}{IB}=2$.\hfill(1)\\ 
			 Lại có $\dfrac{AI}{BI}=\dfrac{AN}{NC}$ do $IN\parallel BC$.\hfill(2)\\
			Từ (1) và (2) suy  ra $\dfrac{AN}{NC}=2$.
		\end{itemchoice}
	 }
\end{ex}

% nguồn: Ok_dungsai_lop11_Lan2
\begin{ex}%[1H4H4-6]%[ma:f253461e]
	Cho hình hộp $ABCD.A'B'C'D'$. Gọi $G_1$, $G_2$ là  trọng tâm của các tam giác $A'BD$, $B'D'C$.
	 \choiceTF
		{\True $A'D'CB$  là hình bình hành}
		 {\True $(A'BD)\parallel(B'D'C)$}
		{\True  $G_1$, $G_2$ cùng thuộc $AC'$}
		 {$G_1 G_2=\dfrac{2}{3}AC'$}
	\loigiai{ 
		\begin{center}
			\begin{tikzpicture}[>=stealth,line join=round,line cap=round,font=\footnotesize,scale=1]
					 \def\h{4cm}
					 \path
						(0,0) coordinate (A) 
						++(-90:\h)  coordinate (A')
						(0:5)  coordinate (B)
						++(-90:\h) coordinate  (B')
						(-120:3)  coordinate (D)
						++(-90:\h) coordinate (D') 
						($(B)+(D)-(A)$)  coordinate (C)
						++(-90:\h)  coordinate (C')
						($(A)!0.5!(C)$) coordinate (O) 
						 ++(-90:\h) coordinate (O')
						($(A)!0.5!(C')$)  coordinate (I)
						($(A)!1/3!(C')$) coordinate  (G_1)
						($(A)!2/3!(C')$) coordinate (G_2); 
					 \draw (A)--(B)--(C)--(D)--cycle
						 (D)--(D')--(C)--(C')
						 (D)--(B)--(B')
						(D')--(C')--(B')--(C)--(A); 
					\foreach \i  in{A,O,D,B,C,C',D',B'}
					 \draw[dashed] (A')--(\i)
						(C)--(O') 
						(B')--(D') 
						 (A)--(C');
					 \foreach \m/\n in{A/90,B/90,C/0,D/180,A'/180,B'/0,C'/0,D'/180,O/90,O'/-90,I/-90,G_1/180,G_2/0}
					\fill[black] (\m) circle  (1pt)+(\n:3mm) node{$\m$};
				\end{tikzpicture}
		\end{center}
		\begin{itemchoice}
			 \itemch \textbf{Đúng}.\\
			Vì $ABCD.A'B'C'D'$ là  hình hộp nên $\heva{&A'D'\parallel BC\\&A'D'=BC} \Rightarrow A'D'CB$ là hình bình hành.
			 \itemch \textbf{Đúng}.\\
			Vì $A'D'CB$ là hình bình hành  suy ra $A'B\parallel CD'\subset(B'D'C) \Rightarrow A'B\parallel(B'D'C)$.\\
			Tương tự, ta có $\heva{&A'B'\parallel CD\\&A'B'=CD} \Rightarrow A'B'CD$ là hình bình hành.\\  Suy ra $A'D\parallel B'C\subset(B'D'C) \Rightarrow A'D\parallel(B'D'C)$.\\
			Ta có $\heva{&A'B\parallel(B'D'C)\\&A'D\parallel(B'D'C)\\&A'B\cap  A'D=A'\\&A'B, A'D\subset (A'BD)}$ suy ra $(A'BD)\parallel(B'D'C)$.
			\itemch \textbf{Đúng}.\\ 
			Gọi $O$, $O'$, $I$ theo thứ tự là tâm của các hình  bình hành $ABCD$, $A'B'C'D'$, $ACC'A'$.\\
			Vì  $G_1$ là trọng tâm tam giác $AB'D$ nên $\dfrac{A'G_1}{A'O}=\dfrac{2}{3}$.\\
			Suy ra $G_1$ là trọng tâm tam giác $A'AC$, suy  ra $G_1=AI\cap A'O$. \hfill(1)\\
			Tương tự, $G_2$ là trọng  tâm tam giác $B'D'C$ nên $\dfrac{CG_2}{CO'}=\dfrac{2}{3}$.\\
			Suy ra $G_2$ là trọng  tâm tam giác $A'C'C$, suy ra $G_2=C'I\cap CO'$.\hfill(2)\\
			Từ (1) và  (2) suy ra $G_1$, $G_2$ cùng thuộc $AC'$.
			 \itemch \textbf{Sai}.\\
			Ta có 
			\begin{itemize}
				\item $\dfrac{AG_1}{AI}=\dfrac{2}{3}  \Rightarrow \dfrac{AG_1}{AC'}=\dfrac{1}{3}$;
				 \item $\dfrac{C'G_2}{C'I}=\dfrac{2}{3}$.
			\end{itemize}
			Suy  ra $\dfrac{C'G_2}{AC'}=\dfrac{1}{3}$.\\
			 Do vậy $AG_1\doteq G_1 G_2=G_2 C'=\dfrac{1}{3}AC'$.
		\end{itemchoice}
	 }
\end{ex}

% nguồn: Ok_dungsai_lop11_Lan1
\begin{ex}%[1H4V4-6]%[ma:15e468fb]
 Cho hình chóp $S.ABCD$ có đáy $ABCD$ là hình bình hành. Gọi $H$, $I$, $K$ lần lượt là trung điểm của $SA$, $SB$, $SC$. Gọi $M$ là giao điểm của $AI$ và $KD$, $N$ là  giao điểm của $ DH$ và $CI$. Khi đó
  \choiceTF
  {\True $HI\parallel(ABCD)$} 
   {\True $(HIK) \parallel(ABCD)$}
  {$SM$ và $HI$ chéo nhau} 
  {$(SMN)$ cắt $(HIK)$} 
  \loigiai{
  \begin{center}
   \begin{tikzpicture}[scale=1,font=\footnotesize,line join=round,line cap=round,>=stealth]
      \def\bc{5}
      \def\ba{4}
      \def\h{7}
     \def\gocB{60} 
     \coordinate[label=below left:$B$] (B)  at (0,0);
     \coordinate[label=above  right:$A$] (A) at (\gocB:\ba);
      \coordinate[label=below:$C$] (C) at (\bc,0);
     \coordinate[label=right:$D$] (D) at ($(C)-(B)+(A)$); 
     \coordinate (O) at ($(A)!.5!(C)$); 
     \coordinate[label=above:$S$] (S)  at ($(O)+(105:\h)$);
      \coordinate[label=above right:$H$] (H) at ($(S)!0.5!(A)$);
      \coordinate[label=below left:$I$] (I) at ($(S)!0.5!(B)$);
     \coordinate[label=above right:$K$] (K) at ($(S)!0.5!(C)$); 
     \coordinate[label=below left:$M$] (M) at (intersection cs:first  line={(A)--(I)},second line={(K)--(D)});
      \coordinate[label=left:$N$] (N) at (intersection cs:first line={(D)--(H)},second line={(C)--(I)});
     \draw (B)--(C)--(D)--(S)--cycle  (S)--(C)
      (S)--(M)--(N)--cycle (I)--(M)--(D) 
      (C)--(I); 
     \draw[dashed]  (I)--(A)--(D)
       (S)--(A)--(B)
       (I)--(H)--(K)--cycle (D)--(N)--(I);
     \foreach \diem in {A,B,C,D,S,H,I,K,M,N} \fill (\diem) circle  (1pt);
    \end{tikzpicture}
  \end{center}
  \begin{itemize}
   \item a) b) Vì $HI$ là đường trung bình của tam giác $SAB$ nên $H I\parallel AB$, \\ 
   mà $AB \subset(ABCD) $  suy ra $ HI \parallel (ABCD)$. (1)\\
   Tương tự ta có $KI \parallel BC, BC \subset(ABCD) $ nên $KI \parallel (ABCD)$.  (2)\\
   Mặt khác $HI \subset(HKI),  KI \subset(HKI), HI \cap KI=I$. (3)\\
   Từ (1), (2) và  (3) suy ra $(HIK) \parallel(ABCD)$.\\
   \item  c) d)
   $$\begin{aligned}
   &  \text {Vì }\heva{&M \in AI, AI \subset(S AB) \\&
    M \in DK, DK \subset(SCD)
   } \Rightarrow M \in(SAB) \cap(SCD)\\ 
    & \Rightarrow SM=(SAB) \cap(SCD)
  \end{aligned}
   $$
   Khi đó
  $$\begin{aligned}
  & \heva{& (SAB)\cap (SCD)=SM \\ & AB\subset (SAB),CD\subset (SCD)}\Rightarrow SM\parallel AB\parallel CD\Rightarrow SM\parallel HI \\ & AB\parallel CD (1)  \\
  & \text{ Vì }\heva{& N\in DH,DH\subset (SAD) \\ & N\in CI,CI\subset (SBC)}\Rightarrow  N\in (SAD)\cap (SBC) \\
  & \Rightarrow SN=(SAD)\cap (SBC). 
 \end{aligned}$$
 Khi đó, ta  có
  $$
 \begin{aligned}
  \heva{&(SAD)\cap(SBC)=SN\\ & AD\subset(SAD),BC\subset(SBC)\\ & AD\parallel  BC}&\Rightarrow SN\parallel AD\parallel BC\Rightarrow SN\parallel KI\quad (2)
 \end{aligned}$$
 Mặt khác ba điểm $S$, $M$,  $N$ không thẳng hàng. (3) \\
 Từ (1), (2), (3) suy ra  $(SMN) \parallel(HIK)$.
 \end{itemize}
} 
\end{ex}

% nguồn: Ok_dungsai_lop11_Lan1
\begin{ex}%[1H4V4-6]%[ma:6b9f094c]
 Cho hai hình bình hành  $ABCD$ và $ABEF$ nằm ở hai mặt phẳng khác nhau. Gọi $M$ là trọng tâm $\triangle ABE$. Gọi $(P)$ là mặt phẳng đi qua $M$ và song song với mặt $(ADF)$. Lấy $N$ là giao điểm của $(P)$ và $AC$. Xét tính đúng sai của các mệnh đề sau.
 \choiceTF 
  {$EFDC$  là hình thang}
  {\True $FD\parallel  EC$}
  {\True  $(ADF) \parallel (BCE)$}
   {$\dfrac{AN}{NC}=3$}
 \loigiai{ 
  \begin{center}
   \begin{tikzpicture}[scale=1,font=\footnotesize,line join=round,line cap=round,>=stealth]
     \path (0,0) coordinate (B)++(3,0) coordinate (C)++(40:2.5) coordinate  (D)
      ($(B)+(D)-(C)$) coordinate (A)++(80:2) coordinate  (F)
      ($(F)+(B)-(A)$) coordinate (E) 
      (barycentric cs:A=1,B=1,E=1)  coordinate (M)
      (barycentric cs:A=1,B=2) coordinate  (K)
      (barycentric cs:E=1,B=1) coordinate (I) 
      (barycentric cs:A=1,C=2) coordinate (N); 
     \draw (B)--(C)--(D)--(F)--(E)--cycle (C)--(E)--(B); 
      \draw[dashed] (B)--(A)--(D)
       (F)--(A)--(E)
       (I)--(A)--(C)
      (M)--(K)--(N); 
     \foreach  \p/\g in {A/45,B/-120,C/-30,D/0,F/90,E/120,M/135,K/-90,I/180,N/45} \fill[black] (\p) circle (1pt)+(\g:0.3) node{$\p$};
    \end{tikzpicture}
  \end{center}
  \begin{itemchoice}
   \itemch Sai. Ta có $EF\parallel AB\parallel CD$  và $EF=AB=CD$ nên $EFDC$ là hình bình hành.
   \itemch Đúng. Vì $EFDC$ là hình  bình hành nên $FD\parallel EC$.
   \itemch Đúng.  Ta có $FD\parallel EC\Rightarrow FD\parallel (BCE)$.\\
    Tương tự, $AD\parallel BC\Rightarrow AD\parallel (BCE)$.\\
   Suy ra $(ADF)\parallel (BCE)$. 
   \itemch Sai. Gọi $I$ là trung điểm của $EB$ và  $K$ là giao điểm của $(P)$ với đường thẳng $AB$.\\
    Ta có $(P)\parallel (ADF)\parallel (BCE)$, $(P)\cap (ABEF)=MK$ và $(BEC)\cap (ABEF)=BE$ nên $MK\parallel AB$.\\
   Suy ra $\dfrac{AK}{AB}=\dfrac{AM}{AI}=\dfrac{2}{3}\Rightarrow \dfrac{AK}{KB}=2$.\\ 
   Ta có $(P)\parallel (ADF)\parallel (BCE)$, $(P)\cap (ABCD)=NK$ và  $(BEC)\cap (ABCD)=BC$ nên $NK\parallel BC$.\\
   Do đó,  $\dfrac{AN}{NC}=\dfrac{AK}{KB}=2$.
  \end{itemchoice}
  }
\end{ex}

% nguồn: Ok_dungsai_lop11_Lan2
\begin{ex}%[1H4V4-6]%[ma:bc68d4ea]
	Cho hình chóp $S.ABCD$ có  đáy $ABCD$ là hình bình hành tâm $O$. Gọi $M$, $N$ lần lượt là trung điểm của $SA$ và $SD$. Khi đó
	 \choiceTF
		{$ON$ chéo nhau  với $SB$}
		{\True  $\left(OMN\right)\parallel\left(SBC\right)$}
		{Gọi $P$ và $Q$ là  trung điểm của $AB$ và $ON$. Khi đó $PQ$ cắt $\left(SBC\right)$}
		{\True Gọi $R$ là trung điểm $AD$. Khi  đó $\left(MOR\right)\parallel\left(SCD\right)$}
	 \loigiai{
		\begin{center}
			\begin{tikzpicture}[scale=1,font=\footnotesize,line join=round,line cap=round,>=stealth]
					\path 
						 (0,0) coordinate (A)
						(-1.5,-1.5) coordinate (B) 
						(5,0) coordinate  (D)
						($(B)+(D)-(A)$) coordinate (C) 
						($(B)!0.5!(D)$)  coordinate (O)
						(1,3)  coordinate (S)
						($(S)!0.5!(A)$) coordinate (M) 
						($(S)!0.5!(D)$) coordinate  (N)
						 ($(B)!0.5!(A)$) coordinate (P)
						($(D)!0.5!(A)$) coordinate  (R)
						($(O)!0.5!(N)$)  coordinate (Q);
					\draw 
						(S)--(B)--(C)--(D)--(S)--(C); 
					\draw[dashed] 
						(S)--(A)--(C) 
						 (B)--(A)--(D)--cycle
						 (M)--(N)--(O)--(P)--(M)--(O)
						(Q)--(P) 
						 (M)--(R)--(O);
					 \foreach \p/\g in {A/170,B/-90,C/-90,D/0,S/90,M/50,N/10,Q/0,R/-80,P/-80,O/-90}
					\draw[fill=black] (\p) circle (1pt) node[shift=(\g:3mm)] {$\p$}; 
				\end{tikzpicture}
		\end{center}
		\begin{itemchoice}
			\itemch Ta  có $OM\parallel SC$ (đường trung bình tam giác $SAC$). Ta có $ON\parallel SB$ (đường trung bình tam giác $SBD$).
			\itemch Ta có $\heva{&ON\parallel SB;\,OM\parallel SC\\&OM,\,ON\subset (OMN),\,OM\cap ON=O\\&SB,\,SC\subset (SBC),\,SB\cap  SC=S} \Rightarrow\left(OMN\right)\parallel\left(SBC\right)$.
			\itemch Gọi $P$  và $Q$ lần lượt là trung điểm của $AB$ và $ON$. Chứng minh $PQ\parallel\left(SBC\right)$.\\
			Ta có  $\heva{&OP\parallel AB\\&AB\parallel MN}
			\Rightarrow OP\parallel MN\Rightarrow  OMPN$ là hình thang $\Rightarrow P\in\left(OMN\right)$.\\
			Ta có $\heva{&NP\subset\left(OMN\right)\\&\left(OMN\right)\parallel\left(SBC\right)} \Rightarrow  NP\parallel\left(SBC\right)$.
			\itemch Gọi $R$ là trung  điểm $AD$. Chứng minh $\left(MOR\right)\parallel\left(SCD\right)$.\\
			Ta có $OR\parallel CD$ (đường trung bình của tam giác  $ACD$).\\
			Ta có $\heva{&OM\parallel  SC\\&OR\parallel CD\\&OM,\,OR\subset (MOR),\,OM\cap OR=O\\&SC,\,SD\subset\left(SCD\right),\,SC\cap SD=S} \Rightarrow\left(MOR\right)\parallel\left(SCD\right)$.\\
		\end{itemchoice}
	 }
\end{ex}

% nguồn: Ok_dungsai_lop11_Lan2
\begin{ex}%[1H4V4-6]%[ma:f56492e6]
	Cho hình chóp $S.ABCD$  có đáy là hình bình hành tâm $O$. Gọi $M$, $N$ lần lượt là trung điểm của $SA$ và $SD$. Khi đó
	\choiceTF 
		{\True  $MN\parallel\left(SBC\right)$}
		{\True $\left(OMN\right)\parallel\left(SBC\right)$} 
		{Gọi $E$ là trung điểm đoạn $AB$ và $F$ là một điểm thuộc đoạn $ON$. Khi đó $EF$ cắt với mặt phẳng  $\left(SBC\right)$}
		{Gọi $G$ là một điểm trên mặt phẳng $\left(ABCD\right)$  cách đều $AB$ và $CD$. Khi đó $GN$ cắt $\left(SAB\right)$}
	 \loigiai{
		\begin{itemchoice}
			\itemch Vì $MN$ là đường trung bình của tam giác $SAD$ nên $MN\parallel AD\Rightarrow MN\parallel  BC$\\
			 $\Rightarrow MN\parallel\left(SBC\right)$.\quad(1).
			\itemch Tương tự ta có $O,N$ theo thứ tự  là trung điểm của $BD$, $SD$ nên $ON$ là đường trung bình của tam giác $SBD\Rightarrow ON\parallel SB\Rightarrow ON\parallel\left(SBC\right)$.\quad(2)\\
			Từ $(1)$ và  $(2)$ suy ra $\left(OMN\right)\parallel\left(SBC\right)$.
			\begin{center}
				\begin{tikzpicture}[scale=1,font=\footnotesize,line join=round,line cap=round,>=stealth]
						\path 
						(0,0) coordinate  (A)
						(-1.5,-1.5) coordinate (B) 
						(4,0) coordinate (D) 
						($(B)+(D)-(A)$)  coordinate (C)
						(1,3) coordinate  (S)
						($(S)!0.5!(A)$) coordinate  (M)
						($(S)!0.5!(D)$)  coordinate (N)
						($(C)!0.5!(A)$) coordinate (O) 
						($(B)!0.5!(A)$)  coordinate (E)
						 ($(O)!0.7!(N)$) coordinate (F);
						 \draw
						 (S)--(B)--(C)--(D)--(S)--(C);
						\draw[dashed] 
						(S)--(A)--(C) 
						(B)--(A)--(D)--cycle 
						(O)--(M)--(N)--(O)--(E)--(M) 
						(E)--(F); 
						\foreach \p/\g  in {A/40,B/-90,C/-90,D/0,S/90,N/45,F/0,O/-90,E/-60,M/40}
						\draw[fill=black] (\p) circle (1pt) node[shift=(\g:3mm)] {$\p$}; 
					\end{tikzpicture}
			\end{center}
			\itemch Ta  có $OE$ là đường trung bình của tam giác $ABD$ nên $OE\parallel AD\Rightarrow OE\parallel MN$.\\
			Do đó $E\in\left(OMN\right)$  mà $F\in ON$, $ON\subset\left(OMN\right)\Rightarrow F\in\left(OMN\right)$.\\
			Ta có $\heva{&EF\subset\left(OMN\right)\\&\left(OMN\right)\parallel\left(SBC\right)} 
			 \Rightarrow EF\parallel\left(SBC\right)$.
			\itemch Vì $G$ thuộc mặt phẳng $\left(ABCD\right)$ và cách đều $AB$, $CD$ nên $G$ thuộc đường trung bình của hình bình hành $ABCD$ (ứng  với hai cạnh $AB$, $CD$).
			\begin{center}
				\begin{tikzpicture}[scale=1,font=\footnotesize,line join=round,line cap=round,>=stealth]
						\path 
						(0,0) coordinate  (A)
						(-1.5,-1.5) coordinate  (B)
						(4,0) coordinate (D) 
						($(B)+(D)-(A)$) coordinate  (C)
						 (1,3) coordinate (S)
						($(S)!0.5!(A)$) coordinate  (M)
						($(S)!0.5!(D)$)  coordinate (N)
						($(C)!0.5!(A)$)  coordinate (O)
						($(B)!0.5!(A)$) coordinate (E) 
						($(O)!0.7!(N)$)  coordinate (F)
						($(B)!0.5!(C)$)  coordinate (I)
						 ($(A)!0.5!(D)$) coordinate (J)
						($(O)!0.5!(J)$)  coordinate (G);
						 \draw
						 (S)--(B)--(C)--(D)--(S)--(C);
						\draw[dashed] 
						(S)--(A)--(C) 
						 (B)--(A)--(D)--cycle
						 (O)--(M)--(N)--(O)--(E)--(M)
						(E)--(F) 
						 (I)--(J)
						 (N)--(G);
						\foreach \p/\g in {A/40,B/-90,C/-90,D/0,S/90,N/45,F/0,O/-90,E/-60,M/40,I/-90,G/0} 
						\draw[fill=black]  (\p) circle (1pt) node[shift=(\g:3mm)] {$\p$};
						\fill[black] (J) circle (1pt); 
					\end{tikzpicture}
			\end{center}
			Gọi $I$ là trung điểm $BC$ thì $I$, $O$, $G$  thẳng hàng.\\
			Ta có $OI$ là đường trung bình  của $\Delta ABC$ nên $OI\parallel AB\Rightarrow OI\parallel\left(SAB\right)$.\quad(3)\\
			Tương tự ta có  $ON\parallel SB\Rightarrow ON\parallel\left(SAB\right)$.\quad(4)\\
			Từ  $(3)$, $(4)$ suy ra $\left(OIN\right)\parallel\left(SAB\right)$ mà $NG\subset\left(OIN\right)$ nên $NG\parallel\left(SAB\right)$.
		\end{itemchoice}
	} 
\end{ex}

% nguồn: Ok_dungsai_lop11_Lan2
\begin{ex}%[1H4V4-6]%[ma:7bd26548]
	Cho hình chóp $ S.ABCD$ có đáy $ ABCD$ là hình bình hành. Gọi $ H$, $I$, $K$ lần lượt là trung điểm của $ SA$, $SB$, $SC$. Gọi $ M$ là giao điểm của $  AI$ và $ KD$, $N$ là giao điểm của $ DH$ và $ CI$. Khi đó
	\choiceTF 
		{\True $  HI\parallel\left(ABCD\right)$}
		{\True $\left(HIK\right)\parallel\left(ABCD\right)$} 
		 {$SM$ và $HI$ chéo nhau}
		{\True $\left(SMN\right)$  cắt $\left(HIK\right)$}
	 \loigiai{
		\begin{center}
			\begin{tikzpicture}[scale=1,font=\footnotesize,line join=round,line cap=round,>=stealth]
					\path 
						 (0,0) coordinate (A)
						(-1.5,-1.5)  coordinate (D)
						 (4,0) coordinate (B)
						($(B)+(D)-(A)$) coordinate  (C)
						(1,3) coordinate  (S)
						($(S)!0.5!(A)$)  coordinate (H)
						($(S)!0.5!(B)$) coordinate (I) 
						 ($(S)!0.5!(C)$) coordinate (K)
						(intersection  of A--I and K--D) coordinate (M)
						(intersection of C--I and D--H)  coordinate (N)
						(intersection of  S--B and D--M) coordinate (i)
						 (intersection of N--C and D--M) coordinate (j);
					 \draw
						 (S)--(D)--(C)--(S)--(M)--(N)--(S)
						 (D)--(M)
						(i)--(B) 
						(j)--(C)--(B); 
					\draw[dashed] 
						 (i)--(S)--(A)
						(B)--(A)--(D) 
						 (D)--(N)--(j)
						 (A)--(M)
						 (H)--(I)--(K)--cycle;
					 \foreach \p/\g in {A/170,D/190,C/-90,B/0,S/100,N/60,M/20,I/60,K/-10,H/-70}
					\draw[fill=black] (\p) circle (1pt) node[shift=(\g:3mm)] {$\p$}; 
				\end{tikzpicture}
		\end{center}
		\begin{itemchoice}
			\itemch Vì $ HI$ là đường trung bình của $\Delta SAB$ nên $ HI\parallel AB$ mà $  AB\subset\left(ABCD\right)$\\
			$\Rightarrow HI\parallel\left(ABCD\right)$.\quad  (1)
			\itemch Tương tự ta có $ KI\parallel  BC$, $BC\subset\left(ABCD\right)\Rightarrow KI\parallel\left(ABCD\right)$.\quad (2)\\
			Mặt  khác $ HI\subset\left(HKI\right)$, $KI\subset\left(HKI\right)$, $HI\cap KI=I$.\quad (3)\\
			Từ $(1)$,  $(2)$ và $(3)$ suy ra $\left(HIK\right)\parallel\left(ABCD\right)$.
			\itemch $\heva{&M\in AI,\,AI\subset\left(SAB\right)\\ 
			&M\in DK,\,DK\subset\left(SCD\right)} 
			\Rightarrow  M\in\left(SAB\right)\cap\left(SCD\right) \Rightarrow SM=\left(SAB\right)\cap\left(SCD\right)$.\\
			Khi đó $\heva{&\left(SAB\right)\cap\left(SCD\right)=SM\\ 
			&AB\subset\left(SAB\right),\,CD\subset\left(SCD\right)\\&AB\parallel CD}\Rightarrow SM\parallel AB\parallel CD\Rightarrow SM\parallel  HI$.\quad(4)\\
			 Vì $\heva{&N\in DH,\,DH\subset\left(SAD\right)\\
			&N\in  CI,\,CI\subset\left(SBC\right)}
			\Rightarrow N\in\left(SAD\right)\cap\left(SBC\right)  \Rightarrow SN=\left(SAD\right)\cap\left(SBC\right)$.\\
			Khi đó ta  có $\heva{&\left(SAD\right)\cap\left(SBC\right)=SN\\
			&AD\subset\left(SAD\right),\,BC\subset\left(SBC\right)\\&AD\parallel  BC}\Rightarrow SN\parallel AD\parallel BC\Rightarrow SN\parallel KI$.\quad(5)\\
			Mặt khác  ba điểm $ S$, $M$, $N$ không thẳng hàng.\quad $(6)$
			\itemch Từ $(4)$, $(5)$, $(6)$  suy ra $\left(SMN\right)\parallel\left(HIK\right)$.
		\end{itemchoice}
	 }
\end{ex}

% nguồn: Ok_dungsai_lop11_Lan2
\begin{ex}%[1H4V4-6]%[ma:3a6e23a3]
	Cho lăng trụ tam giác $  ABC.A'B'C'$ có $ I$, $K$, $G$ lần lượt là trọng tâm các tam giác $ ABC$, $A'B'C'$, $ACC'$. Gọi $ M$, $M'$ lần lượt là trung điểm của $ BC$, $B'C'$. Khi đó
	\choiceTF 
		{\True $ AMM'A'$ là hình bình hành} 
		 {$\dfrac{AI}{AM}=\dfrac{AG}{AN}=\dfrac{1}{3}$}
		{Mặt phẳng $\left(IKG\right)$ cắt $\left(BCC'B'\right)$} 
		{$\left(A'KG\right)\parallel\left(AIB'\right)$} 
	\loigiai{ 
		\begin{center}
			\begin{tikzpicture}[scale=1,font=\footnotesize,line join=round,line cap=round,>=stealth]
					 \path
						(0,0) coordinate (A') 
						(1,.2-1.5)  coordinate (B')
						 (5,0) coordinate (C')
						 (-0.7,-4) coordinate (A)
						($(A)+(B')-(A')$) coordinate (B) 
						($(B)+(C')-(B')$)  coordinate (C)
						($(B)!0.5!(C)$) coordinate  (M)
						($(C')!0.5!(C)$)  coordinate (N)
						 ($(B')!0.5!(C')$) coordinate (M')
						 ($(A)!2/3!(M)$) coordinate (I)
						 ($(A')!2/3!(M')$) coordinate (K)
						($(A)!2/3!(N)$) coordinate  (G);
					 \draw
						(A')--(B')--(C')--cycle 
						(A')--(A)--(B)--(B') 
						(B)--(C)--(C') 
						(A')--(M') 
						 (M)--(N);
					\draw[dashed] 
						(A)--(C) 
						(K)--(I)--(G)--cycle (M)--(A)--(N); 
					\foreach \p/\g  in {A'/120,B'/190,C'/10,A/180,B/-90,C/0,M/-45,I/-90,G/-60,N/0,M'/-60,K/75}
					\draw[fill=black]  (\p) circle (1pt) node[shift=(\g:3mm)] {$\p$};
				\end{tikzpicture}
		\end{center}
		\begin{itemchoice}
			\itemch $ MM'$ là đường trung bình của hình bình hành $ BCC'B'$ nên  $\heva{&MM'\parallel BB'\\&MM'=BB'}
			 \Rightarrow\heva{&MM'\parallel AA'\\&MM'=AA'.}$ Suy ra $AMM'A'$ là hình bình hành.
			\itemch Vì $  I$, $K$ theo thứ tự là trọng tâm các tam giác $ ABC$, $A'B'C'$ nên
			$ IM=KM'=\dfrac{1}{3}A'M'=\dfrac{1}{3}AM$, mà $ IM\parallel  KM'$ nên $ IKM'M$ là hình bình hành.\\
			 Suy ra $ IK\parallel MM'$, $MM'\subset\left(BCC'B'\right)\Rightarrow IK\parallel\left(BCC'B'\right)$.\quad(1)\\
			Gọi $ N$ là trung điểm của $ CC'$ tam  giác $ AMN$ có $\dfrac{AI}{AM}=\dfrac{AG}{AN}=\dfrac{2}{3}$.
			\itemch Vì $\dfrac{AI}{AM}=\dfrac{AG}{AN}=\dfrac{2}{3}$ nên $ IG\parallel MN$ mà $ MN\subset\left(BCC'B'\right)$ nên \break  $ IG\parallel\left(BCC'B'\right)$.\quad(2)\\
			Từ $(1)$ và $(2)$  suy ra $\left(IKG\right)\parallel\left(BCC'B'\right)$.
			\begin{center}
				\begin{tikzpicture}[scale=1,font=\footnotesize,line join=round,line cap=round,>=stealth]
						\path 
						 (0,0) coordinate (A')
						 (1,.2-1.5) coordinate (B')
						(5,0) coordinate  (C')
						(-0.7,-4) coordinate  (A)
						 ($(A)+(B')-(A')$) coordinate (B)
						 ($(B)+(C')-(B')$) coordinate (C)
						($(B)!0.5!(C)$) coordinate (M) 
						($(C')!0.5!(C)$)  coordinate (N)
						($(B')!0.5!(C')$)  coordinate (M')
						($(A)!2/3!(M)$) coordinate (I) 
						 ($(A')!2/3!(M')$) coordinate (K)
						($(A)!2/3!(N)$) coordinate (G) 
						($(A)!0.5!(C')$) coordinate  (O);
						\draw 
						 (A')--(B')--(C')--cycle
						(A')--(A)--(B)--(B') 
						 (A)--(B')--(M)--(N)
						 (B)--(C)--(C')
						(A')--(M'); 
						 \draw[dashed]
						 (A)--(C)
						 (K)--(I)--(G)--cycle (M)--(A)--(N)
						 (K)--(I)--(B')
						(A)--(C') 
						 (A')--(C);
						\foreach  \p/\g in {A'/120,B'/190,C'/10,A/180,B/-90,C/0,M/-45,I/-90,G/70,N/0,M'/-60,K/75,O/-90}
						\draw[fill=black] (\p)  circle (1pt) node[shift=(\g:3mm)] {$\p$};
					\end{tikzpicture}
			\end{center}
			\itemch  Vì $\left(A'KG\right)\equiv\left(A'M'C\right),\left(AIB'\right)\equiv\left(AMB'\right)$ ta cần chứng minh $\left(A'M'C\right)\parallel\left(AMB'\right)$.\\
			Dễ thấy $ AMM'A'$ là hình bình hành nên $ AM\parallel A'M'$ mà $ A'M'\subset\left(A'M'C\right)$ nên  $ AM\parallel\left(A'M'C\right)$.\quad (3)\\
			Ta có $\heva{&CM\parallel  B'M'\\&CM=B'M'}\Rightarrow CMB'M'$ là hình bình hành.\\
			Suy ra $ B'M\parallel CM'$, $CM'\subset\left(A'M'C\right)\Rightarrow B'M\parallel\left(A'M'C\right)$.\quad(4)\\ 
			Từ $(3)$ và $(4)$ suy ra $\left(A'M'C\right)\parallel\left(AMB'\right)$ hay $\left(A'KG\right)\parallel\left(AIB'\right)$. 
		\end{itemchoice}
	} 
\end{ex}

% nguồn: Ok_dungsai_lop11_Lan2
\begin{ex}%[1H4V4-6]%[ma:44939992]
	Cho hình bình hành $ABCD$ và $ABEF$ nằm ở hai mặt phẳng khác nhau. Gọi $M$ là trọng tâm $\Delta ABE$. Gọi $(P)$ là mặt phẳng đi qua $M$ và song song với mặt $\left(ADF\right)$. Gọi $N$ là giao  điểm của $(P)$ và $AC$. Khi đó
	\choiceTF 
		 {$EFDC$ là hình thang}
		{\True $FD\parallel EC$} 
		{\True  $\left(ADF\right)\parallel\left(BCE\right)$}
		{$\dfrac{AN}{NC}=3$} 
	 \loigiai{
		\begin{center}
			\begin{tikzpicture}[scale=1,font=\footnotesize,line join=round,line cap=round,>=stealth]
					\path 
						(0,0) coordinate (B) 
						(1,3) coordinate (E) 
						(4,0)  coordinate (C)
						(2,1) coordinate  (A)
						($(A)+(C)-(B)$) coordinate (D) 
						($(A)+(E)-(B)$)  coordinate (F);
					\draw 
						(B)--(C)--(D)--(F)--(E)--(C) 
						(B)--(E); 
					\draw[dashed] 
						 (B)--(A)--(D)
						(F)--(A); 
					 \foreach \p/\g in {A/120,B/180,C/-30,D/0,E/170,F/90}
					\draw[fill=black]  (\p) circle (1pt) node[shift=(\g:3mm)] {$\p$};
				\end{tikzpicture}
		\end{center}
		\begin{itemchoice}
			\itemch Cho hình  bình hành $ABCD$ và $ABEF$ nằm ở hai mặt phẳng khác nhau.
			\itemch  Ta có $\heva{&EF\parallel CD\left(\parallel AB\right)\\&EF=CD\left(=AB\right)}
			\Rightarrow EFDC $ là hình bình hành $\Rightarrow  FD\parallel EC$.
			\itemch Ta có $\heva{&AD\parallel BC;AF\parallel  BE\\&AD,\,AF\subset\left(ADF\right);\,AD\cap AF=A\\&
			BC,\,BE\subset\left(BEC\right);  \,BC\cap BE=B}
			 \Rightarrow\left(ADF\right)\parallel\left(BCE\right)$.
			 \itemch Tính $\dfrac{AN}{NC}$.
			\begin{center}
				\begin{tikzpicture}[scale=1,font=\footnotesize,line join=round,line cap=round,>=stealth]
						\path 
						(0,0)  coordinate (B)
						(1,3)  coordinate (E)
						(4,0)  coordinate (C)
						 (2,1) coordinate (A)
						 ($(A)+(C)-(B)$) coordinate (D)
						 ($(A)+(E)-(B)$) coordinate (F)
						($(C)!1/3!(D)$) coordinate (K) 
						 ($(E)+(K)-(C)$) coordinate (J)
						 ($(B)+(J)-(E)$) coordinate (I)
						($(I)!0.4!(J)$) coordinate (M) 
						(intersection of  E--M and A--B) coordinate (S)
						(intersection  of A--C and I--K) coordinate (N);
						 \draw
						(B)--(C)--(D)--(F)--(E)--(C) 
						 (B)--(E)
						(J)--(K); 
						\draw[dashed] 
						 (B)--(A)--(D)
						(F)--(A) 
						(J)--(I)--(K) 
						 (C)--(A)--(E)--(S);
						\foreach \p/\g  in {A/-90,B/180,C/-30,D/0,E/170,F/90,J/90,M/0,I/160,S/-10,N/-140,K/-30}
						\draw[fill=black] (\p)  circle (1pt) node[shift=(\g:3mm)] {$\p$};
					\end{tikzpicture}
			\end{center}
			Vẽ mặt phẳng $(P)$ chứa $M$ và $(P)\parallel\left(ADF\right)$ cắt $AB$, $AC$,  $CD$, $EF$ lần lượt tại $I$, $N$, $K$, $J$.\\
			Ta  có $\dfrac{AI}{BI}=\dfrac{AN}{NC}\left(IN\parallel BC\right)$ và $\dfrac{EJ}{IS}=\dfrac{ME}{MS}=2\left(IS\parallel JE\right)$.\\
			Mặt khác $BI=EJ$ (tứ  giác $BIJE$ là hình bình hành)\\
			$\Rightarrow\dfrac{BI}{IS}=2\Rightarrow\dfrac{BI}{2}=\dfrac{IS}{1}=\dfrac{BI+IS}{2+1}=\dfrac{BS}{3}\Rightarrow BI=\dfrac{2}{3}BS;IS=\dfrac{1}{3}BS$.\\ 
			 Ta có $AI=AS+AI=BS+\dfrac{1}{3}BS=\dfrac{4}{3}BS\Rightarrow >\dfrac{AI}{BI}=\dfrac{\dfrac{4}{3}BS}{\dfrac{2}{3}BS}=2\Rightarrow\dfrac{AN}{NC}=2$.
		\end{itemchoice}
	} 
\end{ex}

% nguồn: Ok_dungsai_lop11_Lan2
\begin{ex}%[1H4V4-6]%[ma:5f8b4114]
	Cho hình hộp $ABCD.A'B'C'D'$. Gọi  $G_1$, $G_2$ là trọng tâm của các tam giác $A'BD,B'D'C$. Khi đó
	\choiceTF 
		{\True $A'D'CB$ là hình bình  hành}
		{\True $\left(A'BD\right)\parallel\left(B'D'C\right)$} 
		{\True $G_1$, $G_2$ cùng thuộc  $AC'$}
		 {$G_1G_2=\dfrac{2}{3}AC'$}
	 \loigiai{
		\begin{center}
			\begin{tikzpicture}[scale=1,font=\footnotesize,line join=round,line cap=round,>=stealth]
					 \path
						 (0,0) coordinate (D')
						(1.5,1.5) coordinate (A') 
						(5,0) coordinate (C') 
						 ($(A')+(C')-(D')$) coordinate (B')
						(0,3)  coordinate (D)
						($(A')+(D)-(D')$)  coordinate (A)
						 ($(B')+(A)-(A')$) coordinate (B)
						($(B)+(D)-(A)$) coordinate  (C);
					\draw 
						 (A)--(D)--(C)--(B)--cycle
						 (D)--(D')--(C')--(C)
						(C')--(B')--(B)--(D) 
						(D')--(C)--(B'); 
					 \draw[dashed]
						 (D')--(A')--(B')--cycle
						 (D)--(A')--(B)
						(A)--(A'); 
					\foreach \p/\g in  {A'/180,B'/0,C'/-30,D'/190,A/90,B/65,C/0,D/180}
					 \draw[fill=black] (\p) circle (1pt) node[shift=(\g:3mm)] {$\p$};
				\end{tikzpicture}
		\end{center}
		\begin{itemchoice}
			\itemch Vì $ABCD.A'B'C'D'$ là hình hộp nên $\heva{&A'D'\parallel BC\\&A'D'=BC} \Rightarrow A'D'CB$ là hình bình  hành.\\
			Suy ra  $A'B\parallel CD'\Rightarrow A'B\parallel\left(B'D'C\right)$.\quad(1)\\
			Tương tự ta  có $\heva{&A'B'\parallel CD\\&A'B'=CD}
			\Rightarrow A'B'CD $  là hình bình hành.\\
			Suy ra  $A'D\parallel B'C\Rightarrow A'D\parallel\left(B'D'C\right)$.\quad(2)
			 \itemch Từ $(1)$ và $(2)$ suy ra $\left(A'BD\right)\parallel\left(B'D'C\right)$.
			\itemch Gọi  $O$, $O'$, $I$ theo thứ tự là tâm của các hình bình hành $ABCD$, $A'B'C'D'$, $ACC'A'$.\\
			Vì $G_1$ là trọng tâm tam giác $AB'D$ nên $\dfrac{A'{G_1}}{A'O}=\dfrac{2}{3}$ $\Rightarrow{G_1}$ là trọng tâm tam giác $A'AC$ nên suy  ra $G_1=AI\cap A'O$.\quad (3)\\
			Tương tự, $G_2$  là trọng tâm tam giác $B'D'C$ nên $\dfrac{C{G_2}}{C{O^\prime}}=\dfrac{2}{3}$\\
			$\Rightarrow{G_2}$ là trọng  tâm tam giác $A'C'C$ suy ra $G_2=C'I\cap CO'$.\quad(4)\\
			Từ $(3)$ và $(4)$ suy ra  $G_1,G_2$ cùng thuộc $AC'$.
			\begin{center}
				\begin{tikzpicture}[scale=1,font=\footnotesize,line join=round,line cap=round,>=stealth]
						 \path
						(0,0) coordinate  (D')
						(1.5,1.5)  coordinate (A')
						 (5,0) coordinate (C')
						($(A')+(C')-(D')$) coordinate  (B')
						(0,4) coordinate  (D)
						($(A')+(D)-(D')$)  coordinate (A)
						($(B')+(A)-(A')$)  coordinate (B)
						($(B)+(D)-(A)$) coordinate (C) 
						($(B)!0.5!(D)$) coordinate (O) 
						($(B')!0.5!(D')$)  coordinate (O')
						($(A)!0.5!(C')$) coordinate (I) 
						 (intersection of A--C' and A'--O) coordinate (G_1)
						(intersection  of A--C' and C--O') coordinate (G_2);
						\draw 
						(A)--(D)--(C)--(B)--cycle 
						(D)--(D')--(C')--(C)--(A) 
						 (C')--(B')--(B)--(D)
						 (D')--(C)--(B');
						 \draw[dashed]
						(D')--(A')--(B')--cycle 
						(D)--(A')--(B) 
						(A')--(A)--(C')--cycle 
						 (O)--(A')--(C)--(O');
						\foreach \p/\g in  {A'/180,B'/0,C'/-30,D'/190,A/90,B/65,C/0,D/180,O/90,O'/-90,G_1/180,G_2/0,I/90}
						\draw[fill=black] (\p) circle (1pt)  node[shift=(\g:3mm)] {$\p$};
					\end{tikzpicture}
			\end{center}
			Ta có  $\dfrac{AG_1}{AI}=\dfrac{2}{3}\Rightarrow\dfrac{AG_1}{A{C^\prime}}=\dfrac{1}{3}$; $\dfrac{C'G_2}{C^\prime I}=\dfrac{2}{3}\Rightarrow\dfrac{C'G_2}{A{C^\prime}}=\dfrac{1}{3}$.
			\itemch Do vậy $A{G_1}=G_1G_2=G_2C'=\dfrac{1}{3}AC'$.\\ 
			Vậy $G_1$,  $G_2$ cùng thuộc $AC'$, đồng thời chia $AC'$ thành ba phần bằng nhau.
		\end{itemchoice}
	} 
\end{ex}

% nguồn: Ok_dungsai_lop11_Lan2
\begin{ex}%[1H4V4-6]%[ma:b4c8e734]
	Trong mặt phẳng $(P)$ cho hình bình hành $ ABCD$. Vẽ các nửa đường thẳng song song nhau, nằm về một phía đối với mặt phẳng $(P)$ và đi qua các điểm $A$, $B$, $C$, $D$. Một mặt phẳng $(Q)$ cắt bốn nửa đường thẳng  nói trên tại $ A'$, $B'$, $C'$, $D'$. Khi đó
	 \choiceTF
		{\True Mặt phẳng $\left(AA',BB'\right)$ song song với mặt phẳng $\left(CC',DD'\right)$} 
		{\True $  A'B'\parallel C'D'$}
		{Tứ giác $  A'B'C'D'$ là hình thang}
		{\True Gọi $ O$ và $ O'$ lần lượt là giao điểm của hai đường chéo của $ ABCD$ và $  A'B'C'D'$. Khi đó $ OO'\parallel AA'$}
	 \loigiai{
		\begin{center}
			\begin{tikzpicture}[scale=1,font=\footnotesize,line join=round,line cap=round,>=stealth]
					 \path
						(0,0) coordinate (A) 
						 (-1.5,-1.5) coordinate (B)
						(4,0) coordinate  (D)
						 ($(B)+(D)-(A)$) coordinate (C)
						 (1,4) coordinate (a)
						($(B)+(a)-(A)$) coordinate (b) 
						 ($(b)+(C)-(B)$) coordinate (c)
						($(a)+(c)-(b)$) coordinate (d) 
						($(A)!0.9!(a)$)  coordinate (A')
						($(B)+(A')-(A)$) coordinate  (B')
						($(D)!0.5!(d)$)  coordinate (D')
						($(B')+(D')-(A')$) coordinate  (C')
						 ($(B)!0.5!(D)$) coordinate (O)
						($(B')!0.5!(D')$) coordinate (O'); 
					 \draw (B')--(C')--(D')--(A')--(B')--(D')
						 (A')--(C')
						 (B)--(C)--(D)
						(B)--(b) 
						 (C)--(c)
						(D)--(d) 
						 (A')--(a);
					\draw[dashed] 
						(A')--(A)--(C) 
						(B)--(A)--(D)--cycle 
						(O)--(O'); 
					 \foreach \p/\g in {A/170,B/190,C/-60,D/0,A'/160,B'/180,C'/-20,D'/0,O'/80,O/-90}
					\draw[fill=black]  (\p) circle (1pt) node[shift=(\g:3mm)] {$\p$};
				\end{tikzpicture}
		\end{center}
		\begin{itemchoice}
			\itemch Ta có $ AA'\parallel  DD'$ và $ AB\parallel CD$ nên $\text{mp}\left(AA',BB'\right)\parallel \text{mp}\left(CC',DD'\right)$.
			\itemch Chứng minh  $ A'B'C'D'$ là hình bình hành:\\
			Ta có  $\heva{&\text{mp}\left(AA',BB'\right)\parallel \text{mp}\left(CC',DD'\right)\\
			 &(Q)\cap \text{mp}\left(A{A^\prime},\,B{B^\prime}\right)=A'B'\\
			&(Q)\cap \text{mp}\left(CC',DD'\right)=C'D'} 
			\Rightarrow A'B'\parallel C'D'.\quad(1)$\\ 
			Hoàn toàn tương tự  ta chứng minh được $ A'D'\parallel B'C'.\quad(2)$
			\itemch Từ $(1)$  và $(2)$ suy ra $ A'B'C'D'$ là hình bình hành.
			\itemch Ta có  $\heva{&\left(ACC'A'\right)\cap\left(BDD'B'\right)=OO'\\
			 &AA'\subset\left(ACC'A'\right),BB'\subset\left(BDD'B'\right)\\
			&AA^\prime\parallel BB^\prime} 
			\Rightarrow OO'\parallel AA'\parallel BB'$ hay $ OO'\parallel  AA'$.
		\end{itemchoice}
	 }
\end{ex}


%%%%%%%%%%%%%%%% TRẢ LỜI NGẮN %%%%%%%%%%%%%%%%

% nguồn: Ok_TraLoiNgan_Lop11_Lan1
\begin{ex}%[1H4H4-6]%[ma:0b41164e]
	Cho hình chóp  $S.ABC$ có $G_1$, $G_2$, $G_3$ lần lượt là trọng tâm các tam giác $SAB$, $SBC$, $SAC$. Số giao tuyến của mặt phẳng $(G_1G_2G_3)$ với các mặt của hình chóp $S.ABC$ là
	\shortans{3} 
	 \loigiai{
		\begin{center}
			\begin{tikzpicture}[scale=1,font=\footnotesize,line join=round,line cap=round,>=stealth]
					\def\ac{4} 
					\def\ab{2} 
					 \def\as{4}
					 \def\gocA{50}
					\coordinate[label=left:$A$] (A)  at (0,0);
					\coordinate[label=right:$C$] (C) at  (\ac,0);
					\coordinate[label=below  left:$B$] (B) at (-\gocA:\ab);
					 \coordinate[label=above:$S$] (S) at (70:\as);
					\coordinate[label=left:$M$]  (M) at ($(A)!0.5!(B)$);
					\coordinate[label=below  right:$N$] (N) at ($(B)!0.5!(C)$);
					\coordinate[label=below:$P$] (P)  at ($(A)!0.5!(C)$);
					\coordinate[label=left:$G_1$] (G1)  at ($(S)!2/3!(M)$);
					\coordinate[label=right:$G_2$] (G2) at ($(S)!2/3!(N)$); 
					\coordinate[label=left:$G_3$] (G3) at ($(S)!2/3!(P)$); 
					\draw (A)--(B)--(C)--(S)--cycle  (S)--(B)
						(M)--(S)--(N); 
					 \draw[dashed] (A)--(C)
						(S)--(P) 
						 (M)--(N)--(P)--cycle (G1)--(G2)--(G3)--cycle;
					\foreach \diem in  {A,B,C,S,M,N,P,G1,G2,G3}\fill (\diem) circle (1pt);
				\end{tikzpicture}
		\end{center}
		\begin{itemize}
			\item Gọi $M$, $N$, $P$  theo thứ tự là trung điểm các cạnh $AB$, $BC$, $AC$.\\
			Tam  giác $SMN$ có $\dfrac{SG_1}{SM}=\dfrac{SG_2}{SN}=\dfrac{2}{3}$ (tính chất trọng tâm).\\
			Suy ra $G_1G_2 \parallel MN$ mà $MN  \subset(ABC) \Rightarrow G_1G_2 \parallel (ABC)$. (1)\\
			Tương tự, tam giác  $SMP$ có $\dfrac{SG_1}{S M}=\dfrac{SG_3}{S P}=\dfrac{2}{3}$ (tính chất trọng tâm).\\
			Suy ra  $G_1G_3 \parallel MP$ mà $M P \subset(ABC) \Rightarrow G_1G_3 \parallel (ABC)$. (2)\\
			Từ (1) và  $(2)$ suy ra $(G_1G_2G_3) \parallel (ABC)$.
			\item Vì $\heva{&G_1\in(G_1G_2G_3)\cap (SAB)\\&G_2\in(G_1G_2G_3)\cap (SBC)\\&G_3\in(G_1G_2G_3)\cap (SAC)}$ nên mặt phẳng $(G_1G_2G_3)$ có giao tuyến với các mặt phẳng $(SAB)$, $(SAC)$ và $(SBC)$. 
		\end{itemize}
	} 
\end{ex}

% nguồn: Ok_TraLoiNgan_Lop11_Lan1
\begin{ex}%[1H4H4-6]%[ma:b2cc265d]
	Cho hình chóp $S.ABCD$ có đáy là hình bình hành. $M$ là trung điểm của $SC$. Gọi $I$ là giao điểm của đường thẳng  $AM$ với mặt phẳng $(SBD)$. Tính tỉ số $\dfrac{IA}{IM}$.
	 \shortans{2}
	 \loigiai{
	\immini 
		{Gọi $AC\cap BD=O$ thì $(SAC)\cap(SBD)=SO$.\\ 
		Trong mặt phẳng $(SAC)$, lấy $AM\cap SO=I$\\ $\Rightarrow I=AM\cap(SBD)$.\\ 
		Do trong $\triangle SAC$, $AM$ và $SO$ là hai đường trung tuyến, nên $I$ là trọng tâm  $\triangle SAC$.\\
		 Vậy $IA=2IM$.
		} 
		{\begin{tikzpicture}[scale=0.7,font=\footnotesize,line join=round,line cap=round,>=stealth]
		\path 
			 (0,0) coordinate (A)
			(-1.5,-1.5) coordinate  (D)
			(4,0) coordinate  (B)
			($(B)-(A)+(D)$)  coordinate (C)
			($(A)+(0,4)$) coordinate (S) 
			 ($(S)!1/2!(C)$) coordinate (M);
		\coordinate (O)  at (intersection cs:first line={(C)--(A)},second line={(B)--(D)});
		\coordinate (I) at (intersection cs:first line={(S)--(O)},second  line={(A)--(M)});
		 \draw (S)--(D)--(C)--(B)--cycle (S)--(C);
		\draw[dashed] (A)--(D) 
			 (S)--(A)--(B)
			(A)--(C) 
			 (B)--(D)
			 (S)--(O)
			 (A)--(M);
		\foreach  \x/\g in {A/180,B/0,C/0,D/180,S/90,O/-90,I/180,M/0}\fill (\x) circle (1pt)+(\g:4mm) node[black]{$ \x $};
	\end{tikzpicture}
		} 
	 }
\end{ex}

% nguồn: Ok_TraLoiNgan_Lop11_Lan1
\begin{ex}%[1H4H4-6]%[ma:1eebe7c9]
	Cho hình chóp $S.ABCD$ có đáy $ABCD$ là hình bình hành tâm $O$. Gọi $M$, $N$ lần lượt là trung điểm của $SA$ và $SD$. Tính tỉ số diện tích  $\dfrac{S_{\triangle SBC}}{S_{\triangle ONN}}$.
	 \shortans{4}
	 \loigiai{
	\immini{Ta có $OM$, $ON$ lần lượt là  đường trung bình của tam giác $SAC$ và $SBD$ nên $SC\parallel OM$, $SC=2OM$ và $SB\parallel ON$, $SB=2ON$.\\
		Do đó, $$\dfrac{S_{\triangle SBC}}{S_{\triangle OMN}}=\dfrac{\tfrac{1}{2}\cdot SB\cdot SC\cdot\sin\widehat{BSC}}{\tfrac{1}{2}\cdot ON\cdot OM\cdot\widehat{NOM}}=\dfrac{SB}{ON}\cdot\dfrac{SC}{OM}=4.$$}{\begin{tikzpicture}[scale=1,font=\footnotesize,line join=round,line cap=round,>=stealth]
		\path (0,0) coordinate  (B)++(3,0) coordinate (C)++(1.5,1) coordinate (D)
			($(B)+(D)-(C)$) coordinate (A) 
			(0.5,2.5) coordinate (h) 
			($(A)+(h)$)  coordinate (S)
			 (barycentric cs:S=1,A=1) coordinate (M)
			 (barycentric cs:S=1,D=1) coordinate (N)
			(barycentric cs:C=1,A=1) coordinate  (O);
		\draw  (S)--(B)--(C)--(D)--cycle (S)--(C);
		\draw[dashed] (S)--(A)--(B)--(D)--(A)--(C) 
			(O)--(M)--(N)--cycle; 
		\foreach \p/\g in {A/150,B/240,C/-60,D/0,S/90,M/180,N/30,O/-90} \fill[black] (\p) circle (1pt)+(\g:0.3)  node{$\p$};
	\end{tikzpicture}
		 }
	 }
\end{ex}

% nguồn: Ok_TraLoiNgan_Lop11_Lan1
\begin{ex}%[1H4V4-6]%[ma:e416abfb]
	Cho hình chóp $S.ABC$ có $D$, $E$, $F$ lần lượt là trọng tâm các tam giác $SAB$, $SBC$, $SCA$. Tam giác  $ABC$ có diện tích bằng $2025$ cm$^2$. Mặt phẳng $(DEF)$ cắt các cạnh $SA$, $SB$, $SC$ lần lượt tại $A'$, $B'$, $C'$. Tính diện tích tam giác $A'B'C'$.
	\shortans{900} 
	 \loigiai{
	\immini{Gọi  $M$, $N$, $P$ lần lượt là trung điểm của $AB$, $BC$, $CA$.\\
		Ta có $\dfrac{SD}{SM}=\dfrac{1}{3}=\dfrac{SE}{SN}\Rightarrow  DE\parallel MN$. Mà $MN\subset (ABC)$ nên $DE\parallel (ABC)$.\\
		Tương  tự, $DF\parallel (ABC)$.\\
		Suy ra $(DEF)\parallel  (ABC)$.\\
		Ta  có $(DEF)\equiv (A'B'C')$ nên $(A'B'C')\parallel (ABC)$.\\
		Dễ thấy  $A'B'\parallel AB$ nên $\dfrac{A'B'}{AB}=\dfrac{SD}{SM}=\dfrac{2}{3}$.\\
		Tương tự $\dfrac{A'C'}{AC}=\dfrac{2}{3}$.}{\begin{tikzpicture}[scale=1,font=\footnotesize,line join=round,line cap=round,>=stealth]
		 \path (0,0) coordinate (A)
			(4,0) coordinate  (C)
			(1,-1) coordinate  (B)
			($(A)+(1.5,2.5)$)  coordinate (S)
			 (barycentric cs:S=1,A=2) coordinate (A')
			(barycentric cs:S=1,B=2)  coordinate (B')
			 (barycentric cs:S=1,C=2) coordinate (C')
			 (barycentric cs:S=1,A=1,B=1) coordinate (D)
			(barycentric cs:S=1,B=1,C=1) coordinate (E) 
			(barycentric cs:S=1,A=1,C=1) coordinate (F) 
			 (barycentric cs:A=1,B=1) coordinate (M)
			(barycentric cs:C=1,B=1)  coordinate (N)
			(barycentric  cs:A=1,C=1) coordinate (P);
		 \draw (S)--(A)--(B)--(C)--cycle (M)--(S)--(B)
			 (N)--(S)
			(A')--(B')--(C'); 
		\draw[dashed]  (A)--(C)
			(S)--(P) 
			(A')--(C'); 
		\foreach \p/\g in {A/180,B/-90,C/0,S/90,A'/150,B'/45,C'/30,D/210,E/-30,F/135,M/-120,N/-60,P/-90} \fill[black] (\p) circle (1pt)+(\g:0.3)  node{$\p$};
	\end{tikzpicture}
		 }\noindent
		Ta có $\dfrac{S_{\triangle A'B'C'}}{S_{\triangle ABC}}=\dfrac{\tfrac{1}{2}\cdot A'B'\cdot A'C'\cdot\sin\widehat{B'A'C'}}{\tfrac{1}{2}\cdot AB\cdot  AC\cdot\sin\widehat{BAC}}=\dfrac{A'B'}{AB}\cdot\dfrac{A'C'}{AC}=\dfrac{4}{9}$.\\
		Vậy $S_{\triangle A'B'C'}=\dfrac{4}{9}\cdot S_{\triangle ABC}=\dfrac{4}{9}\cdot  2025=900$ cm$^2$.
	 }
\end{ex}

% nguồn: Ok_TraLoiNgan_Lop11_Lan1
\begin{ex}%[1H4V4-6]%[ma:70ff9621]
	Cho hình chóp $S.ABCD$ có đáy $ABCD$ là hình thang với $AB\parallel CD$, $AB=2CD$, tam giác $SAB$ cân tại $S$, điểm $M$ thuộc cạnh $AD$ sao cho  $MD=2MA$. Mặt phẳng $(\alpha)$ qua $M$ và song song với mặt phẳng $(SAB)$; mặt phẳng $(\alpha)$ cắt các cạnh $BC$, $SC$, $SD$ lần lượt tại $N$, $P$, $Q$. Tính tỉ số $\dfrac{PQ}{MN}$.
	 \shortans{0,2}
	 \loigiai{
		\begin{center}
			\begin{tikzpicture}[scale=1,font=\footnotesize,line join=round,line cap=round,>=stealth]
					 \coordinate (A) at (0,0);
					\coordinate (B) at (4,0); 
					\coordinate (D)  at (-1,-1.7);
					\coordinate (C)  at (1,-1.7);
					 \coordinate (S) at (0.3,4);
					\coordinate (M) at ($(A)!1/3!(D)$); 
					\coordinate (Q) at  ($(S)!1/3!(D)$);
					\coordinate (P) at  ($(S)!1/3!(C)$);
					\coordinate (N) at  ($(B)!1/3!(C)$);
					\coordinate (I) at (intersection of  A--D and B--C);
					\draw ($(A)!1/3!(B)+(90:0.3)$) node[rotate=0]{$6x$}; 
					\draw ($(C)!1/2!(D)+(90:0.3)$) node[rotate=0]{$3x$}; 
					\draw  (S)--(D)--(C)--(B)--(S)--(C)
						(Q)--(P)--(N) 
						 (D)--(I)--(C);
					\draw[dashed] (D)--(A)--(S) 
						 (A)--(B)
						 (Q)--(M)--(N);
					\foreach \x/\g in{S/90,A/45,B/0,C/-45,D/180,M/-45,N/-45,P/45,Q/180,I/-135}\fill (\x)  circle (1pt)+(\g:.35) node{$\x$};
				\end{tikzpicture}
		\end{center}
		Đặt $CD=3x\Rightarrow AB=6x$.\\ 
		Theo giả thiết, ta có $(\alpha)$ đi qua $M$ song song với mặt phẳng  $(SAB)$ cắt các cạnh $BC,SC,SD$ lần lượt tại $N,P,Q \Rightarrow (MNPQ)\equiv (\alpha)$.\\
		Vì $\heva{&(\alpha)\parallel (SAB)\\&(\alpha)\cap (SAD)=MQ\\&(SAB)\cap (SAD)=SA}\Rightarrow MQ\parallel SA$.\\ 
		Tương tự, $MN\parallel AB$; $PQ\parallel  CD \parallel AB$.\\
		Vì $MQ\parallel SA$ nên $\dfrac{MQ}{SA}=\dfrac{DM}{DA}=\dfrac{2}{3}\Rightarrow \dfrac{PQ}{CD}=\dfrac{SQ}{SD}=\dfrac{1}{3}\Rightarrow  PQ=\dfrac{1}{3}CD=x$.\\
		Gọi  $AD\cap BC=I$.\\
		Vì $\heva{&CD\parallel AB\\&CD=\dfrac{1}{2}AB} \Rightarrow CD$ là đường trung bình của $\Delta IAB\Rightarrow D$ là  trung điểm cạnh $IA$.\\
		Theo giả thiết, ta có $MA=\dfrac{1}{3}AD=\dfrac{1}{6}AI\Rightarrow MI=\dfrac{5}{6}IA\Rightarrow MN=\dfrac{5}{6}AB=5x$.\\ 
	Suy ra $\dfrac{PQ}{MN}=\dfrac{x}{5x}=\dfrac{1}{5}=0{,}2$.} 
\end{ex}

% nguồn: Ok_TraLoiNgan_Lop11_Lan1
\begin{ex}%[1H4V4-6]%[ma:347aee86]
	Cho hình lập phương $ABCD.A'B'C'D'$. Các điểm $M$, $N$, $P$ theo thứ tự  là trung điểm các cạnh $A'B'$, $BC$, $DD'$. Số giao tuyến của mặt phẳng $(MNP)$ với các mặt phẳng $(CB'D')$, $(A'BD)$, $(ABCD)$ và $(A'B'C'D')$ là
	 \shortans{2}
	 \loigiai{
		\begin{itemize}
			\item Gọi $O, I$ theo thứ tự là trung điểm của $B'C$  và $B'D$.\\
			Vì $ON$ là đường trung bình của tam giác $  CB'$ nên $\heva{&ON\parallel BB'\\ & ON=\dfrac 12BB'}\Rightarrow\heva{&ON\parallel PD'\\ & ON=PD'}$.\\
			Suy ra $PNOD'$ là hình bình  hành.\\
			Vì vậy $PN\parallel  OD'$ mà $OD'\subset(B'D'C)$ nên $PN\parallel(B'D'C)$.(1)
			\begin{center}
				\begin{tikzpicture}[scale=1,font=\footnotesize,line join=round,line cap=round,>=stealth]
						\def\bc{4} 
						 \def\ba{2}
						\def\gocB{35} 
						 \coordinate[label=below left:$B$] (B) at (0,0);
						\coordinate[label=above left:$A$]  (A) at (\gocB:\ba);
						\coordinate[label=below:$C$] (C) at (\bc,0); 
						\coordinate[label=right:$D$] (D) at ($(C)-(B)+(A)$); 
						\coordinate[label=above left:$A'$]  (A') at ($(A)+(90:\bc)$);
						\coordinate[label=left:$B'$]  (B') at ($(B)-(A)+(A')$);
						 \coordinate[label=above:$C'$] (C') at ($(C)-(A)+(A')$);
						\coordinate[label=right:$D'$] (D') at  ($(D)-(A)+(A')$);
						\coordinate[label=above:$M$] (M)  at ($(A')!0.5!(B')$);
						\coordinate[label=below:$N$] (N) at  ($(B)!0.5!(C)$);
						\coordinate[label=right:$P$] (P)  at ($(D)!0.5!(D')$);
						\coordinate[label=below:$O$]  (O) at ($(C)!0.5!(B')$);
						\coordinate[label=above:$I$]  (I) at ($(D')!0.5!(B')$);
						 \draw (B')--(B)--(C)--(D)--(D')--(A')--(B')--(C')--(D')
						(B')--(C)--(C') 
						(B')--(D'); 
						\draw[dashed] (A')--(A)--(D) 
						(A)--(B) 
						(M)--(N)--(P)--cycle (D')--(O) 
						 (C)--(I);
						\foreach \diem in  {A,B,C,D,A',B',C',D',M,N,P,O,I} \fill (\diem) circle (1pt);
					\end{tikzpicture}
			\end{center}
			 Tương tự, $IM$ là đường trung bình của tam giác $A'B'D'$ nên $\heva{&I M\parallel A'D'\\ & I M=\dfrac 12A'D'}$\\
			$\Rightarrow\heva{&I M  \parallel CN \\& IM=CN}$. Do vậy $IMNC$ là hình bình hành.\\
			 Suy ra $\heva{& MN\parallel IC\\ & IC\subset(B'D'C)}
			 \Rightarrow MN\parallel(B'D'C)$.(2)\\
			Từ (1) và (2) suy  ra $(MNP)\parallel(B'D'C)$.
			\item Tương tự $(MNP)\parallel(A'BD)$. 
			\item Vì $\heva{&M\in(MNP)\cap (A'B'C'D')\\&N\in(MNP)\cap (ABCD)\\&B\in(BMN)\cap (ABCD)}$ nên mặt phẳng $(MNP)$ có giao tuyến với các mặt phẳng $(ABCD)$  và $(A'B'C'D')$.
		\end{itemize}
	} 
\end{ex}

% nguồn: Ok_TraLoiNgan_Lop11_Lan1
\begin{ex}%[1H4V4-6]%[ma:ffd3b3ec]
	Cho hình chóp $S.ABCD$ có đáy là hình thang với $AD$ là đáy lớn và $AD=2 BC$. Gọi ${M, N}$ theo  thứ tự là trung điểm của $SA$ và $AD$. Số giao tuyến hai mặt phẳng $(BMN)$ với các mặt phẳng $(SCD)$, $(SAB)$, $(SAD)$, $(SBC)$ và $(ABCD)$ là
	 \shortans{4}
	 \loigiai{
		\begin{center}
			\begin{tikzpicture}[scale=1,font=\footnotesize,line join=round,line cap=round,>=stealth]
					 \def\ad{4}
					\def\ab{2} 
					\def\bc{2} 
					 \def\as{4}
					 \def\gocA{50}
					 \def\gocB{130}
					\coordinate[label=left:$A$]  (A) at (0,0);
					\coordinate[label=below  left:$B$] (B) at (-\gocA:\ab);
					\coordinate[label=below right:$C$] (C) at  ($(B)+(180-\gocA-\gocB:\bc)$);
					\coordinate[label=right:$D$]  (D) at (\ad,0);
					\coordinate[label=above:$S$] (S) at  (75:\as);
					\coordinate[label=left:$M$]  (M) at ($(S)!0.5!(A)$);
					 \coordinate[label=above right:$N$] (N) at ($(A)!0.5!(D)$);
					\draw (A)--(B)--(C)--(D)--(S)--cycle  (B)--(S)--(C)
						 (M)--(B);
					\draw[dashed] (A)--(D) 
						(M)--(N)--(B); 
					\foreach \diem in {A,B,C,D,S,M,N} \fill (\diem)  circle (1pt);
				\end{tikzpicture}
		\end{center}
		\begin{itemize}
			\item Ta có $DN \parallel BC$ và $DN=BC=\dfrac{1}{2} AD$ nên $BCDN$ là  hình bình hành. \\
			Suy ra $BN \parallel CD  \Rightarrow BN \parallel (SCD)$. (1)\\
			Vì $MN$ là đường trung bình của tam giác $SAD$ nên  $MN \parallel SD \Rightarrow MN \parallel (SCD)$. (2) \\
			Mặt khác $BN \subset(BMN), MN \subset(BMN), BN \cap MN=N$.(3)\\ 
			Từ (1), (2), (3) suy ra  $(BMN) \parallel (SCD)$.
			\item Vì  $\heva{&M\in(BMN)\cap (SAB)\\&M\in(BMN)\cap (SAD)\\&B\in(BMN)\cap (ABCD)}$ nên mặt phẳng $(BMN)$ có giao tuyến với các mặt phẳng $(SAB)$, $(SAD)$, $(SBC)$ và $(ABCD)$.
		\end{itemize}
	 }
\end{ex}

% nguồn: Ok_TraLoiNgan_Lop11_Lan1
\begin{ex}%[1H4V4-6]%[ma:c28a7d09]
	Cho hình chóp $S.ABCD$, đáy $ABCD$ là hình bình hành có $O$ là giao điểm của hai đường chéo. Gọi $M$, $N$ lần lượt là trung điểm  của $SA$, $SD$. Gọi $E$ là trung điểm của $AB$ và $F$ là một điểm thuộc $ON$. Đường thẳng $EF$ cắt bao nhiêu mặt của hình chóp $S.ABCD$.
	 \shortans{4}
	 \loigiai{
		\begin{center}
			\begin{tikzpicture}[scale=1,font=\footnotesize,line join=round,line cap=round,>=stealth]
					 \def\bc{4}
					\def\ba{2} 
					 \def\h{5}
					 \def\gocB{40}
					\coordinate[label=below left:$B$] (B) at  (0,0);
					\coordinate[label=above  right:$A$] (A) at (\gocB:\ba);
					\coordinate[label=below:$C$] (C) at (\bc,0); 
					\coordinate[label=right:$D$]  (D) at ($(C)-(B)+(A)$);
					\coordinate[label=below:$O$] (O) at  ($(A)!.5!(C)$);
					 \coordinate[label=above:$S$] (S) at ($(O)+(90:\h)$);
					\coordinate[label=left:$M$] (M) at  ($(S)!0.5!(A)$);
					 \coordinate[label=right:$N$] (N) at ($(S)!0.5!(D)$);
					\coordinate[label=left:$E$] (E) at  ($(A)!0.5!(B)$);
					\coordinate[label=right:$F$] (F) at  ($(O)!0.65!(N)$);
					 \draw (B)--(C)--(D)--(S)--cycle (S)--(C);
					 \draw[dashed] (C)--(A)--(D)--(B)
						(O)--(S)--(A)--(B) 
						(M)--(N)--(O)--(E)--cycle  (E)--(F);
					\foreach \diem  in {A,B,C,D,S,O,M,N,E,F} \fill (\diem) circle (1pt);
				\end{tikzpicture}
		\end{center}
		\begin{itemize}
			\item  Do $OM$ là đường trung bình của tam giác $SAC$ nên $OM \parallel SC$ (1) .\\
			Lại có,  $ON$ là đường trung bình của tam giác $SDB$ nên $ON \parallel SB$ (2).\\
			Từ  (1) và (2) suy ra $(OMN) \parallel (SBC)$.\\
			Ta có $OE$ là đường trung bình của tam giác $ABD$ nên $OE \parallel AD$, mà $AD \parallel MN$, suy ra  $OE \parallel MN$. Do đó $E$ sẽ thuộc mặt phẳng $(OMN)$. Suy ra $EF \subset(OMN)$. Ta có $(OMN) \parallel (SBC)$ nên $EF \parallel (SBC)$.
			\item Vì $\heva{&E\in(EF)\cap (ABCD)\\&E\in(EF)\cap (SAB)\\&EF\nparallel MN\Rightarrow EF \text{ cắt } (SAD)\\&EF\parallel(SBC) \text{ mà } (SCD)\text{ cắt  } (SBC) \text{ nên } EF \text{ cắt } (SCD)} $ do đó mặt phẳng $EF$ có giao điểm với các mặt phẳng $(SAB)$, $(SAD)$, $(ABCD)$ và $(SCD)$.
		\end{itemize}
	} 
\end{ex}

% nguồn: Ok_TraLoiNgan_Lop11_Lan1
\begin{ex}%[1H4V4-6]%[ma:c1924d05]
	Cho hình chóp $S.ABCD$ có đáy $ABCD$ là hình bình hành tâm $O$, gọi $M$, $N$ lần lượt là trung  điểm của $SA$, $SD$. Số giao tuyến của hai mặt phẳng $(OMN)$ với các mặt phẳng $(SBC)$, $(SAB)$, $(SAD)$ và $(ABCD)$ là
	\shortans{3} 
	 \loigiai{
		\begin{center}
			\begin{tikzpicture}[scale=1,font=\footnotesize,line join=round,line cap=round,>=stealth]
					 \def\bc{4}
					\def\ba{2} 
					\def\h{4} 
					\def\gocB{30} 
					\coordinate[label=below  left:$B$] (B) at (0,0);
					\coordinate[label=above right:$A$]  (A) at (\gocB:\ba);
					\coordinate[label=below:$C$] (C) at (\bc,0); 
					\coordinate[label=right:$D$] (D)  at ($(C)-(B)+(A)$);
					 \coordinate[label=below:$O$] (O) at ($(A)!.5!(C)$);
					\coordinate[label=above:$S$] (S) at ($(O)+(105:\h)$); 
					\coordinate[label=left:$M$] (M) at ($(S)!0.5!(A)$); 
					\coordinate[label=above right:$N$] (N) at  ($(S)!0.5!(D)$);
					\draw (B)--(C)--(D)--(S)--cycle  (S)--(C);
					\draw[dashed] (C)--(A)--(D)--(B) 
						 (S)--(A)--(B)
						(O)--(M)--(N)--cycle; 
					\foreach \diem in {A,B,C,D,S,O} \fill (\diem) circle  (1pt);
				\end{tikzpicture}
		\end{center}
		\begin{itemize}
			\item Ta có $O$, $M$ lần  lượt là trung điểm của $ AC, SA$ nên $OM$ là đường trung bình của $\Delta SAC$ nên $OM \parallel SC \Rightarrow OM \parallel (SBC)$. (1)\\
			$O  N$ là đường trung bình của $\Delta SBD \Rightarrow ON\parallel SB \Rightarrow ON \parallel (SBC)$.(2)\\
			Mặt khác $O M \subset(OMN), ON \subset(OMN)$  và $OM \cap ON=O$.(3)\\
			Từ (1), (2), (3)  suy ra $(OMN) \parallel (SBC)$.
			 \item Vì $\heva{&M\in(OMN)\cap (SAB)\\&M\in(OMN)\cap (SAD)\\&O\in(OMN)\cap (ABCD)}$ nên mặt phẳng $(OMN)$ có giao tuyến với các mặt phẳng $(SAB)$, $(SAD)$ và $(ABCD)$.
		\end{itemize}
	 }
\end{ex}

% nguồn: Ok_TraLoiNgan_Lop11_Lan1
\begin{ex}%[1H4V4-6]%[ma:2c6921e9]
	Cho hình chóp $S.ABC$. Gọi $M$, $N$ lần lượt là trung điểm của $SA$ và $BC$, $P$ là điểm trên  cạnh $AB$ sao cho $\dfrac{AP}{AB}=\dfrac{1}{3}$. Gọi $Q$ là giao điểm của $SC$ với mặt phẳng $(MNP)$. Tính $\dfrac{SC}{SQ}$.
	\shortans{3} 
	 \loigiai{
	\immini 
		{Tìm giao điểm $Q$ của $SC$ với  mặt phẳng $(MNP)$.\\
		Chọn  mặt phẳng phụ $(SAC)$ chứa $SC$.\\
		Trong $(ABC)$ gọi  $H=AC\cap NP$.\\
		Suy  ra $(MNP)\cap(SAC)=HM$.\\
		Khi đó $Q$ là giao điểm của $HM$ và  $SC$.\\
		Gọi $L$ là trung  điểm $AC$.
		} 
		{\begin{tikzpicture}[scale=0.8,font=\footnotesize,line join=round,line cap=round,>=stealth]
		\path 
			(0,0)  coordinate (A)
			(3,0) coordinate (C) 
			 (-50:1.2) coordinate (B)
			 (60:4) coordinate (S)
			($(S)!1/2!(A)$)  coordinate (M)
			($(C)!1/2!(B)$)  coordinate (N)
			($(A)!1/3!(B)$) coordinate  (P)
			($(C)!1/2!(A)$) coordinate (L); 
		\coordinate (H) at (intersection cs:first line={(C)--(A)},second line={(N)--(P)}); 
		\coordinate  (Q) at (intersection cs:first line={(C)--(S)},second line={(H)--(M)});
		\draw (S)--(M)--(P)--(B)--(C)--(S)--(B) 
			(H)--(P) 
			 (H)--(M);
		 \draw[dashed] (M)--(A)--(P)
			 (H)--(C)
			 (P)--(N)
			 (M)--(Q)
			(M)--(L)--(N)--cycle; 
		 \foreach \x/\g in {A/120,B/-90,C/0,S/90,H/180,M/110,N/-45,P/-110,Q/0,L/-110}\fill (\x) circle (1pt)+(\g:3mm) node[black]{$\x$};
	\end{tikzpicture}
		} 
		\noindent Ta có $\dfrac{HA}{HL}=\dfrac{AP}{LN}=\dfrac{\dfrac{1}{3}AB}{\dfrac{1}{2}AB}=\dfrac{2}{3}$ (vì $M$, $N$ là trung điểm của $AC$ và $BC$  nên $LN=\dfrac{1}{2}AB$).\\
		 $\Rightarrow HA=\dfrac{2}{3}HL$.\\
		 Mà $LC=AL=HL-HA=HL-\dfrac{2}{3}HL=\dfrac{1}{3}HL$ nên $ HL=\dfrac{3}{4}HC$.\\
		Mặt khác,  ta có $\dfrac{HC}{HL}=\dfrac{QC}{ML}=\dfrac{4}{3}$ (vì $ ML\parallel SC$).\\
		Mà $2ML=SC$ nên $\dfrac{QC}{SC}=\dfrac{3}{2}\Rightarrow\dfrac{SC}{SQ}=3$. 
	} 
\end{ex}

% nguồn: Ok_TraLoiNgan_Lop11_Lan2
\begin{ex}%[1H4V4-6]%[ma:bc902898]
	Cho hình bình hành $ABCD$. Qua $A$, $B$, $C$, $D$ lần lượt vẽ các nửa đường thẳng $Ax$, $By$, $Cz$, $Dt$ ở cùng phía so với mặt phẳng $\left(ABCD\right)$, song song với nhau và không nằm trong $\left(ABCD\right)$. Một mặt phẳng $(P)$ cắt  $Ax$, $By$, $Cz$, $Dt$ tương ứng tại $A'$, $B'$, $C'$, $D'$ sao cho $AA'=3$, $BB'=5$, $CC'=4$. Tính $DD'$.
	\shortans{2} 
	\loigiai{ 
		\begin{center}
			\begin{tikzpicture}[scale=1,font=\footnotesize,line join=round,line cap=round,>=stealth]
					\path 
						 (0,0) coordinate (B)
						(-1.5,-1.5)  coordinate (A)
						 (4,0) coordinate (C)
						($(A)+(C)-(B)$) coordinate  (D)
						(-0.7,3)  coordinate (x)
						($(B)+(x)-(A)$) coordinate (y) 
						($(C)+(y)-(B)$)  coordinate (z)
						($(z)+(x)-(y)$)  coordinate (t)
						 ($(A)!0.5!(x)$) coordinate (A')
						($(A')+(D)-(A)$)  coordinate (D')
						($(B)!0.7!(y)$)  coordinate (B')
						($(B')+(D')-(A')$) coordinate  (C')
						(intersection  of A--C and B--D) coordinate (O)
						 (intersection of A'--C' and B'--D') coordinate (O');
					\draw 
						(x) node[left]{$x$}--(A)--(D)--(C)--(z) node[right]{$z$} 
						 (A')--(B')--(C')--(D')--cycle (A')--(C')
						 (B')--(D')
						(D)--(t)  node[right]{$t$}
						 (B')--(y) node[left]{$y$};
					\draw[dashed] 
						(A)--(B)--(C)--cycle 
						 (B)--(D)
						(O)--(O') 
						 (B)--(B');
					\foreach \p/\g in {A/200,B/190,C/0,D/-30,O/-90,O'/80,A'/200,B'/180,C'/0,D'/0} 
					\draw[fill=black] (\p) circle (1pt) node[shift=(\g:3mm)] {$\p$}; 
				\end{tikzpicture}
		\end{center}
		Do $(P)$ cắt mặt  phẳng $\left(Ax,By\right)$ theo giao tuyến $A'B'$; cắt mặt phẳng $\left(Cz,Dt\right)$ theo giao tuyến $C'D'$, mà hai mặt phẳng $\left(Ax,By\right)$ và $\left(Cz,Dt\right)$ song song nên $A'B'\parallel C'D'$.\\
		Tương tự có $A'D'\parallel B'C'$ nên $A'B'C'D'$ là hình  bình hành.\\
		Gọi $O$, $O'$ lần lượt là tâm $ABCD$ và $A'B'C'D'$. Dễ  dàng có $OO'$ là đường trung bình của hai hình thang $AA'C'C$ và $BB'D'D$ nên $OO'=\dfrac{AA'+CC'}{2}=\dfrac{BB'+DD'}{2}$.\\
		Từ đó  ta có $DD'=2$.
	} 
\end{ex}

% nguồn: Ok_TraLoiNgan_Lop11_Lan2
\begin{ex}%[1H4V4-6]%[ma:91eb74be]
	Cho hình hộp chữ nhật $ABCD.A'B'C'D'$ có $M$, $N$, $P$ lần lượt là các điểm nằm trên ba  cạnh $AA'$, $BB'$, $CC'$, sao cho $AM=\dfrac{1}{2}AA'$, $BN=\dfrac{1}{3}BB'$, $CP=\dfrac{1}{4}CC'$. Gọi $Q$ là giao điểm của mặt phẳng $\left(MNP\right)$ với đường thẳng $DD'$. Khi đó tỉ số $\dfrac{D'Q}{DD'}$ bằng bao nhiêu?
	\shortans{0,58} 
	 \loigiai{
		\begin{center}
			\begin{tikzpicture}[scale=1,font=\footnotesize,line join=round,line cap=round,>=stealth]
					\path (-0.8,3.8) coordinate  (h)
						(0,0) coordinate (B)++(3.5,0)  coordinate (C)++(0.8,1.2) coordinate (D)
						 ($(B)+(D)-(C)$) coordinate (A);
					\foreach \p in {A,B,C,D}\path ($(\p)+(h)$) coordinate  (\p');
					 \path ($(A)!0.5!(A')$) coordinate (M)
						 ($(B)!1/3!(B')$) coordinate (N)
						($(C)!1/4!(C')$) coordinate  (P)
						($(M)+(P)-(N)$) coordinate  (Q)
						($(D)!0.5!(D')$)  coordinate (M')
						($(C)!1/3!(C')$)  coordinate (N');
					\draw (A')--(B')--(B)--(C)--(D)--(D')--cycle  (B')--(C')--(D')
						 (C)--(C')
						(N)--(P)--(Q) 
						(N)--(N')--(M'); 
					\draw[dashed]  (B)--(A)--(D)
						 (A)--(A')
						(Q)--(M)--(N) 
						 (M)--(M');
					\foreach  \p/\g in {A/60,B/240,C/-60,D/0,A'/90,B'/180,C'/0,D'/45,M/180,N/180,P/0,Q/0,M'/60,N'/150} \fill[black] (\p) circle (1pt)+(\g:0.3) node{$\p$};
				\end{tikzpicture}
		\end{center}
		Lấy $M'$, $N'$  lần lượt trên các cạnh $DD'$ và $CC'$ sao cho $M'D=MA$, $N'C=NB$.\\
		 Ta có
		$\heva{&\left(ABB'A'\right)\parallel \left(CDD'C'\right) \\&\left(\alpha\right)\cap \left(ABB'A'\right)=MN \\&\left(\alpha\right)\cap \left(CDD'C'\right)=PQ}\Rightarrow MN\parallel PQ$.\\ 
		Khi đó 
		 $D'Q=D'M'+QM'=\dfrac{DD'}{2}+\left(N'C-PC\right)=\dfrac{DD'}{2}+\left(\dfrac{DD'}{3}-\dfrac{DD'}{4}\right)=\dfrac{7DD'}{12}$.\\
	Vậy $\dfrac{D'Q}{DD'}=\dfrac{7}{12}$.} 
\end{ex}

% nguồn: Ok_TraLoiNgan_Lop11_Lan2
\begin{ex}%[1H4V4-6]%[ma:3afa3425]
	Cho hình hộp $ABCD.A'B'C'D'$. Một mặt phẳng $\left(\alpha\right)$ cắt các cạnh bên $AA'$, $BB'$, $CC'$, $DD'$ lần lượt tại $M$, $N$, $P$,  $Q$ sao cho $AM=5$, $BN=8$, $CP=7$. Khi đó độ dài đoạn $DQ$ bằng.
	 \shortans{4}
	 \loigiai{
		\begin{center}
			\begin{tikzpicture}[scale=1,font=\footnotesize,line join=round,line cap=round,>=stealth]
					\path (-0.8,4) coordinate  (h)
						(0,0) coordinate  (B)++(3.5,0) coordinate (C)++(0.8,1.5) coordinate (D)
						($(B)+(D)-(C)$)  coordinate (A);
					\foreach \p in {A,B,C,D}\path ($(\p)+(h)$)  coordinate (\p');
					\path ($(A)!0.4!(A')$)  coordinate (M)
						($(B)!0.7!(B')$) coordinate (N) 
						($(C)!0.5!(C')$)  coordinate (P)
						 ($(M)+(P)-(N)$) coordinate (Q)
						($(D)!0.4!(D')$) coordinate (M') 
						 ($(C)!0.7!(C')$) coordinate (N');
					\draw (A')--(B')--(B)--(C)--(D)--(D')--cycle  (B')--(C')--(D')
						(C)--(C') 
						(N)--(P)--(Q) 
						 (N)--(N')--(M');
					 \draw[dashed] (B)--(A)--(D)
						 (A)--(A')
						(Q)--(M)--(N) 
						(M)--(M'); 
					 \foreach \p/\g in {A/150,B/240,C/-60,D/0,A'/90,B'/180,C'/0,D'/45,M/60,N/180,P/-120,Q/30,M'/30,N'/-30} \fill[black] (\p) circle (1pt)+(\g:0.3) node{$\p$};
				\end{tikzpicture}
		\end{center}
		 Ta có
		$\heva{&\left(ABB'A'\right)\parallel \left(CDD'C'\right)  \\&\left(\alpha\right)\cap \left(ABB'A'\right)=MN \\&\left(\alpha\right)\cap \left(CDD'C'\right)=PQ}\Rightarrow MN\parallel PQ.\quad \left(1\right)$\\
		Lấy $M'$, $N'$ lần lượt  trên $DD'$, $CC'$ sao cho $DM'=AM=5$, $CN'=BN=8$.\\
		Ta có 
		$\heva{&BN\parallel CN' \\&BN=CN'}\Rightarrow BNN'C$ là hình bình  hành $\Rightarrow NN'\parallel BC$ và $NN'=BC$.\\
		Từ 
		 $\heva{&DM'\parallel AM \\&DM'=AM}\Rightarrow AMM'D$ là hình bình hành $\Rightarrow AD\parallel MM'$ và $MM'=AD$.\\
		Suy ra $MNN'M'$ là hình  bình hành $M'N'\parallel MN.\quad \left(2\right)$\\
		 Từ $\left(1\right)$, $\left(2\right)$ suy ra $PQM'N'$ là hình bình hành $\Rightarrow QM'=PN'=CN'-CP=8-7=1$.\\
	 Vậy $DQ=DM'-QM'=5-1=4$.}
\end{ex}

% nguồn: Ok_TraLoiNgan_Lop11_Lan2
\begin{ex}%[1H4V4-6]%[ma:1f8d785a]
	Cho hình chóp $S.ABC$ có đáy là tam giác $ABC$ thỏa mãn $AB=AC=4$, $\widehat{BAC}=30^\circ $. Mặt phẳng $(P)$ song song với $\left(ABC\right)$ cắt đoạn $SA$ tại $M$ sao cho $SM=2MA$. Gọi $P$ là giao điểm của mặt phẳng $(P)$ và cạnh $SC$. Tính tỉ số $\dfrac{SP}{SC}$ (kết  quả làm tròn đến chữ số thập phân thứ hai).
	 \shortans{0,67}
	\loigiai{ 
		\begin{center}
			\begin{tikzpicture}[scale=1,font=\footnotesize,line join=round,line cap=round,>=stealth]
					\path 
						(0,0)  coordinate (A)
						(2.5,-1.3)  coordinate (C)
						(4,0) coordinate  (B)
						 (1,3) coordinate (S)
						 ($(S)!2/3!(A)$) coordinate (M)
						($(S)!2/3!(B)$)  coordinate (N)
						 ($(S)!2/3!(C)$) coordinate (P);
					\draw 
						(S)--(A)--(C)--(B)--(S)--(C) 
						(M)--(P)--(N); 
					 \draw[dashed]
						(A)--(B) 
						 (M)--(N);
					\foreach \p/\g  in {A/170,B/0,C/-90,S/90,M/170,P/70,N/10}
					\draw[fill=black] (\p)  circle (1pt) node[shift=(\g:3mm)] {$\p$};
				\end{tikzpicture}
		\end{center}
		Diện tích tam giác $ABC$ là  $S_{\triangle ABC}=\dfrac{1}{2}\cdot AB\cdot AC\cdot \sin\widehat{BAC}=\dfrac{1}{2}\cdot4\cdot4\cdot\sin{30^\circ}=4$.\\
		Gọi $N$, $P$ lần lượt là giao  điểm của mặt phẳng $(P)$ và các cạnh $SB$, $SC$.\\
		Vì $(P)\parallel\left(ABC\right)$ nên theoo định lí  Talet ta có $\dfrac{SM}{SA}=\dfrac{SN}{SB}=\dfrac{SP}{SC}=\dfrac{2}{3}\approx 0{,}67$.
	 }
\end{ex}

% nguồn: Ok_TraLoiNgan_Lop11_Lan2
\begin{ex}%[1H4V4-6]%[ma:211fc141]
	\immini{Một kệ để đồ bằng gỗ có mâm tầng dưới $\left(ABCD\right)$ và mâm tầng trên $\left(EFGH\right)$ song song với nhau. Bác thợ mộc đo được $AE=100$ cm, $CG=120$ cm và muốn  đóng thêm mâm tầng giữa $\left(IJKL\right)$ song song với hai mâm tầng trên, tầng dưới và $EI=42$ cm. Tính độ dài đoạn thẳng $KG$.
	\shortans{50,4}}{ 
	\begin{tikzpicture}[scale=1,font=\footnotesize,line join=round,line cap=round,>=stealth]
			\path 
				(0,0) coordinate (A) 
				 (1,1.3) coordinate (B)
				(4,0) coordinate (D) 
				($(B)+(D)-(A)$) coordinate  (C)
				 (0,4) coordinate (E)
				($(E)+(3,0)$)  coordinate (H)
				($(A)!0.5!(E)$) coordinate  (I)
				 ($(B)+(I)-(A)$) coordinate (L)
				 ($(E)+(L)-(I)$) coordinate (F)
				($(F)+(H)-(E)$)  coordinate (G)
				($(H)!0.5!(D)$) coordinate  (J)
				($(L)+(J)-(I)$)  coordinate (K);
			\filldraw[pattern = north east lines,pattern color = black] (A)--(B)--(C)--(D)--cycle 
				(I)--(L)--(K)--(J)--cycle 
				(E)--(F)--(G)--(H)--cycle; 
			\draw 
				(F)--(E)--(A)--(D)--(C) 
				 (E)--(H)--(G)--(F)
				(I)--(J)--(K) 
				(G)--(C) 
				(H)--(D); 
			 \draw[dashed]
				(A)--(B)--(C) 
				 (B)--(F)
				(I)--(L)--(K); 
			\foreach  \p/\g in {A/200,B/180,C/0,D/-30,I/180,E/180,F/160,L/170,H/0,G/60,J/-10,K/0}
			\draw[fill=black] (\p) circle  (1pt) node[shift=(\g:3mm)] {$\p$};
		\end{tikzpicture}
	 }
	 \loigiai{
		Áp dụng định lý Thales trong  không gian, ta có $\dfrac{KG}{EI}=\dfrac{GC}{EA}\Leftrightarrow KG=\dfrac{42\cdot120}{100}=50{,}4$.\\
		Vậy $KG=50{,}4$ cm. 
	 }
\end{ex}

% nguồn: Ok_TraLoiNgan_Lop11_Lan2
\begin{ex}%[1H4V4-6]%[ma:a164bbca]
	Cho hình chóp $S.ABCD$ có đáy là hình thang, $AB\parallel CD$ và $AB=2CD$. Gọi $O$ là giao điểm của $AC$ và $BD$. Lấy $E$ thuộc cạnh $SA$, $F$ thuộc cạnh $SC$ sao cho $\dfrac{SE}{SA}=\dfrac{SF}{SC}=\dfrac{2}{3}$. Gọi $\left(\alpha\right)$ là mặt phẳng qua $O$ và song song với mặt phẳng $\left(BEF\right)$. Gọi $P$ là giao điểm của $SD$ với $\left(\alpha\right)$. Tính tỉ số $\dfrac{SP}{SD}$  (kết quả làm tròn đến chữ số thập phân thứ hai).
	 \shortans{0,86}
	 \loigiai{
		\begin{center}
			\begin{tikzpicture}[scale=1,font=\footnotesize,line join=round,line cap=round,>=stealth]
					\path 
						(0,0) coordinate  (B)
						(1.5,-2)  coordinate (C)
						(5,0) coordinate (A) 
						(4,-1.8)  coordinate (D)
						(2.5,3.5)  coordinate (S)
						(intersection  of A--C and B--D) coordinate (O)
						($(S)!2/3!(A)$) coordinate  (E)
						($(S)!2/3!(C)$) coordinate  (F)
						(intersection  of S--O and E--F) coordinate (I)
						(intersection of S--D and  B--I) coordinate (N)
						 ($(S)!3/2!(N)$) coordinate (P);
					\draw  (S)--(B)--(C)--(D)--(A)--(S)--(C)--(P)
						 (B)--(F)--(N)--(E)
						(S)--(D); 
					\draw[dashed] 
						(C)--(A)--(B)--(D) 
						(F)--(E)--(B) 
						 (O)--(S)
						 (B)--(N);
					 \foreach \p/\g in {A/0,B/180,C/-120,D/-30,S/90,O/120,F/-40,N/0,E/10,I/140,P/10}
					\draw[fill=black] (\p) circle (1pt)  node[shift=(\g:3mm)] {$\p$};
				\end{tikzpicture}
		\end{center}
		 Vì $\dfrac{SE}{SA}=\dfrac{SF}{SC}=\dfrac{2}{3}$ nên đường thẳng $EF\parallel AC$. Mà $EF\subset\left(BEF\right)$, $AC\not\subset\left(BEF\right)$ nên $AC$ song song với mặt phẳng $\left(BEF\right)$.\\
		Vì $AC$ qua  $O$ và song song với mặt phẳng $\left(BEF\right)$ nên $AC\subset\left(\alpha\right)$.\\
		Trong $\left(SAC\right)$, gọi  $I=SO\cap EF$, trong $\left(SBD\right)$, gọi $N=BI\cap SD$. Suy ra $N$ là giao điểm của đường thẳng $SD$ với mặt phẳng $\left(BEF\right)$.\\
		Hai mặt phẳng song song $\left(BEF\right)$ và $\left(\alpha\right)$ bị cắt bởi mặt phẳng thứ ba là  $\left(SCD\right)$ theo hai giao tuyến lần lượt là $FN$ và $Ct$ nên hai giao tuyến đó song song nhau, tức là $Ct\parallel FN$.\\
		Trong $\left(SCD\right)$, $Ct$ cắt $SD$ tại $P$. Khi đó $P$ là giao điểm của $SD$  với $\left(\alpha\right)$.\\
		Trong hình thang $ABCD$, do  $AB\parallel CD$ và $AB=2CD$ nên $\dfrac{BO}{OD}=\dfrac{AB}{CD}=2\Rightarrow\dfrac{BO}{BD}=\dfrac{2}{3}$.\\
		Trong tam giác $SAC$, có $EF\parallel AC$ nên  $\dfrac{SE}{SA}=\dfrac{SI}{SO}=\dfrac{2}{3}\Rightarrow\dfrac{IS}{IO}=2$.\\
		Xét tam giác $SOD$ với  cát tuyến $NIB$, ta có $\dfrac{NS}{ND}\cdot\dfrac{BD}{BO}\cdot \dfrac{IO}{IS}=1\Rightarrow\dfrac{NS}{ND}=\dfrac{BO}{BD}\cdot \dfrac{IS}{IO}=\dfrac{2}{3}\cdot2=\dfrac{4}{3}$.\\
		Suy ra $\dfrac{SN}{SD}=\dfrac{4}{7}$ và ta  lại có $\dfrac{SN}{SP}=\dfrac{SF}{SC}=\dfrac{2}{3}$ nên suy ra $\dfrac{SP}{SD}=\dfrac{6}{7}\approx 0{,}86$.
	 }
\end{ex}

% nguồn: Ok_TraLoiNgan_Lop11_Lan2
\begin{ex}%[1H4V4-6]%[ma:5f62f4d9]
	Cho hình hộp chữ nhật $ABCD.A'B'C'D'$ có $M$, $N$, $P$ lần lượt là các điểm nằm trên ba cạnh $AA'$, $BB'$, $CC'$, sao cho $AM=\dfrac{1}{2}AA'$, $BN=\dfrac{1}{3}BB'$, $CP=\dfrac{1}{4}CC'$. Gọi $Q$ là giao điểm của mặt phẳng $\left(MNP\right)$ với đường thẳng $DD'$. Khi đó tỉ số $\dfrac{D'Q}{DD'}$  bằng bao nhiêu (kết quả làm tròn đến chữ số thập phân thứ hai)?
	\shortans{0,58} 
	\loigiai{ 
		\begin{center}
			\begin{tikzpicture}[scale=1,font=\footnotesize,line join=round,line cap=round,>=stealth,declare function={ab=2; ad=4; cao=3.5; gocB=35;}]
					\path  (0,0) coordinate (B)
						 (ad,0) coordinate (C)
						(gocB:ab) coordinate  (A)
						($(C)-(B)+(A)$) coordinate (D) 
						($(A)+(90:cao)$) coordinate  (A')
						($(B)+(90:cao)$) coordinate (B') 
						 ($(C)+(90:cao)$) coordinate (C')
						($(D)+(90:cao)$) coordinate  (D')
						($(A)!0.5!(A')$) coordinate (M) 
						($(B)!1/3!(B')$) coordinate (N) 
						($(C)!1/4!(C')$) coordinate  (P)
						($(D')!7/12!(D)$) coordinate (Q); 
					\draw[dashed]  (A')--(A)--(B)
						 (A)--(D)
						 (N)--(M)--(Q);
					\draw  (A')--(B')--(C')--(D')--cycle (N)--(P)--(Q)
						 (D)--(D')--(C')--(C)
						(C')--(B')--(B) 
						 (D)--(C)--(B);
					\foreach \x/\gm in  {A'/120,B'/90,C'/90,D'/90,A/160,B/-90,D/40,C/-90,M/180,N/180,P/0,Q/0} \fill (\x) circle (1pt)
						($(\x)+(\gm:3.5mm)$) node{$\x$}; 
				\end{tikzpicture}
		\end{center}
		Lấy $M'$, $N'$ lần lượt trên các cạnh  $DD'$ và $CC'$ sao cho $M'D=MA$, $N'C=NB$.\\
		 Ta có $\heva{&\left(ABB'A'\right)\parallel \left(CDD'C'\right)\\&\left(\alpha\right)\cap\left(ABB'A'\right)=MN\\&
		\left(\alpha\right)\cap\left(CDD'C'\right)=PQ} \Rightarrow  MN\parallel PQ$.\\
		Khi đó $D'Q=D'M'+QM'=\dfrac{DD'}{2}+\left(N'C-PC\right)=\dfrac{DD'}{2}+\left(\dfrac{DD'}{3}-\dfrac{DD'}{4}\right)=\dfrac{7DD'}{12}$.\\ 
		Vậy $\dfrac{D'Q}{DD'}=\dfrac{7}{12}\approx  0{,}58$.
	 }
\end{ex}

% nguồn: Ok_TraLoiNgan_Lop11_Lan2
\begin{ex}%[1H4V4-6]%[ma:74695932]
	Cho hình chóp $S.ABCD$ có đáy $ABCD$  là hình bình hành, các điểm $M$, $N$ lần lượt thuộc các cạnh $SB$, $AC$ sao cho $\dfrac{BM}{MS}=\dfrac{CN}{NA}=x$, $\left(0 < x\ne 1\right)$. Gọi $G$ là trọng tâm tam giác $SCD$. Tìm $x$ để $\left(MNG\right)\parallel\left(SAD\right)$.
	 \shortans{2}
	 \loigiai{
		\begin{center}
			\begin{tikzpicture}[scale=1,font=\footnotesize,line join=round,line cap=round,>=stealth]
					\path 
						(0,0)  coordinate (B)
						(-1.5,-1.5) coordinate (A) 
						(4,0) coordinate (C) 
						($(A)+(C)-(B)$) coordinate  (D)
						(0.7,3)  coordinate (S)
						($(A)!0.4!(B)$) coordinate (P) 
						($(P)+(D)-(A)$) coordinate  (Q)
						($(S)!0.4!(B)$) coordinate  (M)
						($(S)!0.4!(C)$) coordinate  (K)
						(intersection of  P--Q and A--C) coordinate (N)
						 ($(S)!0.5!(D)$) coordinate (L)
						(intersection of C--L and  K--Q) coordinate (G);
					\draw 
						 (S)--(A)--(D)--(C)--(S)--(D)
						(K)--(Q) 
						 (C)--(L);
					 \draw[dashed]
						(S)--(B) 
						 (A)--(B)--(C)--cycle
						 (Q)--(P)--(M)--(K)
						(M)--(N); 
					 \foreach \p/\g in {A/180,B/-80,C/0,D/-10,S/90,K/20,M/35,P/-60,N/-90,Q/-60,G/30,L/180}
					 \draw[fill=black] (\p) circle (1pt) node[shift=(\g:3mm)] {$\p$};
				\end{tikzpicture}
		\end{center}
		Gọi các giao điểm của $\left(MNG\right)$ với các cạnh hình chóp như hình vẽ.\\ 
		Ta có $\dfrac{BM}{MS}=\dfrac{CN}{NA}=x$, $\left(0 < x\ne 1\right)$ nên $BC$,  $MN$, $SA$ lần lượt nằm trên ba mặt phẳng song song, suy ra $MN\parallel \left(SAD\right)$, $\forall x\in\left(0 < x\ne 1\right)$. Do đó \[\left(MNG\right)\parallel\left(SAD\right)\Leftrightarrow NQ \parallel AD \Leftrightarrow \dfrac{NC}{NA}=\dfrac{QC}{QD}\Leftrightarrow\dfrac{NC}{NA}=\dfrac{GC}{GL}=2.\]
		Vậy với $x=2$ thì $\left(MNG\right)\parallel \left(SAD\right)$. 
	} 
\end{ex}

% nguồn: Ok_TraLoiNgan_Lop11_Lan1
\begin{ex}%[1H4C4-6]%[ma:70316a65]
	Cho hình lăng trụ $ADF.BCE$. Gọi $M$ là trọng tâm $\triangle ABE$. Gọi $(P)$ là mặt  phẳng đi qua $M$ và song song với mặt $(ADF)$. Lấy $N$ là giao điểm của $(P)$ và $AC$. Tính $\dfrac{AN}{NC}$.
	\shortans{2} 
	 \loigiai{
		\begin{center}
			\begin{tikzpicture}[line join = round,line cap = round,font=\footnotesize,>=stealth,scale=1]
					\coordinate  (A) at (3,2);
					 \coordinate (D) at (10,2);
					\coordinate (B) at  (0,0);
					 \coordinate (C) at ($(D)+(B)-(A)$);
					\coordinate (S) at  ($(A)!1/2!(B)$);
					\coordinate (I)  at ($(B)!1/3!(A)$);
					\coordinate  (E) at (2,4);
					 \coordinate (F) at (5,6);
					\coordinate  (J) at ($(E)!1/3!(F)$);
					\coordinate (K) at ($(C)!1/3!(D)$); 
					 \coordinate (N) at ($(C)!1/3!(A)$);
					\coordinate  (M) at ($(I)!1/3!(J)$);
					\draw  (B)--(E)--(C)--(D)
						(B)--(C) 
						(F)--(D) 
						(E)--(F) 
						(K)--(J); 
					 \draw[dashed,thin] (B)--(A)--(F)
						(C)--(A)--(D) 
						(S)--(E) 
						 (K)--(I)--(J)
						(E)--(A); 
					\foreach \i/\g in  {A/90,B/-90,C/-90,D/30,F/90,E/90,S/0,I/-90,J/90,K/-90,N/-90,M/0}{\draw[fill=black](\i) circle (1pt)
						($(\i)+(\g:3mm)$)  node[scale=1]{$\i$};}
				\end{tikzpicture}
		\end{center}
		Vẽ mp$(P)$ chứa $M$ và $(P)\parallel(ADF)$ cắt $AB$, $AC$, $CD$, $EF$ lần lượt tại $I$,  $N$, $K$, $J$.\\
		Ta có: $\dfrac{AI}{BI}=\dfrac{AN}{NC}\;(IN\parallel BC)$; $\dfrac{EJ}{IS}=\dfrac{ME}{MS}=2\;(IS\parallel JE)$\\ 
		$BI=EJ$ (tứ  giác $BIJE$ là hình bình hành)\\
		$\Rightarrow\dfrac{BI}{IS}=2\Rightarrow\dfrac{BI}{2}=\dfrac{IS}{1}=\dfrac{BI+IS}{2+1}=\dfrac{BS}{3}$\\ 
		$\Rightarrow BI=\dfrac{2}{3}BS$; $IS=\dfrac{1}{3}BS$.\\ 
		Ta có: $AI=AS+AI=BS+\dfrac{1}{3}BS=\dfrac{4}{3}BS\Rightarrow\dfrac{AI}{BI}=\dfrac{\dfrac{4}{3}BS}{\dfrac{2}{3}BS}=2\Rightarrow\dfrac{AN}{NC}=2$. 
	} 
\end{ex}

% nguồn: Ok_TraLoiNgan_Lop11_Lan1
\begin{ex}%[1H4C4-6]%[ma:419e26d9]
	Cho hình hộp $ABCD.A'B'C'D'$ có $M$, $N$, $P$ lần lượt là trung điểm của $A'B'$,  $BC$, $DD'$. Gọi $E$, $F$, $G$ lần lượt là giao điểm của $AC'$ với các mặt phẳng $(A'BD)$, $(MNP)$, $(CB'D')$. Tính tỉ số $\dfrac{EG}{EF}$.
	\shortans{2} 
	\loigiai{ 
	\immini{Gọi $H$,  $I$ lần lượt là trung điểm của $B'C$ và $B'D$.\\
		Ta có $HN$ là đường trung bình của tam  giác $CBB'$ nên $HN\parallel BB'\parallel DD'$ và $$HN=\dfrac{1}{2}BB'=\dfrac{1}{2}DD'=D'P.$$
		Suy ra $PNHD'$ là hình bình hành. Do đó $$PN\parallel D'H\Rightarrow PN\parallel (CB'D').$$}{\begin{tikzpicture}[scale=1,font=\footnotesize,line join=round,line cap=round,>=stealth]
		\path (-0.8,3) coordinate  (h)
			(0,0) coordinate (D)++(3.5,0) coordinate (C)++(1.3,1)  coordinate (B)
			($(D)+(B)-(C)$) coordinate  (A);
		\foreach \p in {A,B,C,D}\path ($(\p)+(h)$) coordinate (\p'); 
		\path (barycentric cs:A'=1,B'=1) coordinate  (M)
			 (barycentric cs:C=1,B=1) coordinate (N)
			 (barycentric cs:D=1,D'=1) coordinate (P)
			 (barycentric cs:C=1,B'=1) coordinate (H)
			(barycentric  cs:B'=1,D'=1) coordinate (I);
		\draw (A')--(D')--(D)--(C)--(B)--(B')--cycle  (B')--(D')--(C')--(B')--(C)--(C')
			 (D')--(C)
			 (H)--(N);
		\draw[dashed]  (D)--(A)--(B)
			 (C')--(A)--(A')
			 (M)--(N)--(P)--(M)
			(D')--(H) 
			(C)--(I); 
		\foreach \p/\g in {A/180,D/240,C/-60,B/0,A'/90,D'/180,C'/135,B'/45,M/90,N/-30,P/180,H/30,I/90} \fill[black]  (\p) circle (1pt)+(\g:0.3) node{$\p$};
	\end{tikzpicture}
		}\noindent 
		Chứng minh tương tự,  $MN\parallel (CB'D')$.\\
		 Do đó $(MNP)\parallel (CB'D')$.\\
		 Ta cũng chứng minh được $(MNP)\parallel (A'BD)$.\\
		Ba mặt phẳng $(A'BD)$,  $(MNP)$, $(CB'D')$ song song với nhau nên chắn trên các đường thẳng $A'B'$ và $AC'$ những đoạn thẳng tương ứng tỉ lệ. Tức là $$\dfrac{EG}{A'B'}=\dfrac{EF}{A'M}\Rightarrow\dfrac{EG}{EF}=\dfrac{A'B'}{A'M}=2.$$
	} 
\end{ex}

% nguồn: Ok_TraLoiNgan_Lop11_Lan1
\begin{ex}%[1H4C4-6]%[ma:1fc6084c]
	Cho hình chóp $S.ABCD$ có đáy $ABCD$ là hình bình hành. Gọi $I$, $K$ lần lượt là trung điểm của $BC$ và $CD$. Gọi $M$ là trung điểm của $SB$. Gọi $F$ là giao điểm  của $DM$ và $(SIK)$. Tính tỉ số $\dfrac{MF}{MD}$.
	 \shortans{1}
	\loigiai{ 
		\begin{center}
			\begin{tikzpicture}[line join = round,line cap = round,font=\footnotesize,>=stealth,scale=1]
					\coordinate (A) at  (2,2);
					\coordinate (D) at  (6,2);
					\coordinate  (B) at (0,0);
					\coordinate (C)  at ($(D)+(B)-(A)$);
					 \coordinate (S) at (1,5);
					\coordinate (I)  at ($(B)!0.5!(C)$);
					\coordinate (K) at ($(C)!0.5!(D)$); 
					\coordinate (M) at ($(S)!0.5!(B)$); 
					\coordinate (O)  at ($(A)!0.5!(C)$);
					\coordinate (E) at  ($(I)!0.5!(K)$);
					 \coordinate (F) at ($(M)!-1!(D)$);
					\coordinate (x) at ($(F)!2!(S)$); 
					\draw (B)--(C)--(D) 
						(S)--(B) 
						(S)--(I) 
						 (S)--(C)
						(S)--(K) 
						 (S)--(D)
						(M)--(C) 
						(M)--(F)--(x); 
					\draw[dashed,thin] (D)--(B)--(A)--(D) 
						(I)--(K) 
						(S)--(A)--(C) 
						(A)--(M)--(O) 
						 (M)--(D);
					\foreach \i/\g in {A/90,B/-90,C/-90,D/30,S/90,I/-90,K/-60,M/-145,O/-180,F/-180,E/90}{\draw[fill=black](\i) circle  (1pt)
						($(\i)+(\g:3mm)$) node[scale=1]{$\i$};} 
					\draw (x)  node[below]{$x$};
				\end{tikzpicture}
		\end{center}
		 Ta có $S\in(SIK)\cap(SAC)$.\\
		Trong mặt phẳng $(A B C D)$, gọi $E=IK\cap AC$, suy ra $\heva{&E\in I K\subset(S I K)\\&E\in A  C\subset(S A C).}$\\ Do đó $E\in(S I K)\cap(S A C)$.\\
		Suy ra  $S E=(S I K)\cap(S A C)$.\\
		Ta có $\heva{&S \in(S I K) \cap(S B D) \\ &B D \subset(S B D),\, I  K \subset(S I K)\\&B D\parallel I K} \Rightarrow(S I K) \cap(S B D)=S x,(S x\parallel B D\parallel I K)$.\\
		Trong mp$(S B D)$, gọi $F=S x \cap D M \Rightarrow\left\{\begin{array}{l}S \in D M \\ S \in S x \subset(S I K)
	\end{array} \Rightarrow F=D M \cap(S I K)\right.$.\\
	Ta có  $S F \parallel B D \Rightarrow \dfrac{M F}{M D}=\dfrac{M S}{M B}=1$.
	 }
\end{ex}

% nguồn: Ok_TraLoiNgan_Lop11_Lan2
\begin{ex}%[1H4C4-6]%[ma:b7536289]
	Cho hình chóp $S.ABCD$ có đáy là hình thang có hai đáy là $AB$,  $CD$, $AB=2CD$. Điểm $M$ thuộc cạnh $AD$ ($M$ không trùng với $A$ và $D$) sao cho $\dfrac{MA}{AD}=x$. Gọi $\left(\alpha\right)$ là mặt phẳng qua $M$ và song song với mặt phẳng $\left(SAB\right)$. Tìm $x$ để diện tích thiết diện của hình chóp cắt bởi mặt phẳng $\left(\alpha\right)$ bằng một nửa diện tích tam giác $SAB$.
	 \shortans{0,5}
	 \loigiai{
		\begin{center}
			\begin{tikzpicture}[scale=1,font=\footnotesize,line join=round,line cap=round,>=stealth,declare function={k=0.6;}]
					\path (0,0) coordinate (A)++(3,0) coordinate (D)++(-1.5,-1) coordinate (C) 
						($(A)+(C)-(D)$)  coordinate (b)
						 ($(A)!2!(b)$) coordinate (B)
						(-0.7,2.5) coordinate  (h)
						($(A)+(h)$)  coordinate (S)
						 ($(A)!k!(D)$) coordinate (M)
						 ($(B)!k!(C)$) coordinate (Q)
						 ($(S)!k!(D)$) coordinate (N)
						($(S)!k!(C)$) coordinate  (P)
						(intersection of Q--P and  M--N) coordinate (E);
					\fill[cyan!35] (M)--(N)--(P)--(Q); 
					\draw (S)--(B)--(C)--(D)--cycle  (S)--(C)
						(Q)--(P)--(N) 
						(P)--(E)--(N); 
					\draw[dashed] (S)--(A)--(B) 
						 (A)--(D)
						(Q)--(M)--(N); 
					 \foreach \p/\g in {A/150,B/240,C/-60,D/0,S/90,M/-60,N/30,P/180,Q/-90,E/90} \fill[black] (\p) circle (1pt)+(\g:0.3) node{$\p$};
				\end{tikzpicture}
		\end{center}
		Ta có $\heva{&CD\parallel \left(\alpha\right) \\&CD\subset (ABCD) \\&M\in \left(\alpha\right),M\in (ABCD)}$nên giao tuyến của $\left(\alpha\right)$ và mp$\left(ABCD\right)$ là đường thẳng đi qua $M$ và song  song với $CD$, đường thẳng này cắt $CB$ tại $Q$.\\
		Ta có $\heva{&SA\parallel \left(\alpha\right) \\&SA\subset (SAD) \\&M\in \left(\alpha\right),M\in (SAD)}$nên giao tuyến của $\left(\alpha\right)$ và mp$\left(SAD\right)$ là đường thẳng đi  qua $M$ và song song với $SA$, đường thẳng này cắt $SD$ tại $N$.\\
		Ta có  $\heva{&CD\parallel \left(\alpha\right) \\&CD\subset (SCD) \\&N\in \left(\alpha\right),N\in (SCD)}$nên giao tuyến của $\left(\alpha\right)$ và mp$\left(SCD\right)$ là đường thẳng đi qua $N$ và song song với $CD$, đường thẳng này cắt $SC$ tại $P$.\\
		Ta có $MQ\parallel CD PN\parallel CD$ nên $PN\parallel MQ$. Do đó tứ giác $MNPQ$ là hình  thang.\\
		Thiết diện của hình  chóp cắt bởi mặt phẳng $\left(\alpha\right)$ là hình thang $MNPQ$.\\
		Gọi $E$  là giao điểm của $MN$ và $PQ$.\\
		Ta có $QM=\dfrac{MD}{AD}\cdot  AB+\dfrac{AM}{AD}\cdot CD=\left(1-x\right)AB+xCD=\dfrac{2-x}{2}AB$.\\
		Hai tam giác $SAB$ và $EMQ$ đồng dạng nên $\dfrac{S_{\Delta EMQ}}{S_{\Delta SAB}}=\left(\dfrac{MQ}{AB}\right)^2=\dfrac{\left(2-x\right)^2}{4}$.  $\left(1\right)$\\
		 Vì $\dfrac{NP}{CD}=\dfrac{NS}{SD}=\dfrac{AM}{AD}=x$ nên $NP=xCD=\dfrac{x}{2}AB$.\\
		Do đó $\dfrac{NP}{QM}=\dfrac{x}{2-x}$  và
		$\dfrac{S_{\Delta EPN}}{S_{\Delta EMQ}}=\left(\dfrac{NP}{QM}\right)^2=\dfrac{x^2}{\left(2-x\right)^2}\Rightarrow  \dfrac{S_{MNPQ}}{S_{\Delta EMQ}}=1-\dfrac{x^2}{\left(2-x\right)^2}=\dfrac{4-4x}{\left(2-x\right)^2}$. $\left(2\right)$\\
		Từ $\left(1\right)$  và $\left(2\right)$ suy ra $\dfrac{S_{MNPQ}}{S_{\Delta SAB}}=1-x$.\\
		Do đó $S_{MNPQ}=\dfrac{1}{2}S_{\Delta SAB}\Leftrightarrow  1-x=\dfrac{1}{2}\Leftrightarrow x=\dfrac{1}{2}$.\\
	Vậy  $x=\dfrac{1}{2}$ là giá trị cần tìm.}
\end{ex}

% nguồn: Ok_TraLoiNgan_Lop11_Lan2
\begin{ex}%[1H4C4-6]%[ma:09699b73]
	Cho hình chóp$S.ABCD$ có đáy là  hình thang và $AB\parallel CD$ và $AB=2CD$. Gọi $O$ là giao điểm của $AC$ và $BD$. Lấy$E$ thuộc cạnh $SA$, $F$ thuộc cạnh $SC$ sao cho $\dfrac{SE}{SA}=\dfrac{SF}{SC}=\dfrac{2}{3}$. Gọi $\left(\alpha\right)$ là mặt phẳng qua $O$ và song song với mặt phẳng $\left(BEF\right)$. Gọi $P$ là giao điểm của $SD$ với $(\alpha)$. Tính tỉ số $\dfrac{SP}{SD}$.
	\shortans{0,86} 
	\loigiai{ 
		\begin{center}
			\begin{tikzpicture}[scale=1,font=\footnotesize,line join=round,line cap=round,>=stealth]
					\path (0,0) coordinate (C)++(3,-0.5) coordinate (D)++(1.5,2)  coordinate (A)
						($(A)+(C)-(D)$) coordinate (b) 
						($(A)!2!(b)$) coordinate (B) 
						 (2.5,3) coordinate (h)
						($(B)+(h)$)  coordinate (S)
						 ($(S)!2/3!(A)$) coordinate (E)
						($(S)!2/3!(C)$) coordinate (F) 
						 (intersection of A--C and B--D) coordinate (O)
						(intersection of S--O and E--F)  coordinate (I)
						 (intersection of B--I and S--D) coordinate (N)
						 ($(S)!3/2!(N)$) coordinate (P);
					\draw (S)--(B)--(C)--(D)--(A)--cycle (D)--(S)--(C) 
						(B)--(F)--(N)--(E); 
					\draw[dashed] (C)--(A)--(B)--(D) 
						 (S)--(O)
						(N)--(B)--(E)--(F) 
						 (O)--(P);
					\foreach \p/\g in {A/0,B/180,C/-90,D/-90,S/90,E/30,F/-60,O/-90,I/120,N/0,P/0} \fill[black] (\p)  circle (1pt)+(\g:0.3) node{$\p$};
				\end{tikzpicture}
		\end{center}
		Vì $\dfrac{SE}{SA}=\dfrac{SF}{SC}=\dfrac{2}{3}$  nên đường thẳng $EF\parallel AC$ mà $EF\subset \left(BEF\right)$ và $AC\not\subset \left(BEF\right)$ suy ra $AC\parallel \left(BEF\right)$.\\
		Do  $AC$ qua $O$ và song song với mặt phẳng $\left(BEF\right)$ nên $AC\subset \left(\alpha\right)$.\\
		Trong mặt phẳng $\left(SAC\right)$, gọi $I=SO\cap EF$.\\ 
		Trong $\left(SBD\right)$, gọi $N=BI\cap SD$. Suy ra  $N$ là giao điểm của đường thẳng $SD$ với mặt phẳng $\left(BEF\right)$.\\
		Hai mặt phẳng song song $\left(BEF\right)$ và $\left(\alpha\right)$ bị cắt bởi mặt phẳng thứ ba là $\left(SCD\right)$ theo hai giao tuyến lần lượt là $FN$ và $Ct$ nên hai giao tuyến đó song  song với nhau, tức là $Ct\parallel FN$.\\
		Trong $\left(SCD\right)$,  $Ct$cắt $SD$ tại $P$. Khi đó $P$là giao điểm của $SD$với $\left(\alpha\right)$.\\
		Trong hình thang $ABCD$, do  $AB\parallel CD$ và $AB=2CD$ nên
		$\dfrac{BO}{OD}=\dfrac{AB}{CD}=2\Rightarrow \dfrac{BO}{BD}=\dfrac{2}{3}$.\\ 
		Trong tam giác  $SAC$, có $EF\parallel AC$ nên $\dfrac{SE}{SA}=\dfrac{SI}{SO}=\dfrac{2}{3}\Rightarrow \dfrac{IS}{IO}=2$.\\
		Xét tam giác $SOD$ với cát  tuyến $NIB,$ ta có
		$\dfrac{NS}{ND}\cdot\dfrac{BD}{BO}\cdot\dfrac{IO}{IS}=1\Rightarrow \dfrac{NS}{ND}=\dfrac{BO}{BD}\cdot\dfrac{IS}{IO}=\dfrac{2}{3}\cdot2=\dfrac{4}{3}$.\\ 
		Suy ra $\dfrac{SN}{SD}=\dfrac{4}{7}$. $\left(1\right)$.\\ 
		Ta lại  có $\dfrac{SN}{SP}=\dfrac{SF}{SC}=\dfrac{2}{3}\ \left(\text{Do } CP\parallel FN\right)$. $\left(2\right)$\\
	Từ $\left(1\right)$ và  $\left(2\right)$ suy ra $\dfrac{SP}{SD}=\dfrac{6}{7}\approx 0{,}86$.}
\end{ex}


%%%%%%%%%%%%%%%% TỰ LUẬN %%%%%%%%%%%%%%%%

% nguồn: Ok_TuanLuan_Toan11_Lan2
\begin{ex}%[1H4H4-6]%[ma:954fc58c]
	Cho hình chóp $S.ABCD$ có đáy là hình thang, đáy lớn $AD$ gấp đôi đáy bé $BC$. Gọi $O=AC\cap BD$, $M$ thuộc cạnh $SA$ sao cho $AM=2MS$ và $N$  thuộc cạnh $SB$ sao cho $2BN=NS$.
	\begin{enumerate}
		\item Chứng minh rằng  $\left(OMN\right)\parallel\left(SCD\right)$.
		\item Gọi $d=\left(OMN\right)\cap\left(ABCD\right)$, $P=d\cap AD$, $Q=d\cap  BC$. Chứng minh tứ giác $PQCD$ là hình bình hành.
	\end{enumerate}
	\loigiai{ 
		\begin{enumerate}
			\item 
			Ta  có $AM=2MS\Rightarrow\dfrac{AM}{AS}=\dfrac{2}{3}$ và $2BN=NS\Rightarrow\dfrac{BN}{BS}=\dfrac{1}{3}$.\\
			Xét $\triangle  OAB$ và $\triangle OBC$ có $\heva{
			&AD\parallel  BC\\
			&AD=2BC}\Rightarrow  \heva{
			&OA=2OC\\ 
			 &OD=2OB}\Rightarrow
			\heva{ 
			&\dfrac{AO}{AC}=\dfrac{2}{3}\\ 
			 &\dfrac{BO}{BD}=\dfrac{1}{3}.}$
			\begin{center}
				\begin{tikzpicture}[scale=1,font=\footnotesize,line join=round,line cap=round,>=stealth]
						 \path
						(0,0)  coordinate (A)
						 (2,-1) coordinate (B)
						(5,0) coordinate (D) 
						(4.5,-1) coordinate  (C)
						(A)+(3,4) coordinate (S) 
						(intersection  of A--C and B--D) coordinate (O)
						($(S)!1/3!(A)$) coordinate (M) 
						($(S)!2/3!(B)$)  coordinate (N)
						 ($(B)!1/3!(C)$) coordinate (Q)
						($(Q)+(D)-(C)$) coordinate  (P);
						\draw (S)--(A)--(B)--(C)--(D)--cycle (S)--(B) 
						 (S)--(C)
						 (M)--(N);
						\draw[dashed]  (A)--(D)--(B)
						(A)--(C) 
						(P)--(Q) 
						(O)--(N) 
						(O)--(M); 
						\fill[pattern=dots] (M)--(N)--(O); 
						\fill[pattern=dots]  (S)--(C)--(D);
						 \foreach \p/\r in {S/90,A/180,B/-120,C/-60,D/0,M/120,N/180,P/90,O/-70,Q/-90}
						\fill (\p)  circle (1pt) node[shift={(\r:3mm)}]{$\p$};
					\end{tikzpicture}
			\end{center}
			Trong tam giác $SAC$ có  $\dfrac{AM}{AS}=\dfrac{AO}{AC}=\dfrac{2}{3}$ nên $OM\parallel SC$.\\
			Trong tam giác $SBD$  có $\dfrac{BN}{BS}=\dfrac{BO}{BD}=\dfrac{1}{3}$ nên $ON\parallel SD$.\\
			 Như vậy $\heva{
			 &OM\parallel SC\\
			 &ON\parallel SD\\
			 &OM,\, ON\subset\left(OMN\right)\\
			&SC,\,  SD\subset\left(SCD\right)\\
			&OM\cap ON=O}$  $\Rightarrow\left(OMN\right)\parallel\left(SCD\right)$.
			 \item
			Ta có $\heva{ 
			 &\left(OMN\right)\parallel\left(SCD\right)\\
			 &\left(ABCD\right)\cap\left(SCD\right)=CD\\
			 &\left(ABCD\right)\cap\left(OMN\right)=d}$
			$\Rightarrow d\parallel CD$, trong đó $d$  đi qua $O\in\left(ABCD\right)\cap\left(OMN\right)$.\\
			Xét tứ giác  $PQCD$ có $\heva{
			&PQ\parallel CD\\ 
			 &PD\parallel CQ}\Rightarrow PQCD$ là hình bình hành.
		\end{enumerate}
	 }
\end{ex}

% nguồn: Ok_TuanLuan_Toan11_Lan2
\begin{ex}%[1H4H4-6]%[ma:c70a39af]
	Cho tứ diện $ABCD$.  Lấy $G_1$, $G_2$, $G_3$ lần lượt là trọng tâm của các tam giác $ABC$, $ACD$, $ADB$.
	\begin{enumerate}
		\item Chứng minh rằng $\left(G_1G_2G_3\right)\parallel\left(BCD\right)$. 
		\item  Xác định giao tuyến của mặt phẳng $\left(G_1G_2G_3\right)$ với mặt phẳng $\left(ABC\right)$.
	\end{enumerate}
	 \loigiai{
		\begin{center}
			\begin{tikzpicture}[scale=1,font=\footnotesize,line join=round,line cap=round,>=stealth]
					 \path
						(0,0) coordinate (D) 
						(1,-1) coordinate (B) 
						(4,0) coordinate (C) 
						(2,4)  coordinate (A)
						($(D)!.5!(C)$) coordinate  (N)
						($(D)!.5!(B)$) coordinate (L) 
						 ($(B)!.5!(C)$) coordinate (M)
						($(A)!2/3!(N)$) coordinate  (G_2)
						($(A)!2/3!(L)$) coordinate  (G_3)
						 ($(A)!2/3!(M)$) coordinate (G_1)
						 ($(G_1)+(C)-(M)$) coordinate (x) node[above]{$\Delta$};
					 \draw (A)--(D)--(B)--(C)--cycle (A)--(L)
						 (A)--(M)
						(A)--(B); 
					\draw[dashed] (L)--(N)--(M) 
						 (G_3)--(G_2)--(G_1)--cycle (D)--(C)
						 (A)--(N);
					\draw  (x)--($(x)!2.3!(G_1)$);
					 \foreach \p/\r in {A/90,D/180,B/-90,C/0,M/-45,N/-90,L/-120,G_3/-120,G_1/-45,G_2/120}
					\fill  (\p) circle (1pt) node[shift={(\r:3mm)}]{$\p$};
				\end{tikzpicture}
		\end{center}
		\begin{enumerate}
			\item Gọi $M$, $N$, $L$  lần lượt là trung điểm của các cạnh $BC$, $CD$, $BD$.\\
			Ta có $\dfrac{AG_1}{AM}=\dfrac{AG_2}{AN}=\dfrac{AG_3}{AL}=\dfrac{2}{3}$  suy ra $\heva{
			 &G_1G_2\parallel MN\\
			&G_2G_3\parallel NL.}$\\ 
			 Mà $\heva{
			 &MN\subset\left(BCD\right)\\
			&NL\subset\left(BCD\right)}$ nên  $\heva{
			&G_1G_2\parallel\left(BCD\right)\\ 
			&G_2G_3\parallel\left(BCD\right)}$ suy  ra $\left(G_1G_2G_3\right)\parallel \left(BCD\right)$.
			\item Từ $\left(G_1G_2G_3\right)\parallel\left(BCD\right)$ suy ra  $BC\parallel\left(G_1G_2G_3\right)$.\\
			Ta có  $\heva{
			&G_1\in\left(ABC\right)\cap\left(G_1G_2G_3\right)\\ 
			 &BC\subset\left(ABC\right)\\
			&BC\parallel \left(G_1G_2G_3\right)}$ $\Rightarrow\left(ABC\right)\cap\left(G_1G_2G_3\right)=\Delta  \parallel BC$ với $\Delta$ qua $G_1$.
		\end{enumerate}
	} 
\end{ex}

% nguồn: Ok_TuanLuan_Toan11_Lan2
\begin{ex}%[1H4H4-6]%[ma:444472a8]
	Cho hình chóp $S.ABCD$ có đáy là hình bình hành tâm $O$. Gọi $M$, $N$ lần lượt là trung điểm của $SA$,  $SB$.
	\begin{enumerate}
		\item Chứng minh $\left(OMN\right)\parallel\left(SCD\right)$. 
		\item Gọi $K$ là điểm bất kỳ  trên $MN$. Chứng minh $OK\parallel\left(SCD\right)$.
	\end{enumerate}
	 \loigiai{
		\begin{center}
			\begin{tikzpicture}[scale=1,font=\footnotesize,line join=round,line cap=round,>=stealth]
					\path 
						(0,0) coordinate (A) 
						(-2,-1) coordinate (B) 
						(4,0) coordinate  (D)
						 ($(B)+(D)-(A)$) coordinate (C)
						($(A)!.5!(C)$) coordinate (O) 
						(A)+(2,4)  coordinate (S)
						($(S)!.5!(A)$)  coordinate (M)
						($(S)!.5!(B)$) coordinate  (N)
						 ($(M)!0.3!(N)$) coordinate (K);
					\draw (S)--(B)--(C)--(D)--cycle  (S)--(C);
					\draw[dashed]  (S)--(A)--(C)
						(B)--(A)--(D)--cycle (O)--(M)--(N)--cycle  (O)--(K);
					\foreach  \p/\r in {S/90,A/120,B/-120,C/-60,D/0,N/180,M/0,K/90,O/-90}
					 \fill (\p) circle (1pt) node[shift={(\r:3mm)}]{$\p$};
				\end{tikzpicture}
		\end{center}
		\begin{enumerate}
			\item 
			Ta có $OM$ là đường trung bình  của $\Delta SAC$ nên $OM\parallel SC$.\\
			Khi đó  $\heva{
			 &OM\parallel SC\\
			 &SC\subset\left(SCD\right)\Rightarrow OM\parallel\left(SCD\right)\\
			 &OM\not\subset\left(SCD\right).}$\\
			Tương tự  $\heva{
			 &ON\parallel SD\\
			 &SD\subset\left(SCD\right)\\
			 &ON\not\subset\left(SCD\right)}$ $\Rightarrow ON\parallel\left(SCD\right)$.\\
			Vậy $\heva{ 
			&OM,ON\parallel\left(SCD\right)\\ 
			 &OM,ON\subset\left(OMN\right)\\
			 &OM\cap ON=O}$ $\Rightarrow\left(OMN\right)\parallel\left(SCD\right)$.
			\item Ta  có $\heva{
			&\left(OMN\right)\parallel\left(SCD\right)\\ 
			 &OK\subset\left(OMN\right)}\Rightarrow OK\parallel\left(SCD\right)$.
		\end{enumerate}
	} 
\end{ex}

% nguồn: Ok_TuanLuan_Toan11_Lan2
\begin{ex}%[1H4H4-6]%[ma:02da93c1]
	Cho hình chóp $S.ABC$, trên cạnh $SA$ lấy hai điểm $A_1$, $A_2$ sao cho $A_2A_1=2A_1A$. Gọi $(P)$ và $(Q)$ là hai mặt phẳng lần lượt đi qua $A_1$, $A_2$, đồng thời cùng song song với $\left(ABC\right)$. Mặt phẳng $(P)$ cắt các cạnh $SB$, $SC$ lần lượt tại $B_1$, $C_1$; mặt phẳng $(Q)$ lần lượt cắt các cạnh $SB$, $SC$ lần lượt tại $B_2$, $C_2$. Chứng minh $B_2B_1=2B_1B$ và  $C_2C_1=2C_1C$.
	 \loigiai{
		\begin{center}
			\begin{tikzpicture}[scale=1,font=\footnotesize,line join=round,line cap=round,>=stealth]
					 \path
						(0,0) coordinate (A) node[left]{$A$} 
						(3,-1)  coordinate (B) node[below left]{$B$}
						(4,0) coordinate (C) node[right]{$C$} 
						(2,4) coordinate  (S) node[above]{$S$}
						($(S)!1/3!(A)$) coordinate  (A2) node[left]{$A_2$}
						($(A)!1/3!(A2)$)  coordinate (A1) node[left]{$A_1$}
						($(S)!1/3!(B)$) coordinate (B2) node[below  left]{$B_2$}
						($(B)!1/3!(B2)$)  coordinate (B1) node[below left]{$B_1$}
						($(S)!1/3!(C)$) coordinate (C2)  node[right]{$C_2$}
						($(C)!1/3!(C2)$) coordinate  (C1) node[right]{$C_1$};
					\draw  (S)--(A)--(B)--(C)--cycle (A1)--(B1)--(C1)
						(A2)--(B2)--(C2) 
						(S)--(B); 
					\draw[dashed] (A1)--(C1) 
						 (A2)--(C2)
						(A)--(C); 
					\foreach \x in {A,B,C,A1,A2,B1,B2,C1,C2}{\draw[fill] (\x) circle (1pt);} 
				\end{tikzpicture}
		\end{center}
		Áp dụng định lý Thales cho ba mặt phẳng đôi một song song $(P)$, $(Q)$, $\left(ABC\right)$ và hai cát tuyến $SA$,  $SB$ ta được $\dfrac{A_2A_1}{A_1A}=\dfrac{B_2B_1}{B_1B}$.\\
		 Vì $\dfrac{A_2A_1}{A_1A}=2$ (do $A_2A_1=2A_1A$) nên $\dfrac{B_2B1}{B_1B}=2\Rightarrow{B_2}{B_1}=2B_1B$.\\
		Chứng minh tương tự ta được $C_2C_1=2C_1C$. 
	 }
\end{ex}

% nguồn: Ok_TuanLuan_Toan11_Lan2
\begin{ex}%[1H4H4-6]%[ma:a9d582d8]
	Cho hình chóp tứ giác $S.ABCD$. Lấy các điểm $M$, $N$, $P$  lần lượt là trung điểm của các cạnh $SA$, $SB$, $SC$.
	\begin{enumerate}
		\item  Chứng minh rằng $\left(MNP\right)\parallel\left(ABCD\right)$.
		\item Giả sử mặt phẳng $\left(MNP\right)$ cắt  $SD$ tại $Q$. Chứng minh rằng $Q$ là trung điểm của $SD$.
	\end{enumerate}
	\loigiai{ 
		\begin{center}
			\begin{tikzpicture}[scale=1,font=\footnotesize,line join=round,line cap=round,>=stealth]
					\begin{scope}[scale=.8]
						 \path
						(0,0) coordinate  (A)
						 (4,0) coordinate (B)
						(-2,-1) coordinate  (D)
						 (3,-2) coordinate (C)
						(1,4) coordinate  (S)
						($(S)!.5!(A)$) coordinate (M) 
						($(S)!.5!(B)$)  coordinate (N)
						($(S)!.5!(C)$) coordinate  (P)
						 ($(S)!.5!(D)$) coordinate (Q);
						\draw (S)--(D)--(C)--(B)--cycle  (S)--(C)
						 (Q)--(P)--(N);
						\draw[dashed] (S)--(A)--(C) 
						 (D)--(A)--(B)--cycle (Q)--(M)--(N)
						 (M)--(P);
						\foreach \p/\r  in {S/90,A/-90,D/180,C/-60,B/0,Q/180,M/45,N/60,P/-30}
						\fill  (\p) circle (1pt) node[shift={(\r:3mm)}]{$\p$};
					\end{scope}
				\end{tikzpicture}
		\end{center}
		\begin{enumerate}
			\item Vì $MN$ là đường trung bình  của tam giác $SAB$ nên $MN\parallel AB$.\\
			Vì $AB\subset\left(ABCD\right)$ nên $MN\parallel\left(ABCD\right)$. Chứng minh tương tự ta có $NP\parallel\left(ABCD\right)$ mà $MN$ cắt $NP$ nên ta  có $\left(MNP\right)\parallel\left(ABCD\right)$.
			\item  Vì $Q$ là giao điểm của $SD$ và $\left(MNP\right),M$ là điểm chung của hai mặt phẳng $\left(SAD\right)$ và $\left(MNP\right)$ nên $MQ$ là giao tuyến của $\left(MNP\right)$ và $\left(SAD\right)$.\\
			Do $\left(MNP\right)\parallel\left(ABCD\right)$, $MQ=\left(MNP\right)\cap\left(SAD\right)$, $AD=\left(SAD\right)\cap\left(ABCD\right)$  nên ta có $MQ\parallel AD$.\\
			Trong tam  giác $SAD,M$ là trung điểm của $SA$ và $MQ\parallel AD$ nên $Q$ là trung điểm của $SD$.
		\end{enumerate}
	} 
\end{ex}

% nguồn: Ok_TuanLuan_Toan11_Lan2
\begin{ex}%[1H4V4-6]%[ma:833857b1]
	Cho hình chóp $S.ABCD$ có đáy $ABCD$ là hình thang, $AB=3a$, $AD=CD=a$. Mặt bên $SAB$  là tam giác cân đỉnh $S$ và $SA=2a$, mặt phẳng $(\alpha)$ song song với $(SAB)$ cắt các cạnh $AD$, $BC$, $SC$, $SD$ theo thứ tự tại $M$, $N$, $P$, $Q$.
	\begin{enumerate}[a)]
		\item Chứng minh $MNPQ$ là hình  thang cân.
		\item Đặt $x=AM(0<x<a)$. Tính $x$ để $MNPQ$ là tứ giác ngoại tiếp được một đường tròn. Tính bán  kính đường tròn đó.
		\item Gọi $I=MQ\cap NP$. Tìm tập hợp điểm $I$ khi $M$  di động trên $AD$.
		\item Gọi $J=MP\cap NQ$. Chứng minh $IJ$ có phương không đổi và điểm $J$ luôn thuộc một mặt phẳng cố  định.
	\end{enumerate}
	\loigiai{ 
		\begin{center}
			\begin{tikzpicture}[scale=1,font=\footnotesize,line join=round,line cap=round,>=stealth]
					 \path
						(1,5) coordinate (S) 
						(-2,1)  coordinate (A)
						(4.26,0.98) coordinate (B) 
						(0,-4) coordinate  (E)
						(1.68,-2.04)  coordinate (C)
						 (3.41,-.02) coordinate (N)
						(-.81,-1.98) coordinate (D) 
						 (-1.6,0) coordinate (M)
						 (.39,2.66) coordinate (Q)
						 (.8,3.2) coordinate (I)
						(1.22,2.68)  coordinate (P)
						(.83,2.27) coordinate  (J)
						(1.13,0.99)  coordinate (F)
						(.9,-.01) coordinate (K); 
					\draw(S)--(A)--(E)--(B)--(S)--(E) 
						(S)--(C) 
						 (S)--(D)
						 (M)--(I)--(N);
					\draw[dashed] (A)--(B) 
						(S)--(F)--(E) 
						 (I)--(K)
						(M)--(N) 
						(B)--(C) 
						 (M)--(P)
						 (N)--(Q)--(P)
						(C)--(D); 
					\foreach  \i/\g in {A/-100,B/-90,C/-90,D/-110,M/180,N/0,P/0,Q/180,S/90,I/50,J/0,E/0,F/-60,K/-80}
					 \fill (\i) circle (1pt)+(\g:3mm) node[scale=1]{$\i$};
				\end{tikzpicture}
		\end{center}
		\begin{enumerate}[a)]
			\item Do $\heva{&(\alpha)\parallel (SAB) \\&(ABCD)\cap  (SAB)=AB \\&(ABCD)\cap (\alpha)=MN}\Rightarrow MN\parallel AB\qquad (1)$.\\
			Tương tự $\heva{&(\alpha)\parallel (SAB) \\&(SCD)\cap (ABCD)=CD  \\&(SCD)\cap (\alpha)=PQ}$
			$\Rightarrow PQ\parallel  CD \qquad(2)$.\\
			Lại có $AB\parallel CD \qquad(3)$.\\ 
			Từ $(1)$,  $(2)$ và $(3)$ ta có
			$MN\parallel AB\parallel CD\parallel PQ$ nên $MNPQ$ là hình  thang\qquad (*).\\
			Dễ thấy rằng $MQ\parallel SA$, $NP\parallel  SB$.\\
			Do đó $\dfrac{MQ}{SA}=\dfrac{DM}{DA};\dfrac{NP}{SB}=\dfrac{CN}{CB}$ mà $\dfrac{DM}{DA}=\dfrac{CN}{CB}$ nên $\dfrac{MQ}{SA}=\dfrac{NP}{SB}$.\\ 
			 Mặt khác $\Delta SAB$ cân tại $S\Rightarrow SA=SB\Rightarrow MQ=NP \qquad(**)$.\\
			 Từ $(*)$ và $(**)$ suy ra $MNPQ$ là hình thang cân.
			\item $MNPQ$ là tứ  giác ngoại tiếp $\Leftrightarrow MQ+NP=MN+PQ$.\\
			Ta có  $\dfrac{MQ}{SA}=\dfrac{DM}{DA}=\dfrac{a-x}{a}\Rightarrow MQ=2\left(a-x\right)\Rightarrow NP=2\left(a-x\right)$.\\
			Lại có  $\dfrac{PQ}{CD}=\dfrac{SQ}{SD}=\dfrac{AM}{AD}=\dfrac{x}{a}\Rightarrow PQ=x$.\\
			Không khó khăn ta tính  được $MN=3a-2x$.\\
			Do đó $MQ+NP=MN+PQ\Leftrightarrow 4\left(a-x\right)=3a-2x+x\Leftrightarrow x=\dfrac{a}{3}$.\\ 
			Khi đó tính được  $r=\dfrac{a\sqrt{7}}{6}$.
			\item Gọi $E=AD\cap BC\Rightarrow SE=(SAD)\cap (SBC)$.\\ 
			$I=MP\cap NQ\Rightarrow \heva{&I\in MP\subset (SAD) \\&I\in NQ\subset (SBC)}\Rightarrow  I\in SE$.\\
			 Giới hạn:
			Gọi $I_0$ là giao điểm của $SE$ với mặt phẳng $(\beta)$ đi qua $CD$  và song song với $(SAB)$.\\
			Khi  $M\to D\Rightarrow N\to B\Rightarrow I\to I_0$.\\
			Khi $M\to A\Rightarrow  N\to B\Rightarrow I\to S$.\\
			\item Gọi $K=IJ\cap MN$, vì $MNPQ$ là hình thang cân nên $K$ là trung điểm của $MN$. Gọi $F=EK\cap AB$ thì $F$ là trung điểm của $AB$ nên $F$  cố định.\\
			Dễ thấy $IJ\parallel SF$ suy ra $IJ$ có phương không đổi và điểm $J$ thuộc  mặt phẳng cố định $(SEF)$.
		\end{enumerate}
	} 
\end{ex}

% nguồn: Ok_TuanLuan_Toan11_Lan2
\begin{ex}%[1H4V4-6]%[ma:5ebe44d0]
	Cho mặt phẳng $(\alpha)$ và hai đường thẳng chéo nhau $d_1$, $d_2$ cắt $(\alpha)$ tại $A$, $B$. Đường thẳng $\Delta $ thay đổi luôn song song với $(\alpha)$ cắt $d_1$, $d_2$  lần lượt tại $M$ và $N$. Đường thẳng qua $N$ song song với $d_1$ cắt $(\alpha)$ tại $N'$.
	\begin{enumerate}[a)]
		\item Tứ giác $AMNN'$ là hình  gì? Tìm tập hợp điểm $N'$.
		\item  Xác định vị trí của $\Delta $ để độ dài $MN$ nhỏ nhất.
		\item Gọi $O$ là trung điểm của $AB$, $I$ là trung điểm của  $MN$. Chứng minh $OI$ là đường thẳng nằm trong mặt phẳng cố định khi $M$ di động.
	\end{enumerate}
	 \loigiai{
		\begin{center}
			\begin{tikzpicture}[scale=1,font=\footnotesize,line join=round,line cap=round,>=stealth]
					 \path
						(0,1) coordinate (Q) 
						(-1,1) coordinate  (K)
						(-3.28,-3) coordinate (N1) 
						(4,-3) coordinate (M1) 
						 (6.54,0.98) coordinate (L)
						(5,0.98) coordinate (R) 
						 (0,0) coordinate (A)
						(1.5,-1) coordinate  (O)
						 (3,-2) coordinate (B)
						(,) coordinate (C) 
						(5,0)  coordinate (N')
						 (2.5,0) coordinate (J)
						(0,2)  coordinate (M)
						(2.5,2) coordinate (I) 
						(5,2)  coordinate (N)
						(0,-2.36)  coordinate (A1)
						(0,4) coordinate (A2) 
						(2.66,-2.67) coordinate  (B1)
						 (5.66,3.31) coordinate (B2);
					 \draw(Q)--(K)--(N1)--(M1)--(L)--(R)
						(A)--(A2) 
						(M)--(N)--(N') 
						(B2)--(B) 
						 (A)--(B)--(N')
						(O)--(I); 
					 \draw[dashed] (Q)--(R)
						 (A)--(A1)
						(B)--(B1) 
						(A)--(N') 
						(I)--(J)--(O); 
					\draw  (A2) node[above]{$d_1$};
					 \draw (B2) node[above]{$d_2$};
					 \foreach \i/\g in {A/180,B/-50,N'/0,N/30,I/90,J/-90,M/30,O/-90}
					\fill  (\i) circle (1pt)+(\g:3mm) node[scale=1]{$\i$};
				\end{tikzpicture}
		\end{center}
		\begin{enumerate}[a)]
			\item Ta có  $MA\parallel NN' \qquad(1)$.\\
			Do $\heva{&MN\parallel (\alpha) \\&(AMNN')\cap (\alpha)=AN'}$ 
			$\Rightarrow  AN'\parallel MN \qquad(2)$.\\
			Từ $(1)$, $(2)$ suy ra $AMNN'$ là hình bình  hành.\\
			Gọi $(\beta)$ là mặt phẳng chứa $d_2$ và song song với $d_1$ thì $NN'\subset (\beta)\Rightarrow N'\in (\beta)$ từ đó ta có $N'$ thuộc giao tuyến $d_3$ của $(\alpha)$ và $(\beta)$. 
			\item Ta có $MN=AN'$ nên $MN$ nhỏ nhất khi $AN'$ nhỏ  nhất $\Leftrightarrow AN'\perp d_3$.\\
			 Từ đó ta xác định $\Delta $ như sau:
			\begin{itemize}
				\item Dựng  $(\beta)$ chứa $d_2$ và $(\beta)\parallel d_1$.
				\item Dựng giao tuyến $d_3=(\alpha)\cap (\beta)$. 
				\item Gọi $N'$  là hình chiếu của $A$ trên $d_3$.
				\item Từ $N'$ dựng đường  thẳng song song với $d_1$ cắt $d_2$ tại $N$.
				\item Từ $N$  dựng đường thẳng $\Delta $ song song với $N'A$ thì $\Delta $ là đường thẳng thỏa yêu cầu bài toán.
			\end{itemize}
			\item Gọi $J$ là trung điểm  của $AN'$ thì $(OIJ)\parallel (\beta)$ mà $O$ cố định và $(\beta)$ cố định nên $(OIJ)$ cố định. Vậy $OI$ thuộc mặt phẳng cố định đi qua $O$ và song song với $(\beta)$.
		\end{enumerate}
	} 
\end{ex}

% nguồn: Ok_TuanLuan_Toan11_Lan2
\begin{ex}%[1H4V4-6]%[ma:47ffd383]
	Cho hình hộp $ABCD.A'B'C'D'$. 
	\begin{enumerate}[a)]
		\item Chứng minh $(BDA')\parallel  (B'D'C)$.
		\item Chứng minh đường chéo $AC'$ đi qua trọng  tâm $G_1$, $G_2$ của các tam giác $BDA'$, $B'D'C$ đồng thời chia đường chéo $AC'$ thành ba phần bằng nhau.
		\item Xác định thiết diện của  hình hộp cắt $(A'B'G_2)$. Thiết diện là hình gì?
	\end{enumerate}
	 \loigiai{
		\begin{center}
			\begin{tikzpicture}[scale=1,font=\footnotesize,line join=round,line cap=round,>=stealth]
					\path 
						 (-2,2) coordinate (A)
						(3,2) coordinate  (B)
						(5,5) coordinate (C) 
						(0,5) coordinate  (D)
						(-2,-3) coordinate  (A')
						(3,-3) coordinate (B') 
						 (5,0) coordinate (C')
						(0,0) coordinate  (D')
						(0,2.5) coordinate (F) 
						 (5,2.5) coordinate (E)
						(1.5,3.5)  coordinate (O)
						(1.5,-1.5) coordinate (O') 
						(2.67,.67) coordinate  (G_2);
					\coordinate (G_1) at (intersection of A'--O and A--C'); 
					\coordinate (I) at (intersection of E--F and  D'--C);
					 \draw(A)--(B)--(C)--(D)--(A)--(C)--(C')--(B')--(B)--(A')--(B')--(C)
						(B')--(E) 
						(A)--(A') 
						 (B)--(D);
					 \draw[dashed] (D)--(D')--(C')--(A')--(D')--(B')
						 (A')--(D)
						 (O)--(A')--(F)--(E)
						(O')--(C)--(D') 
						 (A)--(C');
					\foreach \i/\g in  {A/180,B/90,C/0,D/30,A'/180,B'/-70,C'/30,D'/-90,O/90,O'/-90,G_2/110,G_1/20,E/0,F/180,I/90}
					\fill (\i)  circle (1pt)+(\g:3mm) node[scale=1]{$\i$};
				\end{tikzpicture}
		\end{center}
		\begin{enumerate}[a)]
			\item Gọi $O$, $O'$ lần lượt là  trọng tâm các mặt $ABCD$ và $A'B'C'D'$.\\
			Dễ thấy $DBB'D'$ là hình bình  hành \\
			nên  $B'D'\parallel BD\subset (BDA')\Rightarrow B'D'\parallel (BDA')\qquad (1)$.\\
			Tương tự $OCO'A'$ là hình bình hành nên $O'C\parallel OA'\subset (A'BD)$ 
			$\Rightarrow CO'\parallel  (A'BD)\qquad (2)$.\\
			Từ $(1)$, $(2)$ suy ra $(A'BD)\parallel (CB'D')$. 
			\item Ta có $A'O$ là trung tuyến của tam giác $A'BD$ và $\dfrac{G_1O}{G_1A'}=\dfrac{OA}{A'C'}=\dfrac{1}{2}$ nên $G_1$ là trọng tâm  của tam giác $A'BD$.\\
			Tương tự $G_2$  cũng là trọng tâm của tam giác $CB'D'$.\\
			Dễ thấy $OG_1$ và $O'G_2$ là đường trung bình của các tam giác $ACG_2$  và $A'C'G_1$ nên
			$AG_1=G_1G_2=G_1C'=\dfrac{1}{3}AC'$. 
			\item Gọi  $I$ là trung điểm của $CD'$.\\
			Do  $G_2$ là trọng tâm tam giác $CB'D'$ nên $I\in B'G_2\subset (A'B'G_2)$.\\
			Vậy  $\heva{&I\in (A'B'G_2)\cap (CDD'C') \\&A'B'\parallel C'D' \\&A'B'\subset (A'B'G_2) \\&C'D'\subset (CDD'C')}\Rightarrow (A'B'G_2)\cap (CDD'C')=EF\parallel C'D', E\in CC',F\in DD'$.\\
			Thiết diện là hình bình hành  $A'B'EF$.
		\end{enumerate}
	} 
\end{ex}

% nguồn: Ok_TuanLuan_Toan11_Lan2
\begin{ex}%[1H4V4-6]%[ma:a0a326b9]
	Cho hình lăng trụ $ABC.A'B'C'$. 
	\begin{enumerate}[a)]
		\item Gọi $I$, $K$, $G$  lần lượt là trọng tâm các tam giác $ABC$, $A'B'C'$ và $ACC'$. Chứng minh $(IGK)\parallel (BB'C'C)$ và $(A'KG)\parallel (AIB)$.
		\item Gọi $P$, $Q$ lần lượt là trung điểm của  $BB'$ và $CC'$. Hãy dựng đường thẳng đi qua trọng tâm của tam giác $ABC$ cắt $AB'$ và $PQ$.
	\end{enumerate}
	\loigiai{ 
		\begin{center}
			\begin{tikzpicture}[scale=1.2,font=\footnotesize,line join=round,line cap=round,>=stealth]
					 \path
						(0,0) coordinate (C) 
						(0,2.5) coordinate (E) 
						 (0,5) coordinate (B)
						(4,5)  coordinate (B')
						(6,2) coordinate (A') 
						(4,0) coordinate (C') 
						 (4,2.5) coordinate (F)
						(0,2.5) coordinate  (E)
						(4.67,2.33) coordinate (K) 
						(0.67,2.33) coordinate  (I)
						 (3,1) coordinate (O)
						 (2,0.67) coordinate (G)
						 (1,1) coordinate (M)
						(2,2) coordinate (A); 
					 \draw(C')--(C)--(B)--(B')--(C')--(B)
						(C)--(F)--(E)--(B')--(A')--(F) 
						(A')--(C'); 
					 \draw[dashed] (C)--(A)--(A')--(C)
						(B)--(A)--(B') 
						(C')--(M)--(B) 
						(C')--(A) 
						 (I)--(G)--(K);
					\foreach \i/\g in {A/180,B/90,C/180,A'/0,B'/0,C'/0,O/-90,G/-90,E/180,F/60,I/0,K/80,M/180} 
					\fill (\i)  circle (1pt)+(\g:3mm) node[scale=1]{$\i$};
				\end{tikzpicture}
		\end{center}
		\begin{enumerate}[a)]
			\item Gọi $O$, $M$, $E$, $F$lần lượt là trung điểm của $AC'$, $AC$, $BC$,  $B'C'$.\\
			 Chứng minh $(IGK)\parallel (BCC'B')$.\\
			Ta  có $\dfrac{MI}{MB}=\dfrac{MG}{MC'}\Rightarrow IG\parallel CC'\subset (BCC'B')\Rightarrow IG\parallel (BCC'B') \qquad(1)$.\\
			Tương tự $\dfrac{A'G}{A'C}=\dfrac{OA'+\dfrac{1}{3}OA'}{A'C}$ 
			$=\dfrac{\dfrac{4}{3}OA'}{A'C}=\dfrac{2}{3}$.\\ 
			Lại  có $\dfrac{A'K}{A'F}=\dfrac{2}{3}\Rightarrow \dfrac{A'G}{A'C}=\dfrac{A'K}{A'F}$
			$\Rightarrow GK\parallel CF\subset (BCC'B')\Rightarrow GK\parallel (BCC'B') \qquad (2)$.\\ 
			Từ  $(1)$, $(2)$ suy ra $(IGK)\parallel (BCC'B')$.\\
			Chứng minh  $(A'KG)\parallel (AIB')$.\\
			Dễ thấy $AA'FE$ là hình bình hành nên $A'F\parallel  AE$ hay $A'F\parallel (AIB') \qquad(3)$.\\
			Cũng dễ thấy $CF\parallel  EB'\subset (AIB')\Rightarrow CF\parallel (AIB') \qquad(4)$.\\
			Từ $(3)$, $(4)$ suy ra $(A'CF)\parallel (AIB')$ mà $(A'CF)$  chính là $(A'KG)$ nên $(A'KG)\parallel (AIB')$.
			\begin{center}
				\begin{tikzpicture}[scale=.81,font=\footnotesize,line join=round,line cap=round,>=stealth]
						\path 
						(2,2)  coordinate (A)
						(1,1)  coordinate (M)
						(0,0) coordinate (C) 
						(0,2.5)  coordinate (E)
						(0.67,2.33)  coordinate (I)
						(0,5) coordinate  (B)
						 (2.5,5) coordinate (P)
						 (2.5,3.75) coordinate (R)
						(2.5,0)  coordinate (Q)
						 (7,2) coordinate (A')
						(5,5)  coordinate (B')
						(5,0)  coordinate (C')
						(8,8)  coordinate (S)
						 (4.12,5) coordinate (X);
						\draw(C')--(C)--(B)--(B')--(C')--(A')--(B') 
						 (P)--(Q)
						(E)--(B')--(S)--(X); 
						\draw[dashed]  (A')--(A)--(C)
						 (B)--(A)--(B')--(Q)--(I)--(X)
						(A)--(E) 
						 (B)--(M);
						\foreach \i/\g in {A/-90,B/180,C/-90,M/-30,E/180,I/-160,P/90,R/-50,Q/-90,A'/0,B'/0,C'/0,S/-90} 
						\fill (\i) circle  (1pt)+(\g:3mm) node[scale=1]{$\i$};
					\end{tikzpicture}
			\end{center}
			\item Trong  $(BCC'B')$ gọi $R=PQ\cap B'E$
			$\Rightarrow \heva{&R\in PQ  \\&R\in B'E\subset (AB'E)}$.\\
			Trong $(AB'E)$ gọi $S=IR\cap AB'$ thì đường thẳng $IR$ chính là đường thẳng cần dựng. 
		\end{enumerate}
	 }
\end{ex}

% nguồn: Ok_TuanLuan_Toan11_Lan2
\begin{ex}%[1H4V4-6]%[ma:43f8a016]
	Cho hình hộp  $ABCD.A'B'C'D'$ có tất cả các mặt đều là hình vuông cạnh $a$.Trên các cạnh $AB$, $CC'$, $C'D'$và $AA'$ lấy các điểm $M$, $N$, $P$, $Q$ sao cho $AM=C'N=C'P=AQ=x\,\left(0\le x\le a\right)$.
	\begin{enumerate}[a)]
		\item Chứng minh bốn điểm $M$, $N$, $P$, $Q$ đồng phẳng và $MP$,  $NQ$ cắt nhau tại một điểm cố định.
		\item Chứng minh $(MNPQ)$ đi  qua một đường thẳng cố định.
		\item Dựng thiết diện của hình hộp khi cắt bởi $(MNPQ)$. Tìm giá trị lớn nhất và giá trị nhỏ nhất của  chu vi thiết diện.
	\end{enumerate}
	 \loigiai{
		\begin{center}
			\begin{tikzpicture}[scale=1,font=\footnotesize,line join=round,line cap=round,>=stealth]
					 \path
						(0,-2) coordinate  (A)
						(5,-2) coordinate  (B)
						(7,1) coordinate  (C)
						 (2,1) coordinate (D)
						 (0,3) coordinate (A')
						(5,3) coordinate  (B')
						(7,6) coordinate (C') 
						(2,6) coordinate (D') 
						(3.5,2) coordinate  (O)
						(3,-2)  coordinate (M)
						(7,3) coordinate (N) 
						(4,6) coordinate (P) 
						(0,1) coordinate (Q) 
						(6,-.5)  coordinate (R)
						(1,4.5) coordinate  (S);
					\draw(A')--(B')--(C')--(D')--(A')--(A)--(B)--(B')--(A')--(C')--(C)--(B) 
						 (R)--(N)
						(S)--(P) 
						 (Q)--(M);
					\draw[dashed] (A)--(D)--(D') 
						(D)--(C) 
						 (P)--(M)--(N)--(P)--(Q)--(N)
						 (M)--(R)--(S)--(Q);
					\foreach \i/\g  in {A/180,B/-90,C/0,D/30,A'/180,B'/0,C'/30,D'/30,O/40,M/-90,N/0,P/90,Q/130,R/0,S/90}
					\fill (\i) circle (1pt)+(\g:3mm)  node[scale=1]{$\i$};
				\end{tikzpicture}
		\end{center}
		\begin{enumerate}[a)]
			\item Dễ thấy $PN\parallel CD'$ và $QM\parallel A'B$ mà $A'B\parallel C'D$ nên $PN\parallel QM$ hay $M$, $N$, $P$, $Q$ đồng  phẳng.
			\item Do $PC'MA$ là hình bình hành nên $MP$ đi qua  trung điểm $O$ của $AC'$.\\
			$\Rightarrow O\in  (MNPQ)$.\\
			Mặt khác  $A'B\parallel MQ\subset (MNPQ)$
			 $\Rightarrow A'B\parallel (MNPQ)$.\\
			Gọi $\Delta $ là đường thẳng qua  $O$ và song song với $A'B$ thì $\Delta $ cố định và $\Delta \subset (MNPQ)$. Hay $(MNPQ)$ luôn chứa đường thẳng cố định $\Delta $.\\
			$(MNPQ)\parallel (A'BC')\Rightarrow  BC'\parallel (MNPQ)\Rightarrow BC'\parallel NR$
			$\Leftrightarrow \dfrac{BR}{BC}=\dfrac{C'N}{CC'}\Rightarrow x=\dfrac{a}{2}$. Đảo lại $x=\dfrac{a}{2}$, dễ dàng  chứng minh được $(MNPQ)\parallel (A'BC')$.
			\item Dễ thấy $\Delta $ cắt $BC$,  $A'D'$ tại các trung điểm $R$ và $S$ của chúng.\\
			Thiết diện là lục giác $MPNPSQ$. Dễ thấy lục giác có tâm đối xứng  là $O$ nên $MQ=NP$, $MR=NS$, $RN=SQ$ do đó chu vi thiết diện là
			$2p=2\left(RM+MQ+QS\right)$.\\ 
			 Ta có $MR=QS=\sqrt{\dfrac{a^2}{4}+\left(a-x\right)^2}$, $QM=x\sqrt{2}$.\\
			Vậy $2p=2\left(x\sqrt{2}+2\sqrt{\dfrac{a^2}{4}+\left(a-x\right)^2}\right)$.\\ 
			Đặt $f(x)=x\sqrt{2}+\sqrt{a^2+4\left(a-x\right)^2};x\in [0;a]$.\\ 
			Theo Cauchy--Schwarz 
			\[\sqrt{\left(a^2+4\left(a-x\right)^2\right)\left(1^2+1^2\right)}\ge a+2\left(a-x\right)\Rightarrow \sqrt{a^2+4\left(a-x\right)^2}\ge  \dfrac{1}{\sqrt{2}}\left(3a-2x\right)\]
			Nên $f(x)\ge x\sqrt{2}+\dfrac{1}{\sqrt{2}}\left(3a-2x\right)=\dfrac{3a}{\sqrt{2}}$. Đẳng thức xảy ra khi  $x=\dfrac{a}{2}$.\\
			Vậy  $\min \left(2p\right)=3\sqrt{2}a$.\\
			Mặt khác bằng biến đổi  tương đương
			\[x\sqrt{2}+\sqrt{a^2+4\left(a-x\right)^2}\le \sqrt{2}a+a\Leftrightarrow{\left(a-x\right)^2}\left[\left(a-x\right)^2-a^2\right]\le  0\]
			đúng  $\forall x\in \left[0;a\right]$.\\
			Đẳng thức xảy ra khi $x=a$.Vậy $\max \left(2p\right)=2a\left(\sqrt{2}+1\right)$. 
		\end{enumerate}
	 }
\end{ex}

% nguồn: Ok_TuanLuan_Toan11_Lan2
\begin{ex}%[1H4V4-6]%[ma:d65f808a]
	Cho hình chóp cụt $ABC.A'B'C'$. Gọi $M$, $N$, $P$ lần lượt là trung điểm  của các cạnh $A'B'$, $BB'$, $BC$.
	\begin{enumerate}[a)]
		\item Xác định thiết diện của hình chóp cụt với  $(MNP)$.
		\item Gọi $I$ là trung điểm của $AB$. Tìm giao điểm của  $IC'$ với $(MNP)$.
	\end{enumerate}
	 \loigiai{
		\begin{center}
			\begin{tikzpicture}[scale=1.2,font=\footnotesize,line join=round,line cap=round,>=stealth]
					\path 
						(1.54,3.4) coordinate (A') 
						(3.04,1.9) coordinate (B') 
						(5.54,3.4) coordinate (C') 
						(2.87,3.4) coordinate (R) 
						(2.29,2.65) coordinate  (M)
						(3.92,1.44) coordinate  (H)
						(0,0) coordinate  (A)
						(1.5,-1.5) coordinate (I) 
						 (3,-3) coordinate (B)
						(3.75,-3.75)  coordinate (J)
						(5.5,-1.5) coordinate  (P)
						 (8,0) coordinate (C)
						(6.67,0) coordinate  (Q)
						(6.38,-.38) coordinate (K) 
						(3.02,-.55) coordinate  (N);
					 \draw(A')--(B')--(C')--(A')--(A)--(J)--(M)--(A')
						 (I)--(M)--(B')--(B)
						 (R)--(M)--(C')--(C)--(P)--(N)
						(J)--(P); 
					\draw[dashed]  (A)--(C)--(I)--(C')
						(B)--(P)--(Q)--(R) 
						(M)--(K); 
					 \foreach \i/\g in {A/180,B/-90,C/0,A'/90,B'/-50,C'/0,M/90,N/30,P/-30,Q/90,R/90,I/-110,H/0,J/0,K/-70}
					\fill (\i) circle (1pt)+(\g:3mm)  node[scale=1]{$\i$};
				\end{tikzpicture}
		\end{center}
		\begin{enumerate}[a)]
			\item Trong $(ABB'A')$ gọi  $J=MN\cap AB$, trong $(ABC)$ gọi $Q=JP\cap AC$.\\
			Ta có $(ABC)\parallel (A'B'C')$  nên
			$(MNP)\cap (A'B'C')=MR\parallel PQ$.\\ 
			Thiết diện là ngũ  giác $MNPQR$.
			\item Trong $(ABC)$ gọi $K=PQ\cap  IC$ thì $K\in (MNP)\Rightarrow MK\subset (MNP)$.\\
			Do $CI\parallel  C'M$ nên trong $(MICC')$ gọi $H=IC'\cap MK\Rightarrow H=IC'\cap (MNP)$.
		\end{enumerate}
	 }
\end{ex}

% nguồn: Ok_TuanLuan_Toan11_Lan2
\begin{ex}%[1H4V4-6]%[ma:05f7fff9]
	Cho hình chóp $S.ABC$, một mặt phẳng $(\alpha)$ di động luôn song song với $(ABC)$, cắt $SA$, $SB$, $SC$ lần lượt tại  $A'$, $B'$, $C'$. Tìm tập hợp điểm chung của ba mặt phẳng $(A'BC)$, $(B'AC)$, $(C'AB)$.
	\loigiai{ 
	 \immini{\textbf{Bổ đề:}
		Cho tam giác $ABC$ các điểm $M$, $N$ thuộc các cạnh $AB$, $AC$ sao cho $MN\parallel BC$. Gọi $E$, $F$ lần lượt là trung điểm  của $BC$, $MN$ và $I=MB\cap CN$ thì $A$, $F$, $I$, $E$ thẳng hàng.\\
		 }
		{\begin{tikzpicture}[scale=1.2,font=\footnotesize,line join=round,line cap=round,>=stealth]
		 \path
			(1,4) coordinate (A) 
			(0,2.5)  coordinate (M)
			(-1,1)  coordinate (B)
			(2,1) coordinate (E) 
			(5,1) coordinate (C) 
			(3,2.5) coordinate (N) 
			(1.5,2.5) coordinate  (F)
			(1.67,2)  coordinate (I);
		 \draw(A)--(B)--(C)--(A)
			(B)--(N)--(M)--(C) 
			 (A)--(E);
		\foreach  \i/\g in {A/90,B/-90,C/-90,M/110,N/40,E/-90,F/70,I/60}
		\fill (\i) circle  (1pt)+(\g:3mm) node[scale=1]{$\i$};
	\end{tikzpicture}
		} 
		\textbf{Chứng  minh:} Ta có $2\overrightarrow{AE}=\overrightarrow{AB}+\overrightarrow{AC}=\dfrac{AB}{AM}\overrightarrow{AM}+\dfrac{AC}{AN}\overrightarrow{AN}$
		$=k\left(\overrightarrow{AM}+\overrightarrow{AN}\right)=2k\overrightarrow{AF}$.\\ 
		Với  $k=\dfrac{AB}{AM}=\dfrac{AC}{AN}$.\\
		Hay $A$, $E$, $F$ thẳng  hàng.\\
		Mặt khác $2\overrightarrow{IE}=\overrightarrow{IB}+\overrightarrow{IC}=-\dfrac{IB}{IN}\overrightarrow{IN}-\dfrac{IC}{IM}\overrightarrow{IM}$ 
		$=l\left(\overrightarrow{IN}+\overrightarrow{IM}\right)=2l\overrightarrow{IF}$ \\ 
		 với $l=-\dfrac{IB}{IN}=-\dfrac{IC}{IM}\Rightarrow I$, $E$, $F$ thẳng hàng.\\
		Vậy $A$, $F$, $I$, $E$  thẳng hàng.\\
		\begin{center}
			\begin{tikzpicture}[scale=1,font=\footnotesize,line join=round,line cap=round,>=stealth]
					 \path
						(1,5) coordinate (S) 
						(0.48,2.4) coordinate (A') 
						 (0,0) coordinate (A)
						(1,-2) coordinate (F) 
						 (2,-4) coordinate (B)
						(5,-2) coordinate  (E)
						(8,0) coordinate (C) 
						(4.64,2.4) coordinate (C') 
						(1,.21) coordinate  (M)
						(3.74,.21)  coordinate (N)
						(1.52,.31)  coordinate (B')
						(3.33,-1.33) coordinate (G) 
						(2.79,.15) coordinate (I) 
						(3.05,1.58) coordinate (P); 
					 \draw(S)--(A)--(B)--(C)--(S)--(B)
						 (A)--(B')--(A')--(B)--(C')--(B')--(C);
					\draw[dashed] (A')--(C')--(A)--(C)--(F) 
						 (S)--(G)
						(B)--(P) 
						 (A')--(C)
						 (A)--(E);
					\foreach \i/\g in {A/-110,B/-90,C/-90,S/30,A'/180,B'/70,C'/30,M/40,N/40,P/90,I/180,G/-90,E/-60,F/180} 
					\fill (\i)  circle (1pt)+(\g:3mm) node[scale=1]{$\i$};
				\end{tikzpicture}
		\end{center}
		Quay lại bài  toán:
		Gọi $M=AB'\cap BA'$, $P=AC'\cap CA'$, $N=BC'\cap CB'$  và $I=CM\cap AN$.\\
		Khi đó $\heva{&I\in AN\subset (ABC') \\&I\in CM\subset (BCA')}\Rightarrow  I\in BP=(ABC')\cap (BCA')$.\\
		Vậy  $I$ chính là điểm chung của ba mặt phẳng $(A'BC)$, $(B'AC)$, $(C'AB)$.\\
		Gọi $E$, $F$ lần lượt là  trung điểm của $BC$, $BA$.\\
		Theo bổ đề trên ta có $S$, $N$, $E$ thẳng hàng và  $I\in AN$ nên $I\in (SAE)$.\\
		Tương tự $I\in (SCF)$. Gọi $G$ là trọng tâm của  $\Delta ABC$ thì $SG=(SAE)\cap (SCF)$ nên $I\in SG$.\\
		Từ đó dễ dàng lập luận được quỹ tích điểm $I$ là đoạn thẳng $SG$ trừ $S$  và $G$.
	} 
\end{ex}

% nguồn: Ok_TuanLuan_Toan11_Lan2
\begin{ex}%[1H4V4-6]%[ma:d3b53a92]
	Hai hình vuông và $ABEF$ ở trong hai mặt phẳng khác nhau. Trên các đường chéo  $AC$ và $BF$ lần lượt lấy các điểm $M$, $N$ sao cho $AM=BN$. Các đường thẳng song song với $AB$ vẽ từ $M$, $N$ lần lượt cắt $AD$, $AF$ tại $M'$, $N'$.
	\begin{enumerate}[a)]
		\item Chứng minh $(BCE)\parallel  (ADF)$.
		\item Chứng minh $(DEF)\parallel  (MNN'M')$.
		\item Gọi $I$ là trung điểm của $MN$. Tìm tập hợp điểm $I$  khi $M$, $N$ thay đổi trên $AC$ và $BF$.
	\end{enumerate}
	 \loigiai{
		\begin{center}
			\begin{tikzpicture}[scale=1.2,font=\footnotesize,line join=round,line cap=round,>=stealth]
					\path 
						(2,2) coordinate  (A)
						(7,2) coordinate  (B)
						(5,0) coordinate  (C)
						(0,0)  coordinate (D)
						(6,5) coordinate  (E)
						(1,5) coordinate (F) 
						(4.5,2) coordinate (J) 
						(2.98,1.35)  coordinate (M)
						(5.04,2.98)  coordinate (N)
						(6.35,1.35) coordinate  (P)
						 (6.67,2.98) coordinate (Q)
						(1.35,1.35) coordinate (M') 
						(1.67,2.98) coordinate (N') 
						(3.35,2.98) coordinate  (X)
						(4.66,1.35) coordinate  (Y)
						(3,2.5)  coordinate (K)
						 (4.01,2.16) coordinate (I);
					 \draw(C)--(D)--(F)--(E)--(C)--(B)--(E)
						(Q)--(P) 
						 (F)--(C);
					\draw[dashed] (D)--(A)--(B) 
						 (A)--(F)
						 (P)--(M')--(N')--(Q)
						 (F)--(B)
						(F)--(J)--(C)--(A) 
						 (M)--(N)
						(X)--(Y) 
						(J)--(K); 
					\foreach  \i/\g in {A/-90,B/-90,C/-90,D/-90,E/0,F/90,I/0,J/-40,M/-90,N/90,P/0,Q/0,M'/-90,N'/80,X/90,Y/-40,K/-90,I/0}
					\fill (\i) circle (1pt)+(\g:3mm)  node[scale=1]{$\i$};
				\end{tikzpicture}
		\end{center}
		\begin{enumerate}[a)]
			\item Ta có $\heva{&BE\parallel AF \\&AF\subset (ADF)}\Rightarrow EB\parallel  (ADF)$.\\
			Tương tự $BC\parallel  (ADF)$.\\
			Từ đó ta có  $(BCE)\parallel(ADF)$.
			\item Vì $MM'\parallel AB\Rightarrow MM'\parallel CD$ nên theo định lí Thales ta có 
			$\dfrac{AM}{AC}=\dfrac{AM'}{AD} (1)$.\\ 
			Tương tự  $NN'\parallel AB\Rightarrow \dfrac{BN}{BF}=\dfrac{AN'}{AF} (2)$.\\
			Từ $(1)$ và  $(2)$ suy ra $\dfrac{AM'}{AD}=\dfrac{AN'}{AF}$.\\
			$\Rightarrow M'N'\parallel DF\subset  (DEF)\Rightarrow M'N'\parallel (DEF)$.\\
			Lại có $MM'\parallel CD\parallel EF\Rightarrow MM'\parallel (DEF)\Rightarrow (DEF)\parallel (MNN'M')$. 
			\item Gọi $P=MM'\cap BC$, $Q=NN'\cap BE$ và $J$, $K$ lần lượt là trung điểm các đoạn $AB$ và  $CF$.\\
			Gọi $X=N'Q\cap FJ$$, $ $Y=M'P\cap  CJ$ thì $XY=(MPQN')\cap (FCJ)$.\\
			Trong $(M'PQN')$, gọi  $I=XY\cap MN$.\\
			Ta có $\dfrac{YM}{AJ}=\dfrac{CM}{CA}\qquad (3)$ và $\dfrac{XN}{BJ}=\dfrac{FN}{FB} \qquad(4)$ mà $AJ=BJ$, $AC=BF$ nên từ $(3)$, $(4)$ suy ra $YM=XN\Rightarrow XMYN$ là hình bình hành nên  $I$ là trung điểm của $MN$.\\
			Do $\heva{&(M'PQN')\parallel (CEFE) \\&(CFJ)\cap (M'PQN')=XY \\&(CFJ)\cap (CEFE)=CF}\Rightarrow XY\parallel CF$ mà $IX=IY$ nên $I$ thuộc đường trung tuyến $JK$  của tam giác $JCF$.\\
			Giới hạn:\\ 
			Khi  $N\to B\Rightarrow M\to A\Rightarrow I\to J$.\\
			Khi  $N\to F\Rightarrow M\to C\Rightarrow I\to K$.\\
			Vậy  tập hợp điểm $I$ là đường trung tuyến $JK$ của tam giác $JCF$.
		\end{enumerate}
	 }
\end{ex}

% nguồn: Ok_TuanLuan_Toan11_Lan2
\begin{ex}%[1H4V4-6]%[ma:d886033c]
	Cho  hình chóp $S.ABCD$ có đáy $ABCD$ là hình bình hành và $M$, $N$, $P$ lần lượt là trung điểm các cạnh $AB$, $CD$, $SA$.
	\begin{enumerate}
		\item Chứng  minh $(SBN) \parallel (DPM)$.
		\item $Q$ là một điểm thuộc đoạn $SP$ ($Q$ khác $S$, $P$). Xác định thiết diện  của hình chóp cắt bởi $(\alpha)$ đi qua $Q$ và song song với $(SBN)$.
		\item Xác định thiết diện của hình chóp cắt bởi $(\beta)$ đi  qua $MN$ song song với $(SAD)$.
	\end{enumerate}
	\loigiai{ 
		\begin{enumerate}
			\item  Ta có
			 $\heva{& BN \parallel DM \\&DM \subset (DPM)} \Rightarrow BN \parallel (DPM)$.\tagEX{1}
			Tương tự $\heva{& BS \parallel MP \\& MP \subset(DPM)}  \Rightarrow BS \parallel (DPM)$.\tagEX{2}
			Từ (1) và (2) suy ra  $(SBN)\parallel (DPM)$.
			\begin{center}
				\begin{tikzpicture}[scale=1.5,font=\footnotesize,line join=round,line cap=round,>=stealth]
						\def\a{3} 
						\def\b{1.5} 
						\def\g{30} 
						 \def\h{2}
						 \path
						(0:0) coordinate (A)--++(\g:\b) coordinate (B)--++(0:\a) coordinate (C)--++(\g-180:\b) coordinate  (D)
						(A)--++(70:3.5)  coordinate (S);
						 \draw ($(A)!.5!(B)$) coordinate (M);
						\draw ($(C)!.5!(D)$) coordinate (N); 
						 \draw ($(S)!.5!(A)$) coordinate (P);
						\draw  ($(S)!.4!(P)$) coordinate (Q);
						\draw ($(N)!.4!(D)$) coordinate  (K);
						\draw ($(B)!.4!(M)$) coordinate (R); 
						\draw ($(S)!.4!(D)$)  coordinate (L);
						\foreach \x/\g  in {A/180,B/150,C/0,D/-90,S/90,M/160,N/-30,P/160,Q/160,R/-80,K/-80,L/-171}
						\fill(\x) circle  (1pt)
						($(\x)+(\g:2mm)$)  node{$\x$};
						\draw[dashed] (S)--(B)--(A) 
						(C)--(B) 
						 (N)--(B)
						(P)--(M)--(D) 
						(Q)--(R)--(K); 
						 \draw (S)--(A)--(D)--(C)--(S)--(D)
						 (S)--(N)
						(D)--(P) 
						 (K)--(L)--(Q);
					\end{tikzpicture}
			\end{center}
			\item Ta có $\heva{& SB \subset(SBN)\\ & (\alpha) \parallel (SBN)} \Rightarrow SB  \parallel (\alpha)$.\\
			Vậy $\heva{&Q \in (SAB) \cap (\alpha)\\&SB \subset (SAB) \\ & SB \parallel  (\alpha) } \Rightarrow (SAB) \cap (\alpha)=QR \parallel SB, R \in AB$.\\
			Tương tự $(\alpha) \cap (ABCD)=RK \parallel  BN$, $K \in CD$, $(\alpha) \cap (SCD)=KL \parallel SB$, $L \in SD$.\\
			Vậy thiết diện là tứ  giác $QRKL$.
			\begin{center}
				\begin{tikzpicture}[scale=1.5,font=\footnotesize,line join=round,line cap=round,>=stealth]
						\def\a{3} 
						 \def\b{1.5}
						 \def\g{30}
						\def\h{2} 
						 \path
						(0:0) coordinate (A)--++(\g:\b) coordinate (B)--++(0:\a)  coordinate (C)--++(\g-180:\b) coordinate (D)
						(A)--++(70:3.5) coordinate  (S);
						\draw  ($(A)!.5!(B)$) coordinate (M);
						\draw ($(C)!.5!(D)$)  coordinate (N);
						\draw ($(S)!.5!(B)$)  coordinate (F);
						\draw ($(S)!.5!(C)$) coordinate  (E);
						\foreach \x/\g in  {A/180,B/150,C/0,D/-30,S/90,M/160,N/-60,F/-50,E/40}
						 \fill(\x) circle (1pt)
						($(\x)+(\g:3mm)$)  node{$\x$};
						\draw[dashed] (S)--(B)--(A) 
						 (C)--(B)
						(N)--(M)--(F)--(E); 
						 \draw (S)--(A)--(D)--(C)--(S)--(D)
						 (E)--(N);
					\end{tikzpicture}
			\end{center}
			\item Ta có  $ \heva{& M \in(\beta) \cap (SAB) \\ & SA \parallel (\beta) \\& SA \subset (SAB)}\Rightarrow (\beta) \cap(SAB)=MF \parallel SA, F \in SB$.\\
			Tương tự $(\beta) \cap (SCD)=NE \parallel SD, E  \in SC$.
			 \\
			Thiết diện  là hình thang $MNEF$.
		\end{enumerate}
	 }
\end{ex}

% nguồn: Ok_TuanLuan_Toan11_Lan2
\begin{ex}%[1H4V4-6]%[ma:61b2863a]
	Cho hình chóp  $S.ABCD$, đáy là hình bình hành tâm $O$. Gọi $M$, $N$ lần lượt là trung điểm của $SA$ và $CD$.
	\begin{enumerate}
		\item Chứng  minh $(OMN) \parallel (SBC)$.
		\item Gọi $I$ là trung điểm của $SD$, $J$ là một điểm  trên $(ABCD)$ cách đều $AB$ và $CD$. Chứng minh $IJ \parallel (SAB)$.
	\end{enumerate}
	 \loigiai{
		\begin{center}
			\begin{tikzpicture}[scale=1.5,font=\footnotesize,line join=round,line cap=round,>=stealth]
					 \def\a{4}
					 \def\b{1.5}
					 \def\g{35}
					\def\h{2} 
					\path 
						(0:0) coordinate (A)--++(\g:\b) coordinate (B)--++(0:\a)  coordinate (C)--++(\g-180:\b) coordinate (D)
						 (A)--++(70:4) coordinate (S);
					\draw ($(A)!.5!(S)$) coordinate  (M);
					 \draw ($(C)!.5!(D)$) coordinate (N);
					\draw ($(D)!.5!(B)$) coordinate  (O);
					 \draw ($(S)!.5!(D)$) coordinate (I);
					\draw  ($(A)!.5!(D)$) coordinate (H);
					\draw ($(B)!.5!(C)$) coordinate (K); 
					 \draw ($(H)!.3!(K)$) coordinate (J);
					\foreach \x/\g in {A/180,B/150,C/0,D/-30,S/90,M/160,N/-60,O/-70,I/40,H/-90,K/-80,J/-60} 
					\fill(\x) circle  (1pt)
						($(\x)+(\g:2mm)$)  node{$\x$};
					\draw[dashed]  (H)--(K)--(I)--(J)
						 (S)--(B)--(A)
						(C)--(B) 
						(N)--(O)--(M)--cycle  (A)--(C)
						(B)--(D); 
					\draw (H)--(I) 
						 (S)--(A)--(D)--(C)--(S)--(D)
						(I)--(N); 
				\end{tikzpicture}
		\end{center}
		\begin{enumerate}
			\item  Do $O$, $M$ lần lượt là trung điểm của $AC$, $SA$ nên $OM$ là đường trung bình của tam giác $SAC$ ứng với cạnh $SC$. Suy ra $OM \parallel SC$. Mà $SC \subset (SBC) \Rightarrow OM \parallel (SBC)$. \tagEX{1}
			Tương tự $ON \parallel BC  \subset (SBC) \Rightarrow ON \parallel (SBC)$.\tagEX{2}
			Từ (1) và  (2) suy ra $(OMN) \parallel (SBC)$.
			\item Gọi $H$, $K$ lần lượt là  trung điểm của $AD$ và $BC$. Do $J \subset(ABCD)$ và $\mathrm{d}(J,AB)=\mathrm{d}(J,CD)$
			nên $J \in HK  \Rightarrow IJ \subset(IHK)$.\\
			 Ta dễ dàng chứng minh được $(IHK) \parallel (SAB)$.\\
			Vậy $\heva{&IJ\subset(IHK) \\ &(IHK) \parallel (SAB)} \Rightarrow IJ \parallel(SAB)$. 
		\end{enumerate}
	 }
\end{ex}

% nguồn: Ok_TuanLuan_Toan11_Lan2
\begin{ex}%[1H4V4-6]%[ma:40ed5359]
	Cho  hình hộp $ABCD.A'B'C'D'$. Chứng minh
	\begin{enumerate}
		\item[a)] $\left(BDA'\right)\parallel  \left(B'D'C\right)$.
		\item[b)] Đường chéo $AC'$ đi qua trọng tâm $G$ và $G'$  của hai tam giác $BDA'$ và $B'D'C$.
		\item[c)] $G$  và $G'$ chia đoạn $AC'$ thành ba phần bằng nhau.
	\end{enumerate}
	 \loigiai{
		\begin{center}
			\begin{tikzpicture}[scale=0.7,font=\footnotesize,line join=round,line cap=round,>=stealth]
					 \path
						(0,0) coordinate  (D)
						++(-130:3) coordinate (A) 
						++(0:6) coordinate (B) 
						($(D)+(B)-(A)$) coordinate (C) 
						 ($(A)!1/2!(C)$) coordinate (O);
					 \foreach \i in {A,B,C,D,O}{
						\coordinate (\i') at  ($(\i)+(0,-5)$);
					 }
					 \path
						($(A)!1/2!(C')$)  coordinate (I)
						($(A)!1/3!(C')$) coordinate (G) 
						($(A)!2/3!(C')$) coordinate (G'); 
					\draw (A)--(B)--(C)--(D)--cycle; 
					 \draw
						 (B)--(B')--(C)--(C')
						(A)--(A') 
						 (B)--(A')--(B')--(C')
						(A)--(C) 
						 (B)--(D);
					\draw[dashed] 
						 (D)--(A')--(D')--(A)
						 (C)--(D')--(C')
						 (A')--(C')--(A)
						 (D)--(D')--(B')
						(A')--(O) 
						 (C)--(O')
						(A')--(C); 
					\foreach \i/\g  in {A/-180,B/-30,C/60,D/120,O/90,A'/-120,B'/-60,C'/-30,D'/-150,O'/-90,I/90,G/90,G'/-90}
					\fill  (\i) circle (1.3pt)+(\g:4mm) node[scale=.7]{$\i$};
				\end{tikzpicture}
		\end{center}
		\begin{enumerate}
			\item[a)] Ta có $DD'\parallel BB'$ và $DD' = BB'$ (vì $ABCD.A'B'C'D'$ là  hình hộp)..\\
			Suy ra $DD'B'B$ là hình  bình hành.\\
			Do đó, $BD\parallel B'D'$ mà $B'D'\subset (B'D'C)$, suy ra $BD\parallel (B'D'C)$. \tagEX{1} 
			Chứng minh tương tự, ta có $DA'\parallel  B'C$.\\
			Mà $B'C\subset (B'D'C)$, suy ra  $DA' \parallel (B'D'C)$. \tagEX{2}
			Mặt khác, ta có $BD \cap  DA' = D$ và $BD$, $DA' \subset (BDA')$. \tagEX{3}
			Từ $(1)$, $(2)$ và  $(3)$ suy ra $(BDA')\parallel (B'D'C)$.
			\item[b)] Gọi $O$, $O'$ lần  lượt là tâm của hai đáy $ABCD$ và $A'B'C'D'$.\\
			Trong hình bình hành $AA'C'C$ gọi  $I$ là giao điểm của $AC'$ và $A'C$; cắt $A'O$, $CO'$ lần lượt tại $G$ và $G'$.\\
			Trong tam giác $AA'C$, ta có $G$ là giao điểm của hai trung tuyến $AI$ và $A'O$ nên $G$ là  trọng tâm của tam giác $AA'C$.\\
			Suy ra $G$ cũng là trọng tâm  tam giác $A'BD$.\\
			Tương tự, $G'$ là  trọng tâm tam giác $BDC$.\\
			Vậy $AC'$ đi qua trọng tâm của hai tam  giác $BDA'$ và $B'D'C$.
			\item[c)] Ta  có $AG=\dfrac{2}{3}AI=\dfrac{1}{3}AC'$ và $C'G'=\dfrac{2}{3}C'I=\dfrac{1}{3}AC'$.\\
			 Suy ra $G$ và $G'$ chia đoạn $AC'$ thành ba phần bằng nhau.
		\end{enumerate}
	} 
\end{ex}

% nguồn: Ok_TuanLuan_Toan11_Lan2
\begin{ex}%[1H4V4-6]%[ma:cc991066]
	Cho lăng trụ $ABCD.A'B'C'D'$ có đáy là hình thang, $AD=CD=BC=a$, $AB=2a$. Mặt phẳng $(\alpha)$  đi qua $A$ cắt các cạnh $BB'$, $ CC'$, $ DD'$ lần lượt tại $M$, $N$, $P$.
	\begin{enumerate}[a)]
		\item Tứ  giác $AMNP$ là hình gì?
		 \item So sánh $AM$ và $NP$.
	\end{enumerate}
	\loigiai{ 
		\begin{center}
			\begin{tikzpicture}[scale=1.2,font=\footnotesize,line join=round,line cap=round,>=stealth]
					\path 
						(0,0) coordinate  (A)
						(6,0)  coordinate (B)
						(4,-2) coordinate  (C)
						 (1,-2) coordinate (D)
						(1,2) coordinate (D') 
						 (0,4) coordinate (A')
						 (6,4) coordinate (B')
						 (4,2) coordinate (C')
						(1,-.48) coordinate  (P)
						(3,0) coordinate (I) 
						 (3,1) coordinate (J)
						(6,2) coordinate (M) 
						(4,.52) coordinate (N); 
					\draw(A')--(B')--(C')--(D')--(A')--(A)--(D)--(C)--(B)--(B') 
						(C)--(C') 
						(D)--(D') 
						 (A)--(P)--(N)--(M);
					\draw[dashed]  (M)--(A)--(B)
						 (C)--(I)--(J)--(N);
					\foreach \i/\g in  {A/180,B/-90,C/0,D/-30,A'/90,B'/0,C'/0,D'/40,P/-40,M/0,N/0,I/-120,J/90}
					\fill (\i) circle (1pt)+(\g:3mm)  node[scale=1]{$\i$};
				\end{tikzpicture}
		\end{center}
		\begin{enumerate}[a)]
			\item Ta có  $(ABB'A')\parallel (CDD'C')$$, $ $(\alpha)\cap (ABB'A')=AM$ \\
			và  $(\alpha)\cap (CDD'C')=NP\Rightarrow AM\parallel NP \qquad(1)$.\\
			Do đó $AMNP$ là  hình thang.
			\item Gọi $I$, $J$  lần lượt là trung điểm của $AB$, $AM$ thì $IC\parallel AD\Rightarrow IC\parallel (ADD'A')$.\\
			Lại có $IJ\parallel BB'\parallel AA'\Rightarrow IJ\parallel  AA'\subset (ADD'A')\Rightarrow (CIJN)\parallel (ADD'A')$.\\
			Mặt khác $(\alpha)\cap (ADD'A')=AP$ và  $(\alpha)\cap (CIJN)=JN$ nên $JN\parallel AP \qquad(2)$.\\
			Từ  $(1)$, $(2)$ suy ra $APNJ$ là hình bình hành, do đó $PN=AJ=\dfrac{1}{2}AM$.
		\end{enumerate}
	 }
\end{ex}

% nguồn: Ok_TuanLuan_Toan11_Lan2
\begin{ex}%[1H4V4-6]%[ma:f8b4c548]
	Cho hình hộp $ABCD.A'B'C'D'$ có tất cả các mặt đều là hình  vuông cạnh $a$. Các điểm $M$, $N$ nằm trên $AD'$, $BD$ sao cho $AM=DN=x$, $(0<x<a\sqrt{2})$.
	\begin{enumerate}[a)]
		\item Chứng minh khi $x$ biến thiên thì $MN$ luôn song  song với một mặt phẳng cố định.
		\item  Khi $x=\dfrac{a\sqrt{2}}{3}$, chứng minh $MN\parallel A'C$.
	\end{enumerate}
	\loigiai{ 
		\begin{center}
			\begin{tikzpicture}[scale=1,font=\footnotesize,line join=round,line cap=round,>=stealth]
					\path 
						(0,0)  coordinate (A)
						 (6,0) coordinate (B)
						(8,3) coordinate (C) 
						(2,3)  coordinate (D)
						(0,5) coordinate  (A')
						(6,5)  coordinate (B')
						 (8,8) coordinate (C')
						 (2,8) coordinate (D')
						 (4,1.5) coordinate (O)
						(1,1.5) coordinate (I) 
						(0.67,2.67)  coordinate (M)
						(3.33,2)  coordinate (N);
					 \draw(A')--(B')--(C')--(D')--(A')--(A)--(B)--(A')
						 (B')--(B)--(C)--(C');
					\draw[dashed]  (A)--(D)--(C)--(D')--(A)
						 (A')--(I)--(C)
						 (D')--(D)--(B)
						(A)--(C) 
						(M)--(N) 
						(A')--(C); 
					\foreach \i/\g in  {A/180,B/-90,C/0,D/130,A'/180,B'/0,C'/0,D'/120,O/-90,I/-80,M/30,N/-90}
					\fill (\i)  circle (1pt)+(\g:3mm) node[scale=1]{$\i$};
				\end{tikzpicture}
		\end{center}
		\begin{enumerate}[a)]
			\item Gọi $(\alpha)$  là mặt phẳng đi qua $M$ và song song với $(A'D'CB)$ và $N'=(\alpha)\cap BD$.\\
			Ta có  $\dfrac{AM}{AD'}=\dfrac{DN'}{DB} \qquad(1)$.\\
			Ta có $AD'=BD=a\sqrt{2}$ nên $AM=DN'$ mà $AM=DN$.\\ 
			Do đó $ DN=DN'\Rightarrow N\equiv N'$.\\ 
			Vậy $MN\subset (\alpha)\parallel (A'D'CB)$ do đó $MN$ song song với mặt phẳng  cố định $(A'D'CB)$.
			\item Khi $x=\dfrac{a\sqrt{2}}{3}$ thì dễ thấy $M$, $N$ lần lượt là trọng tâm các tam giác $A'AD$ và $CAD$ nên $A'M$ và $CN$ cắt nhau tại trung điểm $I$ của $AD$.\\ 
			Khi  đó $\dfrac{IM}{IA'}=\dfrac{IN}{IC}\Rightarrow MN\parallel A'C$.
		\end{enumerate}
	} 
\end{ex}

% nguồn: Ok_TuanLuan_Toan11_Lan2
\begin{ex}%[1H4V4-6]%[ma:7e64b88e]
	Cho tứ diện đều cạnh $a$. Gọi $I$, $J$ lần  lượt là trọng tâm các tam giác $ABC$ và $DBC$. Mặt phẳng $(\alpha)$ qua $IJ$ cắt các cạnh $AB$, $AC$, $DC$, $DB$ lần lượt tại $M$, $N$, $P$, $Q$.
	\begin{enumerate}[a)]
		\item Chứng minh $MN$, $PQ$,  $BC$đồng quy hoặc song song và $MNPQ$ là hình thang cân.
		\item Đặt $AM=x$, $AN=y$. Chứng minh $a(x+y)=3xy$. Tìm  GTNN và GTLN của $AM+AN$.
		\item Tính diện tích tứ giác $MNPQ$ theo $a$ và $s=x+y$. 
	\end{enumerate}
	 \loigiai{
		\begin{center}
			\begin{tikzpicture}[scale=1.2,font=\footnotesize,line join=round,line cap=round,>=stealth]
					\path 
						(2,4) coordinate (A) 
						(0,0)  coordinate (B)
						(2,-2)  coordinate (C)
						(6,0)  coordinate (D)
						 (3.71,-1.14) coordinate (P)
						 (1,-1) coordinate (E)
						(1.2,0) coordinate (Q) 
						 (2.67,-.67) coordinate (J)
						 (-1,1) coordinate (K)
						(0.4,0.8) coordinate (M) 
						(1.33,0.67)  coordinate (I)
						(2,0.57)  coordinate (N);
					 \draw(P)--(N)--(K)--(C)--(A)--(B)
						(C)--(D)--(A)--(E); 
					\draw[dashed]  (E)--(D)--(P)
						(P)--(K) 
						 (Q)--(M)
						(I)--(J) 
						(B)--(D); 
					\foreach \i/\g in {A/180,B/-90,C/-30,D/30,E/-90,P/-70,Q/-70,K/180,M/120,I/50,N/0,J/-90} 
					\fill (\i) circle (1pt)+(\g:3mm) node[scale=1]{$\i$}; 
				\end{tikzpicture}
		\end{center}
		\begin{enumerate}[a)]
			\item Ta có $(ABC)$, $(DBC)$, $(\alpha)$ đôi một cắt nhau theo các giao tuyến là $BC$, $MN$, $PQ$ nên theo định lí về giao tuyến thì $BC$, $MN$, $PQ$ hoặc  đồng quy hoặc đôi một song song.\\
			Ta chứng minh $MNPQ$ là hình  thang cân trong trường hợp $BC$, $MN$, $PQ$ đồng quy
			Gọi $E$ là trung điểm của $BC$  thì $\dfrac{EI}{EA}=\dfrac{EJ}{ED}\Rightarrow IJ\parallel AD$.\\
			Từ đó ta có $\heva{&IJ\subset (\alpha) \\&AD\subset (ACD) \\&IJ\parallel AD  \\&(\alpha)\cap (ACD)=NP}\Rightarrow NP\parallel IJ$.\\
			Tương tự $MQ\parallel IJ$ nên  $MNPQ$ là hình thang.\\
			 Dễ thấy $DQ=AM=x$, $DP=AN=y$.\\
			Theo định lí cô sinta có 
			 $MN^2=AM^2+AN^2-2AM.AN\cos{60^\circ}=x^2+y^2-xy$.\\
			Tương tự $PQ^2=DP^2+DQ^2-2DP.DQ\cos{60^\circ}=x^2+y^2-xy\Rightarrow  MN=PQ$.\\
			 Vậy $MNPQ$ là hình thang cân.\\
			Trường hợp $BC$, $MN$, $PQ$ song song không có gì khó khăn  bạn đọc tự kiểm tra.
			\item Ta  có
			\begin{eqnarray*}
				 &&S_{AMN}=S_{AIM}+S_{AIN}\\
				 &\Leftrightarrow&\dfrac{1}{2}xy\sin{60^\circ}=\dfrac{1}{2}x\cdot\dfrac{a\sqrt{3}}{3}\sin{30^\circ}+\dfrac{1}{2}y.\dfrac{a\sqrt{3}}{3}\sin{30^\circ}\\
				&\Leftrightarrow&a\left(x+y\right)=3xy. 
			\end{eqnarray*}
			Ta có $AM+AN=x+y$. Theo  BĐT Cauchy ta có
			\begin{eqnarray*}
				 &&a\left(x+y\right)=3xy\le 3\left(\dfrac{x+y}{2}\right)^2\\
				&\Leftrightarrow&3\left(x+y\right)^2-4a\left(x+y\right)\ge 0\\ 
				&\Leftrightarrow&x+y\ge \dfrac{4a}{3}\\ 
				&\Leftrightarrow&AM+AN\ge  \dfrac{4a}{3}.
			\end{eqnarray*}
			Đẳng thức xảy ra khi  $x=y=\dfrac{2a}{3}$, khi đó $(\alpha)$ đi qua $IJ$ và song song với $BC$.\\
			Không giảm tổng quát ta có thể giả sử $x\ge y$ khi đó $x\in  [\dfrac{2a}{3};a]$.\\
			Và 
			\begin{eqnarray*}
				 &&x+y=x+\dfrac{ax}{3x-a}=\dfrac{3x^2}{3x-a}\\
				&\Rightarrow&x+y-\dfrac{3a}{2}=\dfrac{3a^2}{3x-a}-\dfrac{3a}{2}=\dfrac{\left(a-x\right)\left(2a-x\right)}{3x-a}\le 0\\ 
				&\Rightarrow&x+y\le \dfrac{3a}{2}. 
			\end{eqnarray*}
			Đẳng thức xảy ra khi $x=a\Rightarrow y=\dfrac{a}{2}$. Khi  đó $(\alpha)$ đi qua $B$.\\
			Vậy $\min \left(AM+AN\right)=\dfrac{4a}{3}$,  $\max \left(AM+AN\right)=\dfrac{3a}{2}$.
			\item 
			Dễ thấy $MNPQ$ là hình  thang cân có $MQ=a-x$, $NP=a-y$.\\
			Giả sử $x\ge y\Rightarrow  a-x\le a-y$.
			\begin{center}
				\begin{tikzpicture}[scale=1.2,font=\footnotesize,line join=round,line cap=round,>=stealth]
						\path 
						(-1,2) coordinate (M) 
						(-3,-1) coordinate  (N)
						(4,-1)  coordinate (P)
						(2,2) coordinate  (Q)
						 (2,-1) coordinate (K)
						 (-1,-1) coordinate (H);
						\draw(M)--(N)--(P)--(Q)--node[midway,above]{$a-x$} 
						 (M)--(H)
						(Q)--(K)--node[midway,above]{$\frac{x-y}{2}$} 
						 (P);
						\foreach  \i/\g in {M/180,N/-90,P/0,Q/30,H/-90,K/-90}
						 \fill (\i) circle (1pt)+(\g:3mm) node[scale=1]{$\i$};
					\end{tikzpicture}
			\end{center}
			Ta  có
			\begin{eqnarray*}
				&&HN=\dfrac{\left(a-y\right)-\left(a-x\right)}{2}=\dfrac{x-y}{2}\\ 
				&\Rightarrow  &MH^2=MN^2-NH^2\\
				 & &=x^2+y^2-xy-\left(\dfrac{x-y}{2}\right)^2\\
				& &=\dfrac{3\left(x^2+y^2\right)-6xy}{4}=\dfrac{3s^2-8as}{4}\\ 
				 &\Rightarrow &MH=\sqrt{3xy}=\sqrt{a\left(x+y\right)}=\dfrac{\sqrt{3s^2-8as}}{2}.
			\end{eqnarray*}
			Do đó 
			\begin{eqnarray*}
				&S_{MNPQ}&=\dfrac{1}{2}\left(MQ+NP\right)MH\\ 
				&&=\dfrac{1}{2}\left(2a-\left(x+y\right)\right)\sqrt{3s^2-8as}\\ 
				&&=\dfrac{1}{4}\left(2a-s\right)\sqrt{3s^2-8as}. 
			\end{eqnarray*}
		\end{enumerate}
	 }
\end{ex}

% nguồn: Ok_TuanLuan_Toan11_Lan2
\begin{ex}%[1H4V4-6]%[ma:9e4f55c7]
	Cho hình chóp $S.ABCD$ có đáy $ABCD$ là hình chữ nhật và $\Delta SAD$ vuông tại $A$. Qua điểm $M$  trên cạnh $AB$ dựng mặt phẳng $(\alpha)$ song song với $(SAD)$ cắt $CD$, $SC$, $SB$ tại $N$, $P$, $Q$.
	\begin{enumerate}[a)]
		\item  Chứng minh $MNPQ$ là hình thang vuông.
		\item Gọi $I=NP\cap MQ$. Tìm tập hợp điểm $I$ khi  $M$ di động trên cạnh $AB$.
	\end{enumerate}
	 \loigiai{
		\begin{center}
			\begin{tikzpicture}[scale=1.2,font=\footnotesize,line join=round,line cap=round,>=stealth]
					 \path
						(1,5)  coordinate (S)
						(1,1) coordinate  (A)
						(4,1) coordinate  (M)
						(7,1) coordinate  (B)
						(5,-2) coordinate  (C)
						(2,-2) coordinate (N) 
						(-1,-2) coordinate (D) 
						(3,1.5)  coordinate (P)
						(4,3)  coordinate (Q)
						 (4,5) coordinate (I);
					 \draw(S)--(D)--(C)--(B)--(Q)--(I)--(N)
						 (S)--(C)
						 (P)--(Q);
					\draw(S)--(5,5)  node[above]{$d$};
					\draw[dashed] (S)--(A)--(D) 
						 (A)--(B)
						 (S)--(Q)--(M)--(N);
					 \foreach \i/\g in {A/180,B/-90,C/0,D/-30,S/90,M/-80,N/-90,P/0,Q/0,I/40}
					\fill (\i) circle (1pt)+(\g:3mm) node[scale=1]{$\i$}; 
				\end{tikzpicture}
		\end{center}
		\begin{enumerate}[a)]
			\item Ta có $\heva{&(\alpha)\parallel (SAB) \\&(ABCD)\cap (\alpha)=MN \\&(ABCD)\cap (SAB)=AB}\Rightarrow MN\parallel AB$.\\ 
			 Tương tự, ta có $(\alpha)\cap (SAB)=MQ\parallel SA$.\\
			$(\alpha)\cap (SCD)=NP\parallel SD$.\\ 
			 Thiết diện là tứ giác $MNPQ$.\\
			Do $\heva{&MN\parallel BC \\&MN\subset (\alpha) \\&BC\subset (SBC) \\&(SBC)\cap (\alpha)=PQ}\Rightarrow PQ\parallel MN  \qquad(1)$.\\
			Ta có $MN\parallel AD$, $MQ\parallel SA$ mà $AD\perp  SA$ nên $MN\perp MQ \qquad(2)$.\\
			Từ $(1)$, $(2)$ suy ra  $MNPQ$ là hình thang vuông.
			\item Gọi $d=(SAB)\cap (SCD)$, khi đó $I=NP\cap  MQ\Rightarrow \heva{&I\in NP\subset (SCD) \\&I\in MQ\subset (SAB)}\Rightarrow I\in d$.\\
			Từ đây dễ dàng tìm được quĩ  tích của điểm $I$.
		\end{enumerate}
	 }
\end{ex}
